% !TeX program = xelatex
% Standalone English volume. Use XeLaTeX until the references stabilize.
\documentclass[a4paper,11pt]{article}
\usepackage[T1]{fontenc}
\usepackage{lmodern}
\usepackage[english]{babel}
\usepackage[left=26mm,right=26mm,top=25mm,bottom=25mm,headheight=15pt,headsep=7mm,footskip=11mm]{geometry}
\usepackage{amsmath,amssymb,mathtools,amsthm}
\newtheorem{proposition}{Proposition}[section]
\newtheorem{lemma}[proposition]{Lemma}
\newtheorem{theorem}[proposition]{Theorem}
\newtheorem{corollary}[proposition]{Corollary}
\newtheorem{maintheorem}[proposition]{Main Theorem}
\theoremstyle{definition}
\newtheorem{definition}[proposition]{Definition}
\theoremstyle{remark}
\newtheorem{remark}[proposition]{Remark}
\makeatletter
\renewcommand*\l@section{\@dottedtocline{1}{0em}{2.4em}}
\renewcommand*\l@subsection{\@dottedtocline{2}{1.45em}{3.45em}}
\renewcommand*\l@subsubsection{\@dottedtocline{3}{3.8em}{4.55em}}
\makeatother
\usepackage{xcolor,fancyhdr,titlesec,needspace,xurl}
\usepackage{longtable,multicol}
\usepackage{hyperref}
\definecolor{ink}{HTML}{1A3542}
\definecolor{muted}{HTML}{4F606A}
\hypersetup{unicode=true,colorlinks=true,linkcolor=ink,filecolor=ink,urlcolor=ink,citecolor=ink,bookmarksnumbered=true,
 pdfdisplaydoctitle=true,pdflang={en},
 pdftitle={The generalized Riemann hypothesis for Dirichlet L-functions. Edition 930, Volume II},
 pdfsubject={A natural-language reconstruction of the proposed proof},
 pdfauthor={A mathematical reconstruction}}
\titleformat{\section}{\Large\bfseries\color{ink}}{\thesection}{.65em}{}
\titlespacing*{\section}{0pt}{.8em}{.65em}
\setlength{\parindent}{1.2em}
\linespread{1.045}
\setlength{\parskip}{.22em plus .05em minus .03em}
\setlength{\emergencystretch}{2em}
\setcounter{topnumber}{2}
\setcounter{tocdepth}{3}\setcounter{secnumdepth}{3}
\widowpenalty=4500
\clubpenalty=4500
\pagestyle{fancy}\fancyhf{}
\fancyhead[L]{\small\color{muted}Dirichlet GRH \textperiodcentered\ Volume I}
\fancyhead[R]{\small\color{muted}Dirichlet zeros and power intervals}
\fancyfoot[C]{\small\thepage}
\renewcommand{\headrulewidth}{.3pt}
\newcommand{\lead}[1]{\par\smallskip\noindent{\bfseries\color{ink}#1.}\enspace}
\newcommand{\sour}[2]{\unskip\penalty150\hspace{.35em}{\small\hyperlink{ref#1}{[#2]}}}
\titleformat{\subsection}{\large\bfseries}{\thesubsection}{.6em}{}
\titlespacing*{\subsection}{0pt}{1.15em plus .15em}{.4em}
\setlength{\abovedisplayskip}{9pt plus 2pt minus 3pt}
\setlength{\belowdisplayskip}{9pt plus 2pt minus 3pt}
\newcommand{\RR}{\mathbb R}\newcommand{\CC}{\mathbb C}\newcommand{\QQ}{\mathbb Q}
\newcommand{\ZZ}{\mathbb Z}\newcommand{\ord}{\operatorname{ord}}
\newcommand{\Frob}{\operatorname{Frob}}\newcommand{\Gal}{\operatorname{Gal}}
\newcommand{\len}{\operatorname{length}}\newcommand{\Res}{\operatorname*{Res}}
\newcommand{\wmean}{\operatorname{av}}\newcommand{\Ker}{\operatorname{ker}}
\newcommand{\Lh}{\mathcal L}\newcommand{\Gg}{\mathcal G}
\newcommand{\real}{\operatorname{Re}}\newcommand{\imag}{\operatorname{Im}}
\newcommand{\abs}[1]{\left|#1\right|}\newcommand{\eps}{\varepsilon}
\newcommand{\anchoritem}[2]{\par\Needspace{2.8\baselineskip}\hypertarget{ref#1}{}\noindent{\bfseries\color{ink}#2.}\enspace}
\newcommand{\srcloc}[2]{{\small\nolinkurl{#1.lean}}, line~#2}
\numberwithin{equation}{section}
\raggedbottom

\newcommand{\Sel}{\mathcal S}
\newcommand{\Cl}{\operatorname{Cl}}
\newcommand{\Aut}{\operatorname{Aut}}
\newcommand{\GL}{\operatorname{GL}}
\newcommand{\OO}{\mathcal O}
\hypersetup{bookmarksdepth=3}


\addto\captionsenglish{\renewcommand{\partname}{Volume}}


\makeatletter
\expandafter\def\csname r@cond:sizes\endcsname{{1}{26}{}{Item.1}{Dirichlet-GRH-v930-Volume-I.pdf}}
\expandafter\def\csname r@cond:weights\endcsname{{3}{26}{}{Item.3}{Dirichlet-GRH-v930-Volume-I.pdf}}
\expandafter\def\csname r@cor:alg97\endcsname{{87.3}{34}{}{proposition.87.3}{Dirichlet-GRH-v930-Volume-XVII.pdf}}
\expandafter\def\csname r@cor:alltwistcost738\endcsname{{50.10}{10}{}{proposition.50.10}{Dirichlet-GRH-v930-Volume-X.pdf}}
\expandafter\def\csname r@cor:markeddet715\endcsname{{47.37}{102}{}{proposition.47.37}{Dirichlet-GRH-v930-Volume-IX.pdf}}
\expandafter\def\csname r@cor:normalizedtransfer722\endcsname{{45.17}{55}{}{proposition.45.17}{Dirichlet-GRH-v930-Volume-IX.pdf}}
\expandafter\def\csname r@cor:sum98\endcsname{{25.3}{25}{}{proposition.25.3}{Dirichlet-GRH-v930-Volume-V.pdf}}
\expandafter\def\csname r@eq:gap\endcsname{{6.3}{26}{}{equation.6.3}{Dirichlet-GRH-v930-Volume-I.pdf}}
\expandafter\def\csname r@eq:length12\endcsname{{{L}}{30}{}{AMS.3}{Dirichlet-GRH-v930-Volume-I.pdf}}
\expandafter\def\csname r@eq:tail114\endcsname{{19.15}{17}{}{equation.19.15}{Dirichlet-GRH-v930-Volume-IV.pdf}}
\expandafter\def\csname r@lem:av10\endcsname{{6.2}{25}{}{proposition.6.2}{Dirichlet-GRH-v930-Volume-I.pdf}}
\expandafter\def\csname r@lem:det12\endcsname{{6.6}{29}{}{proposition.6.6}{Dirichlet-GRH-v930-Volume-I.pdf}}
\expandafter\def\csname r@lem:good175\endcsname{{72.18}{60}{}{proposition.72.18}{Dirichlet-GRH-v930-Volume-XIV.pdf}}
\expandafter\def\csname r@nav:I-28\endcsname{{6.6}{32}{}{subsection.6.6}{Dirichlet-GRH-v930-Volume-I.pdf}}
\expandafter\def\csname r@nav:III-12\endcsname{{25.1}{24}{}{subsection.25.1}{Dirichlet-GRH-v930-Volume-V.pdf}}
\expandafter\def\csname r@nav:III-13\endcsname{{25.2}{26}{}{subsection.25.2}{Dirichlet-GRH-v930-Volume-V.pdf}}
\expandafter\def\csname r@nav:volume-IV\endcsname{{19}{3}{}{section.19}{Dirichlet-GRH-v930-Volume-IV.pdf}}
\expandafter\def\csname r@prop:balance\endcsname{{6.3}{26}{}{proposition.6.3}{Dirichlet-GRH-v930-Volume-I.pdf}}
\expandafter\def\csname r@prop:blockrows719\endcsname{{45.7}{44}{}{proposition.45.7}{Dirichlet-GRH-v930-Volume-IX.pdf}}
\expandafter\def\csname r@prop:cost177\endcsname{{72.23}{65}{}{proposition.72.23}{Dirichlet-GRH-v930-Volume-XIV.pdf}}
\expandafter\def\csname r@prop:field12\endcsname{{6.4}{27}{}{proposition.6.4}{Dirichlet-GRH-v930-Volume-I.pdf}}
\expandafter\def\csname r@prop:maintamagawa725\endcsname{{45.25}{61}{}{proposition.45.25}{Dirichlet-GRH-v930-Volume-IX.pdf}}
\expandafter\def\csname r@prop:oddcyclo737\endcsname{{50.5}{6}{}{proposition.50.5}{Dirichlet-GRH-v930-Volume-X.pdf}}
\expandafter\def\csname r@prop:spec12\endcsname{{6.5}{28}{}{proposition.6.5}{Dirichlet-GRH-v930-Volume-I.pdf}}
\expandafter\def\csname r@prop:weights14\endcsname{{7.2}{89}{}{proposition.7.2}{Dirichlet-GRH-v930-Volume-I.pdf}}
\expandafter\def\csname r@sec:coho12\endcsname{{6.4}{28}{}{subsection.6.4}{Dirichlet-GRH-v930-Volume-I.pdf}}
\expandafter\def\csname r@sec:levelmass678\endcsname{{130.1}{3}{}{subsection.130.1}{Dirichlet-GRH-v930-Volume-XXV.pdf}}
\expandafter\def\csname r@sec:markedmatrix715\endcsname{{47.12}{100}{}{subsection.47.12}{Dirichlet-GRH-v930-Volume-IX.pdf}}
\expandafter\def\csname r@sec:missing\endcsname{{6.6}{32}{}{subsection.6.6}{Dirichlet-GRH-v930-Volume-I.pdf}}
\expandafter\def\csname r@sec:ray101\endcsname{{24.7}{16}{}{subsection.24.7}{Dirichlet-GRH-v930-Volume-V.pdf}}
\expandafter\def\csname r@sec:relationmatrices719\endcsname{{45.3}{43}{}{subsection.45.3}{Dirichlet-GRH-v930-Volume-IX.pdf}}
\expandafter\def\csname r@sec:tamerows725\endcsname{{45.8}{60}{}{subsection.45.8}{Dirichlet-GRH-v930-Volume-IX.pdf}}
\expandafter\def\csname r@seriesguide\endcsname{{1.1}{4}{}{subsection.1.1}{Dirichlet-GRH-v930-Volume-I.pdf}}
\expandafter\def\csname r@thm:balanced179\endcsname{{104.2}{4}{}{proposition.104.2}{Dirichlet-GRH-v930-Volume-XX.pdf}}
\expandafter\def\csname r@thm:charrows719\endcsname{{45.8}{44}{}{proposition.45.8}{Dirichlet-GRH-v930-Volume-IX.pdf}}
\expandafter\def\csname r@thm:conductor232\endcsname{{19.15}{12}{}{proposition.19.15}{Dirichlet-GRH-v930-Volume-IV.pdf}}
\expandafter\def\csname r@thm:criterion16\endcsname{{1.2}{2}{}{proposition.1.2}{Dirichlet-GRH-v930-Volume-I.pdf}}
\expandafter\def\csname r@thm:critical99\endcsname{{25.5}{27}{}{proposition.25.5}{Dirichlet-GRH-v930-Volume-V.pdf}}
\expandafter\def\csname r@thm:end14\endcsname{{7.3}{90}{}{proposition.7.3}{Dirichlet-GRH-v930-Volume-I.pdf}}
\expandafter\def\csname r@thm:explicit\endcsname{{3.3}{14}{}{proposition.3.3}{Dirichlet-GRH-v930-Volume-I.pdf}}
\expandafter\def\csname r@thm:grh16\endcsname{{1.1}{1}{}{proposition.1.1}{Dirichlet-GRH-v930-Volume-I.pdf}}
\expandafter\def\csname r@thm:integral181\endcsname{{104.9}{8}{}{proposition.104.9}{Dirichlet-GRH-v930-Volume-XX.pdf}}
\expandafter\def\csname r@thm:markedmatrix715\endcsname{{47.36}{101}{}{proposition.47.36}{Dirichlet-GRH-v930-Volume-IX.pdf}}
\expandafter\def\csname r@thm:norm91\endcsname{{94.9}{44}{}{proposition.94.9}{Dirichlet-GRH-v930-Volume-XVIII.pdf}}
\expandafter\def\csname r@thm:ribettrace720\endcsname{{45.11}{48}{}{proposition.45.11}{Dirichlet-GRH-v930-Volume-IX.pdf}}
\expandafter\def\csname r@thm:splitattachment722\endcsname{{45.16}{54}{}{proposition.45.16}{Dirichlet-GRH-v930-Volume-IX.pdf}}
\expandafter\def\csname r@thm:twists\endcsname{{5.1}{19}{}{proposition.5.1}{Dirichlet-GRH-v930-Volume-I.pdf}}
\expandafter\def\csname r@thm:wild98\endcsname{{25.1}{24}{}{proposition.25.1}{Dirichlet-GRH-v930-Volume-V.pdf}}
\makeatother
\begin{document}
\setcounter{part}{1}
\fancyhead[L]{\small\color{muted}Dirichlet GRH \textperiodcentered\ Volume II}
\fancyhead[R]{\small\color{muted}Local conditions and Selmer groups}
\thispagestyle{plain}\begin{center}
{\LARGE\bfseries The generalized Riemann hypothesis\\[2pt]for Dirichlet \(L\)-functions}\\[6pt]
{\small 27 November 2026 \quad\textperiodcentered\quad Edition 930}\end{center}
\part{Strict local conditions and Selmer groups of twists}\label{part:II}
\begin{abstract}\noindent We impose the strict local conditions on both conjugation eigenspaces. Tame inertia and auxiliary twists give tests for the Selmer groups, whose lengths are compared with critical values over a common interval.\end{abstract}
\setcounter{section}{10}
\noindent The \hyperref[seriesguide]{guide in Volume~I} lists all twenty-seven volumes.
\begingroup\small\setlength{\parskip}{0pt}\tableofcontents\endgroup
\clearpage

\section{Strict local conditions and conjugation}\label{nav:g-II-0}\label{nav:volume-II}

\subsection{A uniform bound for one conjugation eigenspace}\label{nav:II-00}
\label{sec:critical13}

Fix the primitive character \(\psi\) of conductor \(F\). For this section use the primes, fields and interval of Proposition~\ref{prop:field12}. At a given \(x\), write \(g=[E:\QQ]\), \(G=\Gal(E/\QQ)\), \(c\in G\) for complex conjugation, and \(t=\lfloor x^{401}\rfloor\). The integer
\[
                       N=2^{l+2}Fqd
\]
is a cyclotomic modulus for \(E\). The letter \(d\) here denotes the second auxiliary prime. The field \(E\) is CM, since it is abelian and contains the imaginary quadratic field \(\QQ(\sqrt{-q})\). For an odd prime \(\ell\) and \(n=r+1\ge2\), put \(e_n=(1+(-1)^nc)/2\in\ZZ_\ell[G]\). We shall bound \(e_n\mathcal S_{\ell,r}(E)\) for all odd primes at once. The sign depends on the twist.

\begin{theorem}[The contribution with sign \((-1)^{r+1}\)]\label{thm:critical13}
Fix \(0<\eps<1/52\). For all sufficiently large \(Y\), take the single pair \(q,d\) of Proposition~\ref{prop:field12}. At every \(2Y\le x\le(2Y)^{3/2}\) its fields satisfy
\begin{equation}
 \sum_{\ell>2\ {\rm prime}}(\log\ell)\sum_{1\le r<t}
          \len_{\ZZ_\ell}\big(e_{r+1}\mathcal S_{\ell,r}(E)\big)
       \le C_{F,\eps} x^{812}
       \le \tfrac14\exp(x^\eps/2).
 \label{eq:critical13}
\end{equation}
Every group in the sum is finite, and for fixed \(x\) only finitely many terms are nonzero.
\end{theorem}

\subsubsection{From strict classes to special values}\label{nav:II-01}
Put \(S_0=\{\infty\}\cup\{p:p\mid N\}\), \(U_\ell=\operatorname{Spec}\mathcal O_E[1/(N\ell)]\), and \(C_\ell=H^2_{\rm et}(U_\ell,\ZZ_\ell(n))\). Let \(\Theta_n\in\QQ[G]\) be the equivariant special value characterized by
\[
     \chi(\Theta_n)=L_{S_0}(1-n,\chi^{-1})\qquad(\chi\in\widehat G).
\]
On the right the Euler factors at \(p\mid N\) are removed. This is the value of the Dirichlet function of the character induced to modulus \(N\). The Bernoulli formula below also constructs \(\Theta_n\) over \(\QQ\). Complex conjugation acts on all cohomology groups through the natural action of the field automorphism.

\begin{lemma}[The relevant Fitting ideal]\label{lem:fit13}
There is a \(G\)-equivariant injection \(\mathcal S_{\ell,r}(E)^\vee\hookrightarrow C_\ell\), where the dual has the contragredient action. If a positive integer \(W\) kills \(H^0(E,\QQ_\ell/\ZZ_\ell(n))\), then \(e_nW\Theta_n\) belongs to the zeroth Fitting ideal of \(e_n C_\ell\) over \(R_\ell=e_n\ZZ_\ell[G]\).
\end{lemma}
\begin{proof}
Poitou--Tate duality pairs the strict kernel in degree one for \(\QQ_\ell/\ZZ_\ell(1-n)\) with the strict kernel in degree two for its Tate dual \(\ZZ_\ell(n)\) on \(\operatorname{Spec}\mathcal O_E[1/\ell]\). One obtains this form of duality by taking direct and inverse limits of the finite-coefficient duality sequence. The finite-level groups are finite, so the inverse systems satisfy the Mittag--Leffler condition. On the discrete side vanishing at the finitely many local places can be tested in the direct limit. The pairing is preserved by field automorphisms, and the archimedean degree-one local terms vanish for these odd primes. Thus the dual embeds into global degree two. Inverting the extra primes above \(N\) gives an injection in degree two by the localization sequence, since the preceding terms at \(w\nmid\ell\) are \(H^0(\kappa_w,\ZZ_\ell(n-1))=0\). Indeed multiplication by \((Nw)^{n-1}-1\) has zero kernel in \(\ZZ_\ell\). This proves the injection. We use continuous étale cohomology and the usual identification with Galois cohomology on rings of \(S\)-integers where \(\ell\) is invertible.~\cite{II-II-II-duality13,II-II-II-gp13}

We use the refined Coates--Sinnott theorem of Greither and Popescu. For an abelian extension of the totally real base \(\QQ\), a set \(S_0\) containing the infinite and ramified primes, \(n\ge2\), and \(\ell>2\), it identifies
\[
 \operatorname{Ann}_{\ZZ_\ell[G]}
    \big(H^1_{\rm et}(U_\ell,\ZZ_\ell(n))_{\rm tors}\big)\,\Theta_n
       = e_n\,\operatorname{Fitt}_{\ZZ_\ell[G]}(C_\ell).
\]
The hypothesis on the cyclotomic \(\mu\)-invariant holds by the Ferrero--Washington theorem, because the maximal CM subfield of \(E(\mu_\ell)\) is abelian over \(\QQ\). The torsion group in the formula is \(H^0(E,\QQ_\ell/\ZZ_\ell(n))\). This follows from the coefficient sequence for \(\ZZ_\ell(n)\subset\QQ_\ell(n)\), since \(H^0(E,\QQ_\ell(n))=0\). Projecting the equality to the direct factor \(R_\ell\) and using base change for Fitting ideals proves the assertion.~\cite[Thm.~6.11 and Lem.~6.9]{II-II-II-gp13}\cite{II-II-II-fw13}
\end{proof}

\subsubsection{One integral determinant for all odd primes}\label{nav:II-02}
Choose two distinct fixed rational primes \(P_1,P_2\nmid2F\). For large \(Y\) these differ from \(q,d\), so \(E\) is unramified at both. Put
\[
 W=(P_1^{gn}-1)(P_2^{gn}-1),\qquad A=W\,nN(n+1)! .
\]
For each \(\ell\) choose \(P_i\ne\ell\). At a place of \(E\) over \(P_i\), the cyclotomic character takes arithmetic Frobenius to \(P_i^f\), where the residue degree \(f\) divides \(g\). Consequently \(P_i^{gn}-1\) kills \(H^0(E,\QQ_\ell/\ZZ_\ell(n))\). This proves the required annihilation by a single \(W\) independent of \(\ell\).

\begin{lemma}[Integrality, nonvanishing and size]\label{lem:int13}
The element \(A\Theta_n\) lies in \(\ZZ[G]\). The integer
\[
 D_n=\prod_{\substack{\chi\in\widehat G\\\chi(c)=(-1)^n}}
             A\,L_{S_0}(1-n,\chi^{-1})
\]
is nonzero and satisfies \(\log|D_n|\ll_F gn\{g+\log(2Nn)\}\).
\end{lemma}
\begin{proof}
The standard special-value formula for residue classes, obtained from the Hurwitz zeta function at \(1-n\), is
\[
 L_{S_0}(1-n,\chi^{-1})=
    -\frac{N^{n-1}}{n}
      \sum_{\substack{1\le a\le N\\(a,N)=1}}\chi^{-1}(a)B_n(a/N).
\]
Here \(B_n(z)=\sum_{j=0}^n\binom{n}{j}B_j z^{n-j}\). The recurrence \(\sum_{j=0}^m\binom{m+1}{j}B_j=0\) for \(m\ge1\), with \(B_0=1\), gives \((m+1)!B_m\in\ZZ\) by induction. Hence \(nN(n+1)!\) clears every rational coefficient in the last display. Replacing each \(\chi^{-1}(a)\) with the inverse of the cyclotomic Artin symbol of \(a\) in \(G\), constructs the asserted integral group-ring element.~\cite{II-II-II-bernoulli13}

For the indicated parity, \(1-n\) is not a trivial zero. Writing \(\chi(-1)=(-1)^a\) with \(a\in\{0,1\}\), the gamma argument \((1-n+a)/2\) is a half-integer for this parity. The completed functional equation relates the value to that of the conjugate character at \(n>1\), where the Euler product is nonzero. For the trivial character the same conclusion follows from the zeta functional equation. Every removed factor of a primitive unramified character has the form \(1-\chi^{-1}(p)p^{n-1}\), which is nonzero. The displayed product is Galois-invariant because conjugation permutes its characters. Each factor is an algebraic integer by the first assertion, proving that \(D_n\) is a nonzero rational integer.

On the circle \(|z|=1\), the generating function \(ze^{vz}/(e^z-1)\) is bounded uniformly for \(0\le v\le1\). Cauchy's estimate gives \(|B_n(v)|\ll n!\). The special-value formula therefore bounds each absolute value by \(O(N^n n!)\) before multiplication by \(A\). Since \(\log W\ll_F gn\), \(\log((n+1)!)\ll n\log(2n)\), and the number of factors is at most \(g\), the stated logarithmic estimate follows.
\end{proof}

\subsubsection{The integral group ring at primes dividing the degree}\label{nav:II-03}
Write \(M_\ell=e_n C_\ell\) and \(b_\ell=e_nA\Theta_n\). The finitely generated \(\ZZ_\ell\)-module \(R_\ell\) is free. On \(R_\ell\otimes\overline{\QQ}_\ell\), multiplication by \(b_\ell\) is diagonal with eigenvalues indexed by the characters of the indicated parity. Its determinant is \(D_n\). Thus \(R_\ell/b_\ell R_\ell\) has \(\ZZ_\ell\)-length \(v_\ell(D_n)\) by the elementary divisor theorem. Lemma~\ref{lem:fit13} puts \(b_\ell\) in the zeroth Fitting ideal of \(M_\ell\), which annihilates \(M_\ell\).

If \(\ell\nmid g\), \(R_\ell\) is a product of rings of integers of unramified extensions of \(\QQ_\ell\). Indeed the cyclic group-ring factors are defined by \(X^m-1\) with \(m\) prime to \(\ell\), whose reduction has distinct roots. The abelian group ring and its direct factors are therefore finite étale over \(\ZZ_\ell\). In each discrete valuation ring the valuation of the zeroth Fitting ideal equals the length of the module. Membership of \(b_\ell\) bounds that length by the valuation of the corresponding component of \(b_\ell\). Weighting by the residue degrees and summing gives
\[
                           \len_{\ZZ_\ell}M_\ell\le v_\ell(D_n).
\]

For the other odd primes, \(\ell\mid g\) implies \(\ell\mid o\), where \(o\) is the fixed order of \(\psi\). To see this, \([K_\psi C_l:\QQ]=o\,2^l/[K_\psi\cap C_l:\QQ]\), and adjoining the two quadratic fields multiplies this by four. We bound the number of \(\ZZ_\ell\)-generators of \(C_\ell\). Its reduction modulo \(\ell\) injects into \(H^2_{\rm et}(U_\ell,\mathbf F_\ell(n))\) by the coefficient sequence. Put \(K=E(\mu_\ell)\) and pass to \(\operatorname{Spec}\mathcal O_K[1/(N\ell)]\). Restriction is injective on the latter cohomology by trace, since the finite étale cover has degree dividing \(\ell-1\). Over \(K\) choose an isomorphism of \(\mathbf F_\ell(n)\) with \(\mu_\ell\). Kummer theory bounds this degree-two dimension by the sum of the dimensions of the class group modulo \(\ell\) and the \(\ell\)-torsion in the Brauer group of the open ring. Its class group is a quotient of \({\rm Cl}(K)\). By the global invariant sequence for the Brauer group, its \(\ell\)-torsion injects into the local \(\ell\)-torsion at the missing finite places. Each such place contributes dimension at most one, and there are at most \([K:\QQ](1+\log_2 N)\) of them.~\cite{II-II-II-duality13,II-II-II-milne}

Writing \(a=[K:\QQ]\), we have \(a\le(\ell-1)g\) and \(|\operatorname{disc}K|\le(N\ell)^a\) by the abelian conductor--discriminant formula. Minkowski's bound represents each ideal class by an integral ideal of norm at most \(B=2^a|\operatorname{disc}K|^{1/2}\). The number of such ideals is at most \(B^2\zeta_K(2)\le B^2\zeta(2)^a\), the last inequality following from the Euler factors. Thus \(\log|{\rm Cl}(K)|\ll_F g\log(2N)\). Nakayama's lemma now bounds the number of \(\ZZ_\ell\)-generators of \(C_\ell\), hence of its summand \(M_\ell\), by \(C_Fg\log(2N)\). Finite generation here is the standard finiteness theorem for continuous étale cohomology of an open ring of integers.~\cite{II-II-II-gp13}

These generators also generate over \(R_\ell\). As \(b_\ell M_\ell=0\), a direct sum of that many copies of \(R_\ell/b_\ell R_\ell\) surjects onto \(M_\ell\). Its length is therefore at most \(C_Fg\log(2N)v_\ell(D_n)\). This argument includes the primes at which the group ring has nonsemisimple reduction.

\begin{proof}[Proof of Theorem~\ref{thm:critical13}]
The injection of Lemma~\ref{lem:fit13} preserves the sign on complex conjugation, since \(c^{-1}=c\) on the dual. Every term in the sum of lengths for fixed \(r\) is now finite. Combining the two cases above and summing \(v_\ell(D_n)\log\ell\) gives
\[
 \sum_{\ell>2}(\log\ell)\len_{\ZZ_\ell}(e_n\mathcal S_{\ell,r}(E))
 \ll_F (1+g\log(2N))\,gn\{g+\log(2Nn)\}.
\]
The sum of valuations over all primes is \(\log|D_n|\), so the estimate also gives finite support. On the stated family \(g\ll_F x^3\) and \(\log(2Nt)\ll_{F,\eps} x\), uniformly, by the bounds on \(q,d,2^l\) and \(Y\le x\). When \(n\le t\) the right side is \(O_{F,\eps}(x^{10}n)\). Finally \(\sum_{2\le n\le t}n\le t^2\le x^{802}\). This even gives \(O_{F,\eps}(x^{812})\), and any fixed power of \(x\) is \(o(\exp(x^\eps/2))\).
\end{proof}
The complementary eigenspace has \(\chi(c)=(-1)^{n-1}\). Its primitive \(L\)-function has a trivial zero at \(1-n\), as follows from the pole of the gamma factor there and the functional equation. The determinant in Lemma~\ref{lem:int13} would vanish on this component. Controlling its local lengths requires a different arithmetic estimate.

\subsection{Both signs at finite level and the exceptional twists at two}\label{nav:II-08}
\label{sec:two15}

Continue with the same \(\psi,\eps,Y,q,d\), and with the fields of Proposition~\ref{prop:field12}. Write \(t(x)=\lfloor x^{401}\rfloor\) and \(I_Y=[2Y,(2Y)^{3/2}]\). We first establish finiteness of the full sum in \eqref{eq:length12} at each \(x\). We then give uniform estimates which also see the complementary sign.

\begin{proposition}[Finiteness of the full sum]\label{prop:finite15}
For every \(x\in I_Y\) and \(r\ge1\) one has
\[
 \sum_\ell(\log\ell)\len_{\ZZ_\ell}\mathcal S_{\ell,r}(E)
                  \le \log |K_{2r}(\mathcal O_E)|<\infty .
\]
Here \(K_i\) denotes Quillen's algebraic \(K\)-group. In particular the supremum in \eqref{eq:length12} is a finite sum at each \(x\).
\end{proposition}
\begin{proof}
The Poitou--Tate argument of Lemma~\ref{lem:fit13}, before the extra primes are inverted, embeds the dual of the strict group into \(H^2_{\rm et}(\mathcal O_E[1/\ell],\ZZ_\ell(r+1))\). It applies at two as well because \(E\) contains an imaginary quadratic field and has no real places. In the duality sequence the modified local groups at complex places are zero.~\cite{II-II-II-duality13}

Use the number-field Quillen--Lichtenbaum comparison in even degree. For a ring of \(S\)-integers with all primes over \(\ell\) inverted it identifies the \(\ell\)-adic completion of \(K_{2r}\) with \(H^2_{\rm et}(-,\ZZ_\ell(r+1))\), for \(r\ge1\) and odd \(\ell\), and also for \(\ell=2\) in the totally imaginary case. In the descent spectral sequence the only other possible term in this degree is \(H^0(-,\ZZ_\ell(r))\), which vanishes because the cyclotomic image over a number field is open and \(r>0\). The cohomological dimension is at most two under these hypotheses.~\cite[Ch.~VI, Sec.~8]{II-II-II-weibel15}\cite[Thm.~1]{II-II-II-kahn15}

In Quillen's localization sequence for \(\mathcal O_E\subseteq\mathcal O_E[1/\ell]\) the groups surrounding the map in degree \(2r\) are the sums, over \(v\mid\ell\), of \(K_{2r}(\kappa_v)=0\) and \(K_{2r-1}(\kappa_v)\), respectively. The latter are cyclic of order \(|\kappa_v|^r-1\), prime to \(\ell\). Thus the map is an isomorphism after tensoring with \(\ZZ_\ell\). The group \(K_{2r}(\mathcal O_E)\) is finite by Quillen's finite generation theorem and Borel's rank theorem in positive even degree. Its \(\ell\)-part therefore bounds the order of the strict group. Taking logarithms and summing proves the claim. There are finitely many indices \(r<t(x)\) at a fixed \(x\).~\cite{II-II-II-weibel15}
\end{proof}

\begin{theorem}[A uniform range for both signs]\label{thm:trunc15}
Let \(P=\exp(x^\eps/12)\) and \(b_0=\lfloor\exp(x^\eps/8)\rfloor\). For every sufficiently large \(Y\) and every \(x\in I_Y\),
\[
 \sum_{\ell\le P}(\log\ell)
       \sum_{1\le r<t(x)}
       \len_{\ZZ_\ell}\!\left(\mathcal S_{\ell,r}(E)[\ell^{b_0}]\right)
                    \le \exp(x^\eps/3).
\]
The sum is over rational primes, including two. The bracket denotes the subgroup killed by the indicated power.
\end{theorem}
\begin{proof}
Write \(g=[E:\QQ]\) and \(N=2^{l+2}Fqd\) as before. Let \(d_{\ell,r}\) be the dimension of \(\mathcal S_{\ell,r}(E)[\ell]\) over \(\mathbf F_\ell\). The multiplication-by-\(\ell\) sequence on \(\QQ_\ell/\ZZ_\ell(-r)\) maps \(H^1(G_{E,S_\ell},\mathbf F_\ell(-r))\) onto the \(\ell\)-torsion of the ambient degree-one group. Therefore \(d_{\ell,r}\) is at most the dimension of this finite-coefficient cohomology.

Put \(K=E(\mu_\ell)\), \(a=[K:\QQ]\le g\ell\). Restriction to \(K\) is injective on the finite-coefficient group by corestriction, since \([K:E]\) divides \(\ell-1\) and the extension is unramified outside \(\ell\). The module becomes isomorphic to \(\mu_\ell\). Let \(S_K\) be the primes over \(\ell\) and \(R_K=\mathcal O_K[1/\ell]\). The Kummer sequence gives
\[
 0\longrightarrow R_K^\times/(R_K^\times)^\ell
 \longrightarrow H^1(G_{K,S_K},\mu_\ell)
 \longrightarrow {\rm Cl}(R_K)[\ell]\longrightarrow0 .
\]
Dirichlet's \(S\)-unit theorem bounds the first dimension by \(a+|S_K|\le2a\), including roots of unity. The last dimension is at most \(\log|{\rm Cl}(K)|/\log\ell\). The conductor of \(K\) divides \(N\ell\). The same Minkowski and ideal-count argument as in Section~\ref{sec:critical13} gives \(\log|{\rm Cl}(K)|\ll a\log(2N\ell)\) with an absolute constant. We obtain \(d_{\ell,r}\log\ell\ll g\ell\log(2N\ell)\) for every twist.

In a finite abelian \(\ell\)-group with \(d_{\ell,r}\) cyclic summands, the subgroup killed by \(\ell^{b_0}\) has length at most \(b_0 d_{\ell,r}\). Sum over \(r<t(x)\) and use \(\sum_{\ell\le P}\ell\le P(P+1)\). The result is \(O(t(x)b_0 g P^2\log(2NP))\). The factors \(t(x),g,\log(2NP)\) are bounded by fixed powers of \(x\) on this family. Since \(1/8+2/12=7/24<1/3\), the claimed bound follows throughout \(I_Y\).
\end{proof}

\subsubsection{One dyadic module for the whole interval}\label{nav:II-09}
Put \(k_\infty=K_\psi(\mu_{2^\infty})\), \(B=K_\psi(\sqrt{-q},\sqrt{-d})\), and \(L_\infty=B(\mu_{2^\infty})\). For every \(x\in I_Y\) one has \(E_x=BC_l\) and \(E_x(\mu_{2^\infty})=L_\infty\). In this subsection the symbol \(B\) refers to this base field. Let \(M\) be the Galois group over \(L_\infty\) of its maximal everywhere unramified abelian pro-\(2\) extension.

\begin{lemma}[A uniform generator bound]\label{lem:generators15}
The compact \(\ZZ_2\)-module \(M\) can be generated by at most \(R\) elements, for a constant \(R\) depending only on \(\psi\), independently of \(q,d,Y\).
\end{lemma}
\begin{proof}
The fixed field \(k_\infty\) is the cyclotomic \(\ZZ_2\)-extension of \(K_\psi(i)\). The latter contains \(\mu_4\), so its cyclotomic image is an open subgroup of \(1+4\ZZ_2\simeq\ZZ_2\). Iwasawa's class module theorem identifies the maximal everywhere unramified abelian pro-\(2\) Galois group over \(k_\infty\) with the inverse limit of the dyadic ideal class groups in this tower. It is a finitely generated torsion module over the one-variable Iwasawa algebra. Its \(\mu\)-invariant vanishes by Ferrero--Washington, since \(K_\psi(i)/\QQ\) is abelian. The structure theorem over \(\ZZ_2[[T]]\) then gives finite generation over \(\ZZ_2\). Indeed, its distinguished polynomial factors give finite free \(\ZZ_2\)-modules before taking powers and its pseudo-isomorphism contributes only finite groups. Denote by \(r_0\) the dimension of its quotient modulo two. This number is independent of \(Y\).~\cite{II-II-II-iwasawa15,II-II-II-fw13}

The number of primes above \(q,d\) in \(k_\infty\) is bounded independently of these primes. In \(\QQ(\mu_{2^\infty})/\QQ\) the decomposition group of an odd prime \(p\) is the closure of \(\langle p\rangle\) in \(\ZZ_2^\times\), since the extension is unramified there and Frobenius acts by \(p\). Put \(a=v_2(p^2-1)\). For \(q\equiv3\bmod8\) one has \(a=3\), and for \(d\equiv7\bmod16\) one has \(a=4\). The powers of \(p^2\) generate \(1+2^a\ZZ_2\) topologically by the dyadic logarithm. Thus the index of each decomposition group is at most \(2^{a-1}\le8\). Passing to \(k_\infty\) multiplies the number of primes by at most \(o=[K_\psi:\QQ]\). In total there are at most \(16o\) such places.

Let \(\mathcal P\) be the maximal pro-\(2\) Galois group over \(k_\infty\) unramified away from these places. Its maximal elementary abelian quotient is its Frattini quotient, obtained by killing squares and commutators and taking the closure. On additionally killing all inertias above \(q,d\), this quotient has dimension \(r_0\). At a place of odd residue characteristic \(p\), the dyadic image of inertia is procyclic. Indeed wild inertia is pro-\(p\) and the pro-\(2\) quotient of tame inertia is procyclic. Here the cyclotomic tower over the finite base is unramified at \(q,d\), so the absolute inertia groups can equally be computed over that finite base. Conjugation has no effect in the elementary abelian quotient. Each of the at most \(16o\) places thus contributes at most one further generator. The pro-\(2\) basis theorem gives at most \(r_0+16o\) topological generators for \(\mathcal P\).

The extension \(L_\infty/k_\infty\) has degree at most four, and is unramified away from \(q,d\), by the quadratic discriminants. It corresponds to a subgroup of \(\mathcal P\) of index at most four. Every finite unramified abelian \(2\)-extension of \(L_\infty\) lies in the corresponding maximal extension: its finitely many conjugates over \(k_\infty\) have compositum a Galois \(2\)-extension of \(k_\infty\) unramified away from \(q,d\). Therefore \(M\) is a quotient of that subgroup. The Schreier generator bound for an index-\(e\) subgroup of a \(d\)-generated pro-\(2\) group is \(1+e(d-1)\). It follows by choosing coset representatives along a spanning tree in the coset graph and taking the remaining Schreier generators, then taking the closure. We may take \(R=1+4(r_0+16o)\).
\end{proof}

\begin{theorem}[A common set of exceptional twists at two]\label{thm:exception15}
There is a constant \(R=R(\psi)\) such that, for each sufficiently large \(Y\) and the single chosen pair \(q,d\), there exists one set \(\mathcal R_Y\) of at most \(R\) positive integers with
\[
 \len_{\ZZ_2}\mathcal S_{2,r}(E_x)\ll_{\psi,\eps}\log(2x)
 \qquad (x\in I_Y,\ 1\le r<t(x),\ r\notin\mathcal R_Y).
\]
In particular the contribution of all these twists at two to \eqref{eq:length12} is \(O_{\psi,\eps}(t(x)\log(2x))\).
\end{theorem}
\begin{proof}
The group \(\Gamma_*=\Gal(L_\infty/B(i))\) acts on \(M\) by conjugation. The action is well-defined because the maximal unramified abelian extension is invariant and \(M\) is abelian. The cyclotomic character embeds \(\Gamma_*\) as an open subgroup of \(1+4\ZZ_2\). Choose a topological generator \(\gamma\) and write \(u=\kappa(\gamma)\), \(e_*=v_2(u-1)\). Its index in \(1+4\ZZ_2\) is \(2^{e_*-2}\) and is at most \([B(i):\QQ]\le8o\), by intersecting \(B(i)\) with the cyclotomic extension. In particular \(e_*\) is uniformly bounded.

Take a minimal surjection \(\pi:\ZZ_2^D\to M\), so \(D\le R\) and its reduction modulo two is an isomorphism onto \(M/2M\). If \(D=0\), Proposition~\ref{prop:spec12} already bounds every length by one. Otherwise lift the images under \(\gamma\) of these \(D\) generators to columns of a matrix \(A\) over \(\ZZ_2\). Then \(\pi A=\gamma\pi\) and \(A\) is invertible since its reduction is. This also handles finite torsion in \(M\). The image of the pro-\(2\) group \(\Gamma_*\) on \(M/2M\) is a finite \(2\)-group, so the reduction of \(A\) is unipotent. Fix \(h_*=\lceil\log_2 R\rceil+1\). In characteristic two, \((I+U)^{2^k}=I+U^{2^k}\) for a nilpotent \(D\)-by-\(D\) matrix \(U\). Consequently \(A^{2^{h_*-1}}\equiv I\bmod2\). Squaring gives \(A'=A^{2^{h_*}}\equiv I\bmod4\).

The subgroup \(\Gal(L_\infty/E_x(i))\) is generated topologically by \(\gamma^{2^{h_x}}\), where
\[
  2^{h_x}=[B(i)C_l:B(i)]=[C_l:C_l\cap B(i)].
\]
Thus \(2^l/(8o)\le2^{h_x}\le2^l\). Throughout \(I_Y\), \(h_x\ge h_*\) for large \(Y\) and \(h_x\ll_\psi\log(2x)\).
The module \(X'_2(E_x)\) in Section~\ref{sec:coho12} is a quotient of \(M\). Its defining extensions over \(L_\infty\) are unramified outside two and split completely above two, hence are everywhere unramified. The quotient map respects the action of this subgroup. On imposing the specialization relations of Proposition~\ref{prop:spec12} we therefore have a surjection
\[
  \operatorname{coker}\!\left(A^{2^{h_x}}-u^{-r2^{h_x}}I\right)
                    \longrightarrow X'_r .
\]
Imposing one generator relation suffices for this surjection because all the relations defining \(X'_r\) are imposed in its target.

Let \(\alpha_1,\ldots,\alpha_D\) be the eigenvalues of \(A'\), with multiplicity, in \(\overline{\QQ}_2\), normalizing \(v_2(2)=1\). Each satisfies \(v_2(\alpha_i-1)\ge2\), since \((A'-I)/4\) is an integral matrix. Put \(b_r=u^{-r2^{h_*}}\), also congruent to one modulo four. At each squaring step \(v_2(a+b)=1\) for two such units, so \(v_2(a^2-b^2)=v_2(a-b)+1\). The expression also holds with infinite valuation at equality. Applying this to each eigenvalue, without any diagonalizability assumption, gives
\[
 v_2\!\det\!\left(A^{2^{h_x}}-u^{-r2^{h_x}}I\right)
       = D(h_x-h_*)+\sum_{i=1}^D v_2(\alpha_i-b_r).
\]

Put \(T_Y=\lfloor(2Y)^{(3/2)401}\rfloor\), \(J_Y=\lceil\log_2(T_Y+1)\rceil\).
For distinct \(1\le r,s\le T_Y\),
\[
  v_2(b_r-b_s)= e_*+h_*+v_2(r-s)
                         \le e_*+h_*+J_Y .
\]
The equality is the principal-unit binomial identity proved in Lemma~\ref{lem:det12}. Consequently, for each \(\alpha_i\), at most one integer \(1\le r\le T_Y\) has \(v_2(\alpha_i-b_r)>e_*+h_*+J_Y\). Otherwise the ultrametric inequality would contradict the last display. Collect these integers in \(\mathcal R_Y\). The choices of \(A'\), \(u\), and \(T_Y\) precede the choice of \(x\). The same \(\mathcal R_Y\) thus works throughout \(I_Y\).

For every remaining twist the determinant is nonzero and its valuation is at most \(D(h_x+e_*+J_Y)\). This is \(O_{\psi,\eps}(\log(2x))\), since \(x\ge2Y\). The elementary divisor theorem bounds the length of \(X'_r\) by that valuation. Proposition~\ref{prop:spec12} adds at most one for the kernel of restriction, proving both conclusions.
\end{proof}

The exceptional set gives a precise limitation of the spacing argument. Even a unit \(\alpha\) arbitrarily close to a single \(b_r\) satisfies all of the pairwise spacing inequalities. Proposition~\ref{prop:finite15} proves finiteness at that twist. Theorem~\ref{thm:trunc15} bounds its subgroup of bounded exponent. Sections~\ref{sec:common16} and~\ref{sec:dyadic17} combine a shorter interval, descent, and an effective bounded-degree estimate to handle their full lengths. At odd primes the missing terms are the complementary sign from Section~\ref{sec:critical13} beyond the prime or exponent cutoff of Theorem~\ref{thm:trunc15}.

\subsection{The remaining twists on a common power interval}\label{nav:II-10}
\label{sec:common16}

The proof of Theorem~\ref{thm:exception15} leaves at most \(R(\psi)\) dyadic twists. We now descend their full local lengths to a field of bounded degree and choose one interval on which every such twist that actually occurs is small. Both conclusions concern the groups in \eqref{eq:length12} themselves.

\begin{proposition}[Descent without a torsion error]\label{prop:desc16}
Use \(B,L_\infty,\Gamma_*,h_*,h_x\) from Section~\ref{sec:two15}. There is a field \(F_*\subset L_\infty\) containing \(B(i)\), of degree bounded solely in terms of \(\psi\), such that for every \(r\ge1\) and every sufficiently large \(x\) with \(h_x\ge h_*\),
\[
 \len_{\ZZ_2}\mathcal S_{2,r}(E_x)
 \le 1+ R(\psi)(h_x-h_*)+
                   \len_{\ZZ_2}\mathcal S_{2,r}(F_*).
\]
Moreover \(\log(2|\operatorname{disc}F_*|)\ll_\psi\log(2Fqd)\).
\end{proposition}
\begin{proof}
Take \(F_*\) to be the fixed field of \(\gamma^{2^{h_*}}\), meaning the closed subgroup it generates in \(\Gamma_*\). Its index over \(B(i)\) is \(2^{h_*}\), hence its degree is bounded. The cyclotomic image of this subgroup is \(1+2^{e_*+h_*}\ZZ_2\). Thus \(F_*=B(i,\mu_{2^{e_*+h_*}})\) and its conductor divides \(2^{e_*+h_*}Fqd\). The bounds for \(e_*,h_*\) and the conductor--discriminant formula give the discriminant estimate.

Let \(X\) be the maximal abelian unramified pro-\(2\) Galois group over \(L_\infty\) on which all decomposition groups over two vanish. This is an intrinsic quotient of \(M\) from Lemma~\ref{lem:generators15}. It is the group \(X'_2\) in Section~\ref{sec:coho12} for each of the base fields \(F_*,E_x\). To see this, at any \(v\nmid2\) the cyclotomic extension of either base is unramified, so its absolute inertia subgroup at \(v\) is unchanged on passage to \(L_\infty\). An extension cut out by a quotient defining \(X\), together with all its conjugates over that base, kills these inertia subgroups. It is therefore contained in the base field's maximal extension unramified away from two. Conversely, the extensions defining \(X'_2\) are unramified there and split over two. These descriptions give the same quotient and conjugation action.

Both \(F_*\) and its dyadic completions contain \(i\). Their cyclotomic Galois groups are open subgroups of \(1+4\ZZ_2\), hence are procyclic and isomorphic to \(\ZZ_2\). The degree-one inflation and degree-two obstruction in Proposition~\ref{prop:spec12} then vanish, exactly as in its procyclic calculation. Locally the same holds at every prime over two. Consequently, with \(b_r=u^{-r2^{h_*}}\),
\[
 \mathcal S_{2,r}(F_*)^\vee
             \simeq X/(\gamma^{2^{h_*}}-b_r)X .
\]
For \(E_x\) the length of its strict group is at most one plus the length of \(X/(\gamma^{2^{h_x}}-u^{-r2^{h_x}})X\), since the relations from the full cyclotomic Galois group include this one.

Apply the minimal-generator construction in Theorem~\ref{thm:exception15} to \(X\) itself. Write \(\pi:\ZZ_2^D\to X\) for the surjection and \(A\) for a lift of \(\gamma\). Now \(D\le R\), and \(A'=A^{2^{h_*}}\equiv I\bmod4\) as before. Let \(a=h_x-h_*\). Factoring a difference of \(2^a\)-th powers gives
\[
 (A')^{2^a}-b_r^{2^a}I
    =2^a(A'-b_r I)
       \prod_{j=0}^{a-1}\frac{(A')^{2^j}+b_r^{2^j}I}{2}.
\]
Every factor in the product is an integral matrix congruent to \(I\bmod2\), hence invertible over \(\ZZ_2\). Its product therefore maps the entire free cover onto itself. After applying \(\pi\), the image of the left hand side is exactly \(2^a Z\) with \(Z=(\gamma^{2^{h_*}}-b_r)X\). This uses equality of images on a free cover and holds even if \(X\) has finite torsion. Now \(Z\) has at most \(D\) generators since it is an image of \(X\). The kernel of \(X/2^aZ\to X/Z\) is \(Z/2^aZ\) of length at most \(aD\). The length of \(X/Z\) is finite by Proposition~\ref{prop:finite15} and its proof, which applies to the totally imaginary field \(F_*\). This proves the proposition, including \(D=0\).
\end{proof}

\begin{theorem}[Small exceptional twists on the interval used for the count]\label{thm:common16}
Fix \(0<\eta<401\). For arbitrarily large \(Z\) there exist primes \(q,d\), a field \(F_*\) as above, a real \(V\ge Z\), and at most \(R(\psi)\) exceptional integers such that, simultaneously for every \(V\le x\le V^{3/2}\), the following hold.
\begin{enumerate}
\item The fields, characters, and integral prime-sum weights satisfy conditions~\ref{cond:sizes}--\ref{cond:weights}. The odd-prime bound of Theorem~\ref{thm:critical13} and the finite-exponent bound of Theorem~\ref{thm:trunc15} hold there.
\item At two the length is \(O_{\psi,\eps,\eta}(\log(2x))\) for each nonexceptional \(r<t(x)\). Every exceptional \(r<t(x)\) satisfies \(r\le x^\eta\) and the bound of Proposition~\ref{prop:desc16} with this single \(F_*\).
\end{enumerate}
\end{theorem}
\begin{proof}
Put \(c_0=401\), \(b_1=c_0/\eta\), \(R=R(\psi)\), and fix \(K=2(2b_1)^R\) before choosing \(Z,q,d\). Choose the primes by Theorem~\ref{thm:twists} with initial scale \(Z\) and power \(K+1\). In the range \(Z\le x\le Z^K\) the field is \(E_x=BC_l\) with the same \(q,d\), and \(\log(qd)\ll Z^{\eps/5}(\log Z)^2\). All conductor, degree, and dyadic generator bounds used in Sections~\ref{sec:critical13}--\ref{sec:two15} are consequently uniform in this larger range. In Theorem~\ref{thm:exception15} make the single replacement \(T_Y=\lfloor Z^{Kc_0}\rfloor\). Its proof produces at most \(R\) exceptional integers for this entire range. The inequality \(J_Y\ll_{\psi,\eps,\eta}\log(2x)\) follows from \(x\ge Z\).

Here is the interval selection in detail. Starting with \(U=K\log Z\), attach to each exceptional integer \(r\) the point \(a_r=(\log r)/c_0\). If none lies in \((U/(2b_1),U]\), take \(v=U/2\). For \(v\le u\le2v\) and \(a_r\le u\), one then has \(b_1a_r\le v\le u\). Otherwise discard one point in that half-open interval, replace \(U\) by \(U/(2b_1)\), and repeat. A discarded point exceeds every future upper endpoint. After at most \(R\) divisions the procedure terminates. For the resulting \(v\), \(v\ge K\log Z/(2(2b_1)^R)=\log Z\) and \(2v\le K\log Z\). Take \(V=e^v\). On \([V,V^2]\) this gives
\[
             r\le x^{401}\ \Longrightarrow\ r\le x^\eta
                   \qquad(r\text{ exceptional}).
\]
Use the subinterval \([V,V^{3/2}]\).

At its left endpoint it is convenient to choose \(a_t=\lceil bx\rceil\) in Proposition~\ref{prop:weights14}. This makes \(x\le a_t/b\le x+1/b\), so the character estimates apply even if \(V=Z\). At the upper endpoint \(a_t/b\le Z^K+1\le Z^{K+1}\). The proof for integral weights, including its derivative bound, applies to this rounding. All the size, count-weight, and cohomology deductions cited in the statement now hold on this one interval. The primes and \(F_*\) were chosen before restricting the interval.
\end{proof}

In particular, the dyadic remainder in this construction can be bounded by the sum of at most \(R(\psi)\) orders on a bounded-degree field, counted logarithmically. Explicitly, putting \(v_2(2)=1\), one obtains uniformly
\[
 \sum_{1\le r<t(x)}\len_{\ZZ_2}\mathcal S_{2,r}(E_x)
 \ll_{\psi,\eps,\eta} t(x)\log(2x)
   +\sum_{\substack{r\ {\rm exceptional}\\1\le r\le x^\eta}}
           v_2\!\left(|K_{2r}(\mathcal O_{F_*})|\right).
\]
This follows by summing Proposition~\ref{prop:desc16} and applying Proposition~\ref{prop:finite15} at \(F_*\). The dependence on \(\operatorname{disc}F_*\) still matters here because the field changes with the chosen auxiliary primes.

\subsubsection{The implication still needed for Dirichlet GRH}\label{nav:II-11}
The prime sum itself detects whether a period comparison on this family would finish the proof. Take \(T_2=A_0\), \(m=\lfloor\sqrt x\rfloor\), \(a_E=xb\), \(k=xm\), the integral weights of Proposition~\ref{prop:weights14}, and \(a_t=\lceil bx\rceil\) for all \(t\). These are the numerical choices in Theorem~\ref{thm:common16}. An interval in the following proposition always uses one pair \(q,d\).

\begin{proposition}\label{prop:actual16}
With \(Q=hmb\) and \(\delta=1/o\), at each \(x\) on the selected interval and for every translate \(u\) one has
\[
       \frac{2\Delta_u}{Q\delta}
          =\Gg_\psi(x)+O_{\psi,\eps}(x^{1/2+\eps/2}).
\]
Consequently, to prove Main Theorem~\ref{thm:grh16} it suffices to establish \eqref{eq:gap} for these counts on arbitrarily distant selected intervals for every primitive \(\psi\) and arbitrarily small \(\eps>0\). Conversely, Main Theorem~\ref{thm:grh16}, if established, would imply \eqref{eq:gap} for these numerical choices on every sufficiently distant selected interval.
\end{proposition}
\begin{proof}
Every distinguished set \(L_u\) is the same because \(T_2=A_0\). The idealized contribution for each translate therefore equals its average in Lemma~\ref{lem:av10}. The exact coset proportions of Proposition~\ref{prop:field12} give the weights of \(\Gg_\psi(y)\) on the retained characters, with \(y=\lceil bx\rceil/b\). The other \(3o\) surviving characters have weights of modulus at most one after normalizing the sum over \(T\). Their total is at most \(3o\,y^{1/2+\eps/2}\) by Theorem~\ref{thm:twists}. The comparison of actual and idealized weights in condition~\ref{cond:weights} costs at most \(4/(\delta x)\) after multiplication by \(2/(Q\delta)\). Finally \(0\le y-x\le1/b\), and the slope bound in the proof of Proposition~\ref{prop:balance} gives \(|\Gg_\psi(y)-\Gg_\psi(x)|\ll_\psi(1+\log^2(2x))/b\). This proves the formula.

The sufficiency statement is Proposition~\ref{prop:balance} followed by Theorem~\ref{thm:criterion16}. In the other direction, the critical-line assertion would give \(|\Gg_\psi(x)|\ll_\psi\sqrt x\) by Theorem~\ref{thm:explicit}. Since \(\mathcal B/Q\ge h a_E/(hmb)=x/m\ge\sqrt x\), the formula in the proposition would yield \(\Delta_u\le x^{5\eps}\mathcal B\) for every large \(x\) on those intervals.
\end{proof}

Even for the primitive character of conductor one, the retained term here is the smoothed prime sum for the Riemann zeta function. A proof of \eqref{eq:gap} must bound that term independently of its zeros. In the proposed theta argument the local divisibility must measure the Frobenius-weighted count for the specified projected period maps, and the norm and index estimates must hold for the same maps. The endpoint comparison in Theorem~\ref{thm:end14} and the local length estimates above establish parts of this argument. The required upper bound for these projected periods and the remaining part of \eqref{eq:length12} have not been derived. Section~\ref{sec:dyadic17} completes the dyadic part of the local cohomology estimate.

\subsection{The full local-length bound at the prime two}\label{nav:II-12}
\label{sec:dyadic17}

We now bound the exceptional twists of Theorem~\ref{thm:common16}. The necessary input is an estimate in terms of the degree and discriminant of \(F_*\). Its degree is bounded before the auxiliary primes are chosen. Write \(D_L=|\operatorname{disc}L|\) for any number field \(L\).

\begin{lemma}[A dyadic bound at bounded degree]\label{lem:coarse17}
Fix \(d_0\ge2\). If \(L\) is totally imaginary of degree \(d\le d_0\), then for all \(r\ge1\)
\[
 \len_{\ZZ_2}\mathcal S_{2,r}(L)
       \le \exp\!\left(C_{d_0}(r+2)^8\log(2(r+2)D_L)\right).
\]
\end{lemma}
\begin{proof}
Set \(i=2r\), \(n=2i+8\), \(\Gamma=\operatorname{GL}_n(\mathcal O_L)\), and \(D=dn^2\). We first bound an integer annihilating the torsion in \(H_i(\Gamma,\ZZ)\). Use the contractible well-rounded retract \(W\) of Ash for the free lattice \(\mathcal O_L^n\). Its cells are indexed by the sets of shortest nonzero vectors for positive archimedean Hermitian metrics, normalized to have squared trace minimum one. A face has a larger set of shortest vectors than a cell containing it, with equality only for the same cell. Stabilizers are finite. The metrics form an open cone in a real vector space of dimension at most \(D\). The well-rounded condition is that the shortest vectors span \(L^n\) over \(L\).~\cite[Sec.~2.3]{II-II-II-beh17}

Bayer-Fluckiger, Emery and Houriet give a finite set \(\Phi\subset\mathcal O_L^n\) into which some \(\gamma\in\Gamma\) simultaneously moves all shortest vectors of any one well-rounded form. Their Corollary 5.2 applies to free lattices of rank at least five over a field other than \(\QQ\), as here. Its explicit estimate implies
\[
    \log(2+|\Phi|)\le C_{d_0}(n+1)^3\log(2nD_L).
\]
We use this assertion about shortest vectors, which has no restriction on the residue characteristic.~\cite[Cor.~5.2]{II-II-II-beh17}

We count chains of incident cells. For a vector \(v\), the equation that it has squared trace length one is affine linear in the metric. For the shortest-vector set of any cell, the corresponding simultaneous affine equations are consistent and are implied by at most \(D+1\) of them, by linear algebra. Such a choice from its shortest vectors determines the cell among cells for which it spans the simultaneous equations. Indeed, if two cells yield the same choice, every shortest vector of either has length one for a form in the other. Normalization then makes their full shortest-vector sets equal. In a chain of cells take its smallest face and move its shortest vectors into \(\Phi\). After this single move, every shortest-vector set in the chain is contained in \(\Phi\). Each member can therefore be encoded by a list of at most \(D+1\) vectors from \(\Phi\). A strict chain has at most \(D+1\) cells by dimension. Allow a padding symbol in the lists and for missing cells. The number of \(\Gamma\)-orbits of simplices of the barycentric subdivision of \(W\) is at most \((2+|\Phi|)^{(D+2)^2}\). This argument permits arbitrary incidences among the shortest-vector equations.

Let \(\Gamma_3\) be the kernel of reducing the underlying \(\ZZ\)-linear action on \(\mathcal O_L^n\) modulo three. Its index \(J\) is at most \(3^{(dn)^2}\) and it is torsion-free by the binomial argument in Theorem~\ref{thm:end14}. Thus it acts freely on the contractible subdivided retract. Its quotient computes integral group homology. This quotient has at most
\[
                  C=J(2+|\Phi|)^{(D+2)^2}
\]
oriented simplices when one orientation is chosen for each simplex. Indeed a \(\Gamma\)-orbit splits into at most \(J\) orbits under the subgroup. Equivalently one may use the quotient as a simplicial cell complex, allowing face identifications. Each column of the boundary in degree \(i+1\) has sum of absolute entries at most \(i+2\), from the alternating sum of its faces. If the boundary matrix has rank \(k>0\), the order of the torsion in its cokernel divides the absolute value of any nonzero \(k\)-rowed minor by Smith normal form. Hadamard's inequality bounds such a minor by \((i+2)^k\). The rank-zero case contributes order one. Homology embeds in the cokernel. Consequently \(T=|H_i(\Gamma_3,\ZZ)_{\rm tors}|\) satisfies \(\log T\le C\log(i+2)\). Transfer followed by inclusion is multiplication by \(J\) on the homology of \(\Gamma\). It follows that \(JT\) kills every torsion class in \(H_i(\Gamma,\ZZ)\). In particular
\[
       \log(JT+1)
             \le \exp\!\left(C_{d_0}'(n+1)^7\log(2nD_L)\right),
\]
since \(D=dn^2\).

Homological stability for general linear groups over a Dedekind ring identifies \(H_i(\Gamma,\ZZ)\) with \(H_i(\operatorname{GL}(\mathcal O_L),\ZZ)\) in this range. The stable range \(n\ge2i+1\) used for rings of integers in Section 3.1 of the cited paper suffices. Quillen's plus construction gives the same integral homology for the connected space \(B\operatorname{GL}(\mathcal O_L)^+\) and identifies its \(i\)-th homotopy group with \(K_i(\mathcal O_L)\).~\cite[Sec.~3.1]{II-II-II-beh17}\cite{II-II-II-weibel15}

We also need the Hurewicz kernel at two. Mathew's Theorem 1.3(a) says that \(2^{\lceil i/2\rceil+3}\) kills the dyadic part of the kernel of the stable Hurewicz map of any connective spectrum in positive degree \(i\). Apply it to the connective algebraic \(K\)-theory spectrum. The space \(B\operatorname{GL}(\mathcal O_L)^+\) is its basepoint component in the infinite loop space. The adjunction map from the suspension spectrum of that component to the \(K\)-theory spectrum takes its ordinary Hurewicz image to the stable Hurewicz image. Indeed both are obtained from a representative sphere map by applying integral homology, followed in the second case by this adjunction map. Hence a dyadic homotopy class killed by ordinary Hurewicz is killed by \(2^{r+3}\). As \(K_{2r}(\mathcal O_L)\) is finite, its ordinary Hurewicz image is torsion. We have proved that \(2^{r+3}JT\) annihilates its dyadic part.~\cite[Thm.~1.3(a)]{II-II-II-mathew17}

By the duality and comparison in the proof of Proposition~\ref{prop:finite15}, the same integer kills the strict group. To bound its length from its exponent, its dimension modulo two, equivalently its dimension of elements killed by two, is also needed. The Kummer argument of Theorem~\ref{thm:trunc15} at \(\ell=2\) applies directly to \(L\), since \(\mu_2\subset L\). Dirichlet's unit theorem bounds the unit contribution by \(2d\). Minkowski's bound and the estimate \(\zeta_L(2)\le\zeta(2)^d\) bound \(\log|{\rm Cl}(L)|\) by \(O(d+\log(2D_L))\), including for nonabelian \(L\). The number of cyclic summands is therefore \(O_{d_0}(\log(2D_L))\). Multiplying this by \(r+3+\log_2(JT)\) bounds the length. Since \(n=4r+8\), the asserted exponent eight absorbs the last factors.
\end{proof}

\begin{theorem}[All dyadic twists on the common interval]\label{thm:two17}
For every fixed primitive \(\psi\) and \(0<\eps<1/52\), there are arbitrarily distant intervals \([V,V^{3/2}]\) and a single pair of auxiliary primes on each interval for which, at every \(x\) therein,
\[
 (\log2)\sum_{1\le r<t(x)}\len_{\ZZ_2}\mathcal S_{2,r}(E_x)
                   \le \tfrac14\exp(x^\eps/2).
\]
On these same fields and intervals the bound \eqref{eq:critical13} holds at every odd prime after summing the specified eigenspace, and the integral prime-sum weights satisfy conditions~\ref{cond:sizes}--\ref{cond:weights}.
\end{theorem}
\begin{proof}
Use Theorem~\ref{thm:common16} with \(\eta=\eps/64\). The degree of \(F_*\) is bounded in terms of \(\psi\). Its discriminant satisfies \(\log(2D_{F_*})\ll_\psi\log(2Fqd)\ll_{\psi,\eps} x^{\eps/4}\) throughout the selected interval by Proposition~\ref{prop:desc16} and the bounds for \(q,d\). For an exceptional \(1\le r<t(x)\), the interval choice gives \(r\le x^\eta\). The exponent in Lemma~\ref{lem:coarse17} is then at most \(C_{\psi,\eps}x^{3\eps/8}\), using \(8\eta=\eps/8\). Apply the lemma to the strict group over \(F_*\) in Proposition~\ref{prop:desc16}. The descent contributes \(O_\psi(\log(2x))\) for each such twist. There are at most \(R(\psi)\) of them. For every other twist Theorem~\ref{thm:common16} gives \(O_{\psi,\eps}(\log(2x))\) directly. Thus the sum before its factor \(\log2\) is bounded by
\[
 C_{\psi,\eps}\!\left(t(x)\log(2x)
                     +\exp(C_{\psi,\eps}x^{3\eps/8})\right).
\]
As \(t(x)\le x^{401}\) and \(3\eps/8<\eps\), this gives the stated quarter of the allowance for all sufficiently large \(V\). All the other assertions were included in the selection theorem.
\end{proof}

The full dyadic term in \eqref{eq:length12} and the odd part in \eqref{eq:critical13} now use at most half its allowance. The remaining cohomological task is to bound the \((-1)^r\)-eigenspace for odd primes beyond the ranges in Theorem~\ref{thm:trunc15}. The geometric task is to prove \eqref{eq:gap} for the projected integral theta periods on the same interval. Proposition~\ref{prop:actual16} shows why that estimate, if established, would imply the intended main theorem.

\subsection{Local cohomology at the primes of the theta tests}\label{nav:II-21}
\label{sec:relax183}

Write \(E=E_x\), \(F=E^+\), \(f=[F:\QQ]\), and
\(d=[E:\QQ]=2f\) for the common family of Section \ref{nav:I-28}. Let \(N\) be the cyclotomic modulus there.
The two split prime ideals \(\mathfrak p,\mathfrak q\)
in the theta construction represent the narrow classes of
the absolute different \(\mathfrak d_F\) and the ideal
\(\mathfrak a_0\) below \((\sqrt{-d_0})\).
We shall use the strict cohomology groups of the original
field throughout. The auxiliary class fields below enter
only in the choice of the two primes.

\begin{proposition}[Split primes above a prescribed threshold]
\label{prop:high183}
Let \(R_0\) be the product of the distinct rational primes
dividing \(2Ng_1g_2T\). For every real \(Y\ge2\), the two
prime ideals may be chosen above distinct rational primes
\(\ell_1,\ell_2>Y\), outside \(R_0\), which split
completely in \(E\). With an absolute implied constant,
\[
 \log\ell_1+\log\ell_2\ll
      f^2\big(1+\log|D_F|+\log(2Y)
         +\log(2R_0)(1+\log(2f))\big).
\]
In particular their logarithmic sizes are polynomial in
\(x\) uniformly whenever \(\log(2Y)\) is so bounded.
\end{proposition}
\begin{proof}
First consider one of the two classes. Write \(U\) for the
set of rational primes to avoid, initially those dividing
\(R_0\). Put \(R_U=\prod_{p\in U}p\). Take \(M=\prod_{p\in U}
p^{k_p}\) as in
Proposition \ref{prop:cost177},
with \(k_p\) the least positive integer such that
\(\varphi(p^{k_p})>f\). Thus
\(\log M\le\log R_U+|U|\log(2f)\).
The extension \(L_0=H^+(F)F(\zeta_M)\) is abelian over
\(F\) and ramifies above all the excluded primes.

Choose any prime \(l\mid (R_U+1)\), so \(l\notin U\),
\(l\nmid N\), and \(\log l\le\log(2R_U)\).
Let \(a\ge1\) be least such that \(l^a>Y+1\).
Then \(\log(l^a)\le\log(Y+1)+\log l\).
At every place of \(F\) above \(l\) the field \(F\) is
unramified over \(\QQ_l\). The totally ramified local
cyclotomic extension \(\QQ_l(\zeta_{l^a})/\QQ_l\) remains
totally ramified of full degree after this unramified
base change. Hence \(F(\zeta_{l^a})/F\) has the entire
group \((\ZZ/l^a\ZZ)^\times\) as its inertia there.
The extension \(L_0/F\) is unramified at \(l\) since the narrow
Hilbert class field is unramified at finite places and
\(l\nmid M\). Their intersection is \(F\), since every
nontrivial subextension of \(F(\zeta_{l^a})/F\) is
ramified there. Thus we may prescribe a lift of the narrow
class on \(L_0\) and prescribe \(-1\bmod l^a\) independently
on \(F(\zeta_{l^a})\).

The finite conductor norm of any character on the composite
is at most \((Ml^a)^f\), by the character extension and
local norm calculation in Proposition \ref{prop:cost177}.
The unconditional theorem of
\href{https://arxiv.org/pdf/1604.01750}{Thorner and Zaman,
Theorem 1.1} gives an unramified degree-one prime of \(F\)
in the prescribed Frobenius class, with
\[
    \log\ell\ll1+\log|D_F|+
               f(\log M+\log(l^a)+\log(2f)).
\]
Here \(\ell\) is its rational prime norm. The action on
roots of unity forces \(\ell\equiv-1\pmod{l^a}\), so
\(\ell>Y\). The rational prime avoids \(U\), and in
particular \(N\), hence is unramified in \(E/\QQ\).
The chosen narrow class restricts trivially to \(E/F\),
by Lemma \ref{lem:good175}. The prime of \(F\) therefore
splits in \(E\). As \(E/\QQ\) is Galois, the existence of
one prime there of absolute residue degree one proves
complete splitting of \(\ell\).

For the second class add \(\ell_1\) to \(U\) and repeat.
At this step \(\log R_U=\log R_0+\log\ell_1\) and
\(|U|\le1+\log R_0/\log2\). Substitution of the first
bound into the second yields the displayed estimate. All
original family parameters on its right have logarithms
or degrees polynomially bounded as asserted.
\end{proof}

For a number field \(K\), a finite place \(w\) and a
coefficient prime \(p\), put
\(A_{p,r}=(\QQ_p/\ZZ_p)(-r)\) for \(r\ge1\). Write
\(G_w\) and \(I_w\) for the local Galois and inertia groups.
Use arithmetic Frobenius and write \(q_w=Nw\). Say that a
local class is unramified when its restriction to \(I_w\)
vanishes. Impose the original archimedean conditions in
every group. In our application \(E\) is totally imaginary.

\begin{lemma}[Local restrictions at a negative twist]
\label{lem:loc183}
If \(w\nmid p\), the unramified part of \(H^1(K_w,A_{p,r})\)
is zero and the whole local group is cyclic of order
\(p^{v_p(q_w^{r+1}-1)}\). At \(w\mid p\), writing
\(e_w=e(K_w/\QQ_p)\), the unramified part has logarithmic
order at most
\[
        (2+v_p(r)+v_p(e_w))\log p .
\]
At \(K_w=\QQ_p\) with \(p\) odd and \(p-1\nmid r\),
this unramified group is zero.
\end{lemma}
\begin{proof}
Away from \(p\) the coefficient is unramified.
The quotient of the decomposition group by inertia is
procyclic, of \(p\)-cohomological dimension one.
Inflation--restriction computes its degree-one term as
\(A_{p,r}/(q_w^{-r}-1)A_{p,r}\). The multiplier is a
nonzero \(p\)-adic integer and the module is divisible,
so this group vanishes. The full local cohomology is
thus \(\operatorname{Hom}_{\rm cont}(I_w,A_{p,r})^{\Frob_w}\).
The wild inertia and prime-to-\(p\) factors of tame inertia
contribute zero. On the pro-\(p\) factor \(\ZZ_p(1)\),
Frobenius acts as multiplication by \(q_w\). The induced
action on its homomorphisms to \(A_{p,r}\) is therefore
multiplication by \(q_w^{-r-1}\). Its fixed group is the
torsion killed by \(q_w^{r+1}-1\) in \(\QQ_p/\ZZ_p\), giving
the stated cyclic group.

At \(w\mid p\), unramified cohomology is a quotient of
\(A_{p,r}^{I_w}\), by the same procyclic computation.
The image of the local cyclotomic character on inertia
over \(\QQ_p\) is all of \(\ZZ_p^\times\), since the
\(p\)-power cyclotomic extensions are totally ramified.
Its image on \(I_w\) has finite index \(h\) dividing \(e_w\).
Indeed the index of \(I_w\) in the absolute inertia of
\(\QQ_p\) is \(e_w\), and passing to images gives a
divisor of that index. For odd \(p\), the element
\((1+p)^h\) belongs to the image. Invariance implies
annihilation by \((1+p)^{rh}-1\), of valuation
\(1+v_p(rh)\). At two use \(5^h\) and valuation
\(2+v_2(rh)\). These valuations follow by binomial
expansion, first at an exponent prime to \(p\) and then
at each successive multiplication of the exponent by \(p\).
The order of the unramified quotient has the asserted
upper bound. If \(K_w=\QQ_p\) and \(p-1\nmid r\), a
Teichm\"uller root of unity in inertia acts on the coefficient
by a scalar whose difference from one is a unit. There
are then no inertia invariants.
\end{proof}

Let \(R\) be any squarefree product of rational primes
and let \(W\) consist of all places of \(E\) over its
divisors. Write \(\Sel_{p,r}(E)\) for the usual strict
group. Its classes are unramified away from \(p\) and
vanish on every decomposition group above \(p\). This is
the local condition in
\srcloc{Formalism/GaloisCohomology/Statement}{75}.
Define \(\Sel_{p,r}^{W}(E)\) by allowing the full local
group at \(w\in W\), \(w\nmid p\), and unramified classes
at \(w\in W\), \(w\mid p\). Keep the strict conditions
at all other places.

\begin{theorem}[The cost of the finite test places at every twist]
\label{thm:relax183}
For each integer \(t\ge2\) one has
\begin{equation}
  0\le \sum_{1\le r<t}\sum_p
       \log |\Sel^W_{p,r}(E)/\Sel_{p,r}(E)|
       \ll d\, t^2\log(2Rdt) .                          \label{eq:relax183}
\end{equation}
Every original strict class already restricts to zero at
every finite place. At odd primes the same upper bound
holds for the total increase in the logarithmic orders
after applying the sign \((1+(-1)^r c)/2\) and the operators
\(p^{k_{p,r}}(1-\pi_{p,r})\) of equation \ref{eq:tail114}.
\end{theorem}
\begin{proof}
Away from \(p\) a strict local class is unramified and
hence zero by the lemma. At \(p\) it is zero by definition.
Embed the global quotient in the direct sum of its
allowed local groups at \(W\) by restriction. The kernel
consists exactly of the original strict classes. The
quotient is finite since these local groups are finite.
For a fixed \(w\in W\) the sum over coefficient primes
different from its residue characteristic is at most
\((r+1)\log q_w\), by factoring the positive integer
\(q_w^{r+1}-1\). The sum of \(\log q_w\) over all \(W\)
is at most \(d\log R\), by the residue degree formula.
The resulting contribution over all twists is at most
\(d(t-1)(t+2)\log R/2\).

For the coefficient primes equal to the residue characteristic
use the second part of the lemma. For each rational prime
there are at most \(d\) places, all with the same
ramification index dividing \(d\), since \(E/\QQ\) is
Galois. Consequently at a single \(r\) the remaining
sum is at most \(d(2\log R+\log r+\log d)\), by factoring
\(r\) and \(d\). This proves \eqref{eq:relax183}, including
the coefficient prime two. In particular all prime sums
for these quotients are finite.

Both global groups are stable under \(\Gal(E/\QQ)\),
since \(W\) contains whole rational fibers. The sign and
the indicated averaging projectors are integral group-ring
operators at the coefficient primes where they are used.
For any such combined endomorphism \(T_0\), the map from
the quotient of the two groups onto the quotient of their
images under \(T_0\) is surjective, by applying \(T_0\)
to representatives. The order of the latter quotient is
therefore bounded by the former. Additivity of the
logarithmic order gives the last assertion.
\end{proof}

\begin{corollary}[The cohomological and theta choices on the same interval]
\leavevmode\label{cor:both183}
Take \(t=\lfloor x^{401}\rfloor\), \(b=\exp(x^\epsilon)\),
and \(Y\) at least \(t+2\) and the largest prime-power cutoff
in the count of Section \ref{nav:I-28}. The two auxiliary
primes can be chosen above this \(Y\) with logarithmic norms
polynomial in \(x\) if \(\log(2Y)\) is bounded by a power of
\(x\). Such a threshold exists on the original intervals.
For \(R=R_0\ell_1\ell_2\) the left side of \eqref{eq:relax183}
is \(O_{\psi,\epsilon}(x^C)\), hence
\(o(b^{1/2})\), there. The last real integral comparison of
Theorem \ref{thm:integral181}, and its finite-depth construction
continue to use these very fields and intervals.
\end{corollary}
\begin{proof}
Each cutoff is of the form \(a_j/w\) with positive prescribed
weight \(w\). The weights in the construction are eventually
at least one and its numerators are \(O(xb)\), so one may
take \(Y\ge\max(t+2,Cxb)\) for a sufficiently large fixed
\(C\). Its logarithm is bounded by a power of \(x\). The
prime selection gives the original narrow classes, so the
separately unramified oscillator data, the bounded generators
of Theorem \ref{thm:balanced179}, and the ensuing theta calculations
all apply. The two new rational primes split completely and
exceed the cutoffs. At their own coefficient characteristics
the relaxed unramified group is zero since \(p>t+2\).
The original exceptional primes and both auxiliary primes
enter the cohomological estimate through \(\log R\). This
quantity, \(d\), and \(t\) are all polynomially bounded in \(x\). This proves the corollary.
\end{proof}

The estimate bounds the change of local conditions at all
these primes. The classes already satisfying the original
strict conditions lie in the kernel of their restriction
to this entire set. Their remaining total order is the
global quantity in equation \ref{eq:tail114}.

\subsection{The remaining local estimate in terms of finite group orders}\label{nav:II-27}
Keep the pair of auxiliary primes and fields on the intervals of
Section~\ref{sec:missing}. Write \(E=E_x\), \(F_+=E^+\),
\(n_E=[E:\QQ]\), and \(t=\lfloor x^{401}\rfloor\). Let \(N\) be
the cyclotomic modulus of the field construction, so that each
character of \(\Gal(E/\QQ)\) has conductor dividing \(N\).
For a finite abelian group \(B\), write
\(|B|_{\rm odd}=|B|/2^{v_2(|B|)}\).
Consider the following sum of logarithmic orders of finite groups.
\begin{equation}
\begin{split}
 {\cal A}(E,t)
  ={}&\sum_{1\le 2j<t}\log |K_{4j}(\mathcal O_{F_+})|_{\rm odd}\\
     &+\sum_{\substack{j\ge1\\2j-1<t}}
        \log \frac{|K_{4j-2}(\mathcal O_E)|_{\rm odd}}
                       {|K_{4j-2}(\mathcal O_{F_+})|_{\rm odd}} .
\end{split}                                                    \label{eq:A100}
\end{equation}
Here \(j\) is an integer and \(K_i\) denotes Quillen's algebraic
\(K\)-group. All the quotients inside the logarithms are positive
integers, as the next theorem shows.

\begin{theorem}[The odd group orders required on the common interval]
\label{thm:res100}
Write \(\mathcal L(x)\) for the left side of \eqref{eq:length12}
and put \(\mathcal L_2(x)=\sum_{1\le r<t}\log|\mathcal S_{2,r}(E)|\).
On the intervals of Section~\ref{sec:missing} one has
\[
 \left|\mathcal L(x)-{\cal A}(E,t)-\mathcal L_2(x)\right|
                 \ll n_E^2t^2\log(4n_EtN).
\]
In particular, uniformly on those intervals as their left endpoints
tend to infinity,
\begin{equation}
        \frac{\mathcal L(x)-{\cal A}(E,t)}{\exp(x^\eps/2)}
                        \longrightarrow 0.                       \label{eq:small100}
\end{equation}
Thus any fixed \(\delta>0\) with
\({\cal A}(E,t)\le (1-\delta)\exp(x^\eps/2)\) throughout
arbitrarily distant such intervals suffices for \eqref{eq:length12}.
\end{theorem}
\begin{proof}
At an odd prime the action of complex conjugation on a finite
group splits into its two eigenspaces by the idempotents
\((1+c)/2\) and \((1-c)/2\). At every \(r\ge1\), restriction
identifies the plus eigenspace of
\(K_{2r}(\mathcal O_E)\) at odd primes with
\(K_{2r}(\mathcal O_{F_+})\) at those primes. The hypotheses for
this descent are satisfied because \(E/F_+\) is unramified at
finite places by Theorem~\ref{thm:norm91}. The descent proof and
the comparison with restricted Galois cohomology are in
Sections \ref{nav:III-12}--\ref{nav:III-13}.
In particular the ratio of the two odd orders is the order of the
minus eigenspace. The sign \((-1)^r\) is plus for \(r=2j\) and
minus for \(r=2j-1\). Adding its logarithmic orders produces
exactly \eqref{eq:A100}.

For clarity write \(U(E,t)\) for the sum of the logarithmic orders
of the opposite eigenspaces, with sign \((-1)^{r+1}\), at odd
primes. Let \(C_{\rm odd}(E,t)\) be the sum over \(r<t,\ r\ge1\)
and odd \(\ell\) of
\[
   \sum_{v\mid\ell}\log |H^0(E_v,(\QQ_\ell/\ZZ_\ell)(-r))|
            -\log |H^0(E,(\QQ_\ell/\ZZ_\ell)(-r))|.
\]
The exact order formula of Theorem \ref{thm:wild98}, gives the
identity
\[
       \mathcal L(x)={\cal A}(E,t)+\mathcal L_2(x)
                                  +U(E,t)-C_{\rm odd}(E,t).
\]
Theorem \ref{thm:critical99}, bounds \(U(E,t)\) by
\(O(n_Et^2\log(4n_EtN))\), and Corollary~\ref{cor:sum98} bounds
\(C_{\rm odd}(E,t)\) by \(O(n_E^2t^2\log(4n_Et))\).
Both terms are nonnegative. This proves the first bound.

On the selected intervals one has \(n_E\ll_\psi x^3\) and
\(\log(2N)\ll_{\psi,\eps}x\). Hence the right side is a
polynomial bound in \(x\). In the proof of
Theorem \ref{thm:two17},
the full dyadic term on these intervals is at most
\[
   C_{\psi,\eps}\bigl(t\log(2x)+\exp(C_{\psi,\eps}x^{3\eps/8})\bigr).
\]
The interval selection in
\hyperref[nav:volume-IV]{Volume~IV} retains this
bound. Since \(3\eps/8<\eps\), division by \(\exp(x^\eps/2)\)
sends it and each fixed polynomial in \(x\) to zero uniformly
as \(x\) increases. This proves \eqref{eq:small100} and the last
assertion.
\end{proof}

\begin{corollary}[The summand from the rational integers]
\label{cor:transfer100}
For every \(j\ge1\) and odd \(\ell\nmid n_E\), the group
\(K_{4j}(\ZZ)\{\ell\}\) embeds as a direct summand of the dual of
the actual strict group \(\mathcal S_{\ell,2j}(E)\).
Its image lies in the complementary sign at twist \(2j\).
For the primitive character of conductor one, this holds at every
odd prime. On each of its prescribed fields one consequently has
\[
          \sum_{1\le 2j<t}\log |K_{4j}(\ZZ)|_{\rm odd}
                                      \ \le\ \mathcal L(x).
\]
\end{corollary}
\begin{proof}
For \(\QQ\) at odd \(\ell\), the finite-coefficient Poitou--Tate
argument of Theorem \ref{thm:wild98}, applies with the one real place.
Its modified cohomology at this place vanishes because two is a
unit on all the coefficients. The cyclotomic characters of
\(\QQ\) and \(\QQ_\ell\) both have image \(\ZZ_\ell^\times\).
Indeed \(\Phi_{\ell^a}(1+X)\) is Eisenstein at \(\ell\) for every
\(a\ge1\), so that the full global degree is already obtained in
the completion. At twist \(-2j\) the global invariant subgroup
therefore equals the local one. Its restriction map is the
identity under these identifications. In the duality sequence the
map from the local degree-two group to the dual of the global
invariants is consequently an isomorphism. Thus every class of
\(K_{4j}(\ZZ)\{\ell\}\) has zero localization at \(\ell\).

Under extension of scalars to \(\mathcal O_E\) the number-ring
comparison commutes with restriction. Each localization above
\(\ell\) is the restriction of the zero localization just proved.
The image is therefore in the kernel identified with
\(\mathcal S_{\ell,2j}(E)^\vee\). It is fixed by \(c\), because
it comes from \(\ZZ\). This is sign \((-1)^{2j}\).

For the splitting, follow extension of scalars by the transfer that
forgets the \(\mathcal O_E\)-module structure. The ring
\(\mathcal O_E\) is free of rank \(n_E\) over \(\ZZ\). After a
choice of integral basis the composite functor on finitely
generated projectives over \(\ZZ\) is naturally isomorphic to
the direct sum of \(n_E\) copies of the identity. Additivity of
algebraic \(K\)-theory gives multiplication by \(n_E\) on
\(K_{4j}(\ZZ)\). This multiplication is invertible on the finite
\(\ell\)-primary subgroup. Composing its inverse with transfer
therefore retracts the strict subgroup onto the indicated copy.
For the character of conductor one, \(E\) is the compositum of
a real cyclotomic field of two-power degree and two quadratic
fields. Its degree is a power of two. The splitting applies to all
odd primes, and summing the logarithms proves the inequality.
\end{proof}

\begin{proposition}[The rational summand modulo an odd prime]
\label{prop:realclass100}
Let \(\ell\) be an odd prime, \(L=\QQ(\zeta_\ell)\), \(L^+=L\cap\RR\),
and \(\Delta=\Gal(L/\QQ)\). Denote the character describing its
action on \(\mu_\ell\) by \(\omega:\Delta\to\mathbb F_\ell^\times\).
Let \(V_\ell=\operatorname{Cl}(L^+)/\ell\operatorname{Cl}(L^+)\),
with the action of \(\Delta\) through restriction to \(L^+\).
For every \(j\ge1\),
\[
 \dim_{\mathbb F_\ell}K_{4j}(\ZZ)/\ell K_{4j}(\ZZ)
       =\dim_{\mathbb F_\ell} V_\ell^{\,\omega^{-2j}},
\]
where the superscript denotes the character eigenspace.
Consequently, on the intervals for the character of conductor one,
if \(J=\lfloor(t-1)/2\rfloor\), then
\[
       \mathcal L(x)\ \ge\
       \sum_{\ell>2\ {\rm prime}}
       \left\lfloor\frac{2J}{\ell-1}\right\rfloor
                (\log\ell)\dim_{\mathbb F_\ell} V_\ell .
\]
In particular the rational summands at a fixed \(\ell\) vanish for
every \(j\) exactly when \(\ell\) does not divide the ordinary
class number of \(L^+\).
\end{proposition}
\begin{proof}
Jordan, Klagsbrun, Poonen, Skinner and Zaytman give the rational
class-group formula in \cite[Thm.~5.1]{II-II-I-jkpsz100}. We derive its
even-character case from the same strict groups used in the
interval bound.
Write \(r=2j\), \(M=(\QQ_\ell/\ZZ_\ell)(-r)\),
and \(A=M[\ell]=\mathbb F_\ell(-r)\). The strict group over
\(\QQ\) is the Pontryagin dual of \(K_{4j}(\ZZ)\{\ell\}\) by the
first paragraph of the preceding proof. Its \(\ell\)-torsion
therefore has the dimension on the left of the first formula.
Multiplication by \(\ell\) on \(M\) gives the exact row
\[
 0\longrightarrow H^0(M)/\ell H^0(M)
  \longrightarrow H^1(A)\longrightarrow H^1(M)[\ell]
  \longrightarrow0
\]
in global restricted cohomology, and also in local cohomology
at \(\ell\). Restriction on the left is an isomorphism by equality
of the global and local cyclotomic images. The kernels of the two
other vertical maps are consequently isomorphic by the snake
lemma. Thus the \(\ell\)-torsion of the strict group is computed
exactly using the finite coefficient \(A\), including when \(A\)
is a trivial character.

On \(L\) the action on \(A\) is trivial. The extension \(L/\QQ\)
is unramified away from \(\ell\) and infinity. Its degree
\(\ell-1\) is a unit in \(\mathbb F_\ell\), so
inflation--restriction and averaging identify global classes
with \(\Delta\)-equivariant continuous homomorphisms to \(A\)
over \(L\), unramified outside \(\ell\). The identical
inflation--restriction argument works at \(\ell\), whose
decomposition group in \(\Delta\) is all of \(\Delta\).
Hence a strict homomorphism is trivial on the full decomposition
group above \(\ell\). It is unramified everywhere, and the prime
above \(\ell\) splits completely in the extension it defines.
By the Hilbert class field theorem its source is the ideal class
group of \(L\) modulo the class of that prime. The prime is
principal, generated by \(1-\zeta_\ell\). We have thus obtained
\(\operatorname{Hom}_\Delta(\operatorname{Cl}(L),\mathbb F_\ell(-r))\).

The group algebra of \(\Delta\) over \(\mathbb F_\ell\) splits
into its character eigenspaces, since \(\Delta\) is cyclic of
order \(\ell-1\). An equivariant homomorphism to \(A\) is an
arbitrary linear functional on the \(\omega^{-r}\)-eigenspace of
\(\operatorname{Cl}(L)/\ell\) and vanishes on the others.
As \(r\) is even, complex conjugation acts by \(+1\) on that
eigenspace. Ideal extension and norm between \(L^+\) and \(L\)
have composites \(2\) and \(1+c\), respectively. Division by two
on the odd-primary groups identifies the plus part with the
class group of \(L^+\) at \(\ell\), equivariantly for \(\Delta\).
This proves the dimension formula.

In every \((\ell-1)/2\) consecutive values of \(j\), the
characters \(\omega^{-2j}\) run once through all even characters
of \(\Delta\). Their dimensions sum to \(\dim V_\ell\).
Every copy of \(\mathbb F_\ell\) in the indicated reduction gives
at least \(\log\ell\) to the logarithm of the order of the
corresponding finite group. Apply Corollary~\ref{cor:transfer100}
and retain only complete blocks of \(j\le J\) for each prime.
The displayed sum is finite since its floors vanish for
\(\ell>2J+1\). A finite \(\ell\)-group is zero exactly when its
reduction modulo \(\ell\) is zero, which gives the last assertion.
\end{proof}

The estimates above express the remaining cohomology bound using
global orders on the specified fields. Equation~\eqref{eq:A100}
still needs a uniform upper bound of the size in
Theorem~\ref{thm:res100}. On the same intervals
Corollary~\ref{cor:alg97} still leaves the height of the
arithmetically normalized theta-period map to be controlled.

\subsection{The strict local condition at every coefficient level}\label{nav:III-16}
\label{sec:every113}

At the prime in the coefficients, a class which is zero with
divisible coefficients can have a nonzero finite lift. We compute
the permitted lifts over the field cut out by the finite character.
This also treats coefficient primes ramifying in the original field.

Let \(K\) be a number field, let \(p\) be an odd prime and let
\(O\) be the ring of integers of a finite unramified extension \(J\)
of \(\QQ_p\). Fix \(r\ge1\) and a continuous character \(\eta\) of
\(G_K=\operatorname{Gal}(\overline K/K)\) of finite order prime to
\(p\) with values in \(O^\times\). Write
\[
 \theta=\eta\kappa_p^{-r},\qquad D=(J/O)(\theta),\qquad
 A_j=D[p^j]\simeq (O/p^j)(\theta_j)\quad(j\ge1),
\]
where \(\kappa_p\) is the cyclotomic character and \(\theta_j\)
is the reduction of \(\theta\) modulo \(p^j\). The displayed
identification sends the class of \(p^{-j}a\) to \(a\bmod p^j\).
Let \(S_K(\theta)\subset H^1(K,D)\) consist of classes zero on
inertia at all finite \(v\nmid p\) and zero on \(G_{K_v}\) at
every \(v\mid p\). All cohomology and homomorphisms from Galois
groups are continuous. Set
\[
 F_j=\overline K^{\ker\theta_j},\qquad
 \Gamma_j=\operatorname{Gal}(F_j/K),\qquad R_j=O/p^j.
\]

\begin{lemma}[Restriction to the field of the character]\label{lem:cyclic113}
For each subgroup \(H\subset\Gamma_j\) the groups
\(H^1(H,A_j)\) and \(H^2(H,A_j)\) vanish. Thus restriction gives
\[
 H^1(K,A_j)\ \simeq\
   \operatorname{Hom}_{\Gamma_j}(G_{F_j}^{\rm ab},A_j).
\]
The same assertion holds on decomposition groups, with their own
finite images and local fields.
\end{lemma}
\begin{proof}
The finite image is contained in the product of the cyclic
Teichm\"uller group of the residue field of \(O\) and the image of
\(1+p\ZZ_p\) modulo \(p^j\). The second group is cyclic by the
\(p\)-adic logarithm, and the two orders are coprime. In particular
\(H\) is cyclic. For a generator acting by \(u\), with group order
\(n\), put \(T=1+u+\cdots+u^{n-1}\). The cyclic resolution has
alternate maps \(u-1,T\). If \(u-1\) is a unit it is exact in
positive degrees. If \(H\) is trivial the conclusion also follows.
In the remaining case \(u\) lifts to a rational principal unit with
\(a=v_p(u-1)\) satisfying \(1\le a<j\). Its exact order modulo
\(p^j\) is \(p^{j-a}\), since
\(v_p(u^{p^i}-1)=a+i\). The latter identity follows by the binomial
theorem and induction for odd \(p\). Hence \(v_p(T)=j-a\).
On \(R_j\) the kernel of \(u-1\) is \(p^{j-a}R_j=TR_j\), and
the kernel of \(T\) is \(p^aR_j=(u-1)R_j\).
This proves both vanishings. The inflation--restriction sequence
and the trivial action over \(F_j\) give the two assertions.
\end{proof}

For \(v\mid p\) with \(\bar\theta|_{G_{K_v}}=1\), define the
integer \(h_v\ge1\) by
\(\theta(G_{K_v})=1+p^{h_v}\ZZ_p\). Indeed the residual condition
cancels the Teichm\"uller part, leaving the principal part of
\(\kappa_p^{-r}\). The image of \(\kappa_p\) is open in
\(\ZZ_p^\times\) because \(K_v/\QQ_p\) is finite, so this principal
image is open. At each \(w\mid v\) of \(F_j\) use local reciprocity
to define
\begin{equation}
 \lambda_{w,j}(z)=
     p^{-\max(h_v,j)}\log\theta(\operatorname{rec}_{F_{j,w}}z)
            \pmod {p^{\min(h_v,j)}},\qquad z\in F_{j,w}^\times .
 \label{eq:log113}
\end{equation}
The formula uses the restriction of the global character to the
local abelian Galois group. Either choice of sign for reciprocity
gives the same kernel.

\begin{lemma}[The permitted local homomorphisms]\label{lem:local113}
At \(v\mid p\) with nontrivial residual character the finite lift of
a strict class must restrict to zero over \(F_{j,w}\). At a place
with trivial residual character, the allowed homomorphisms on
\(G_{F_{j,w}}^{\rm ab}\), viewed by local reciprocity on
\(F_{j,w}^\times\), are precisely the homomorphisms that factor
through \(\lambda_{w,j}\). This map is onto
\(\ZZ/p^{\min(h_v,j)}\). At \(v\nmid p\) the finite local
condition is vanishing on inertia over \(F_{j,w}\).
\end{lemma}
\begin{proof}
For any of the local groups \(G\), the kernel of
\(H^1(G,A_j)\to H^1(G,D)\) is the image of
\(D^G/p^jD^G\) under the connecting homomorphism. At \(v\nmid p\)
apply this to inertia. Its image under \(\theta\) is finite of
order prime to \(p\). Thus its invariants are either zero or all
of \(D\), and the kernel vanishes. Restriction to inertia over
\(F_{j,w}\) is injective on the cohomology of inertia downstairs
by Lemma~\ref{lem:cyclic113} and inflation--restriction. At
\(v\mid p\) with nontrivial residual character the full local
invariants vanish, which gives the first assertion in the same way.

At the other places over \(p\), put \(h=h_v\), \(s=\max(h,j)\)
and \(k=\min(h,j)\). The local invariants of \(D\) are
\(p^{-h}O/O\). A lift of \(p^{-h}a\) under multiplication by
\(p^j\) is \(p^{-h-j}a\). Over \(F_{j,w}\) its boundary in
\(A_j\) is the homomorphism
\[
 g\longmapsto a\,\frac{\theta(g)-1}{p^h}
       \equiv a\,p^{s-h}\frac{\log\theta(g)}{p^s}
                         \pmod {p^j}.
\]
The image of \(\theta\) over this field is exactly
\(1+p^s\ZZ_p\), the kernel of reduction modulo \(p^j\) on its
local image. If \(z\in p^s\ZZ_p\), then
\(\log(1+z)-z\in p^{2s}\ZZ_p\), term by term since
\((n-2)s\ge v_p(n)\) for \(n\ge2\) and \(p\) odd. As
\(2s-h\ge j\), the congruence follows. The logarithm maps this
principal unit group onto \(p^s\ZZ_p\). Local reciprocity
therefore gives all values modulo \(p^k\) in
\eqref{eq:log113}. A homomorphism from \(\ZZ/p^k\) to \(R_j\)
may send \(1\) to any element of \(p^{j-k}R_j\), and
\(j-k=s-h\). These are exactly the boundaries just computed.
The local restriction is injective by Lemma~\ref{lem:cyclic113},
so the condition on the restricted homomorphism is sufficient
as well.
\end{proof}

Let \(\Sigma\) be the primes of \(F_j\) above \(p\). Define
\(B_j\) to be the \(p\)-primary part of the following idele class
quotient of \(F_j\). Divide by \(F_j^\times\), the full factors at
infinity, and the local unit groups away from \(\Sigma\). At
\(w\in\Sigma\) divide by all of \(F_{j,w}^\times\) when
\(\bar\theta|_{G_{K_v}}\ne1\), and by \(\ker\lambda_{w,j}\)
otherwise. These subgroups are carried to one another by
\(\Gamma_j\), since the global character is invariant under
conjugation. If \(\Sigma_0\) consists of the latter primes, then
\begin{equation}
 \mathcal O_{F_j,\Sigma}^{\times}\longrightarrow
   \bigoplus_{w\in\Sigma_0}\ZZ/p^{\min(h_v,j)}
   \longrightarrow B_j\longrightarrow
   \Cl(\mathcal O_{F_j,\Sigma})\{p\}\longrightarrow0 .
 \label{eq:idele113}
\end{equation}
Indeed on removing the full factors at \(\Sigma\) the quotient
becomes the ideal class group of the ring of \(\Sigma\)-integers.
Retaining the indicated local factors adds the displayed finite
groups, modulo the images of principal ideles which are units
outside \(\Sigma\). These come from the global \(\Sigma\)-units.
The quotient is finite, and taking primary parts gives the sequence.

\begin{theorem}[The finite class group formula at every level]
\label{thm:strict113}
With no restriction on the ramification of \(p\) in \(K\), there
is an exact sequence
\begin{equation}
 0\longrightarrow D^{G_K}/p^jD^{G_K}
 \longrightarrow \operatorname{Hom}_{O[\Gamma_j]}
             (B_j\otimes_{\ZZ_p}O,R_j(\theta_j))
 \longrightarrow S_K(\theta)[p^j]\longrightarrow0 .
 \label{eq:strict113}
\end{equation}
In particular, when \(O=\ZZ_p\),
\begin{equation}
 \log|S_K(\theta)[p^j]|
 \ \le\ \log|\Cl(F_j)|+
           (\log p)\sum_{w\in\Sigma_0}\min(h_v,j).
 \label{eq:size113}
\end{equation}
\end{theorem}
\begin{proof}
Multiplication by \(p^j\) on \(D\) is onto, so every class killed
by \(p^j\) in \(H^1(K,D)\) has a lift in \(H^1(K,A_j)\).
Under Lemma~\ref{lem:cyclic113} the lifts with image in the strict
group correspond to the equivariant homomorphisms with the local
conditions of Lemma~\ref{lem:local113}. Global Artin reciprocity,
with its compatibility with local reciprocity, identifies these
with the homomorphisms from the idele class quotient just defined.
The targets are \(p\)-primary, so passage to \(B_j\) and extension
of scalars give the middle term of \eqref{eq:strict113}.
The archimedean local factors have no nontrivial finite quotient
of odd order. This application uses local and global class field
theory.~\cite{II-II-V-cft101}

The kernel on global cohomology is the boundary of
\(D^{G_K}/p^jD^{G_K}\). Every such boundary meets all the local
conditions since its divisible-coefficient class is zero. This
proves exactness. For \(O=\ZZ_p\), the number of homomorphisms
from a finite abelian \(p\)-group \(B\) to \(\ZZ/p^j\) is at
most \(|B|\), by cyclic decomposition. Dropping equivariance,
applying \eqref{eq:idele113} and using the surjection of the
ordinary class group onto the \(\Sigma\)-class group proves
\eqref{eq:size113}.
\end{proof}

\begin{corollary}[Recovering successive orders]\label{cor:levels113}
Let \(\gamma\) generate \(\Gamma_j\), and form the finite module
\[
 C_j=(B_j\otimes_{\ZZ_p}O)/
       \bigl(p^j,\gamma-\theta_j(\gamma)\bigr)(B_j\otimes_{\ZZ_p}O).
\]
Any lift to \(O\) of the scalar \(\theta_j(\gamma)\) gives the same
quotient. Put \(a_j=\len_O C_j-\len_O(D^{G_K}/p^jD^{G_K})\)
and \(a_0=0\). Then \(a_j=\len_O S_K(\theta)[p^j]\). Equality
\(a_{j+1}=a_j\) implies \(S_K(\theta)=S_K(\theta)[p^j]\).
The subtracted term is zero for a nontrivial global residual
character. For a trivial global residual character it equals
\(\min(j,h_K)\), with \(\theta(G_K)=1+p^{h_K}\ZZ_p\).
\end{corollary}
\begin{proof}
Equivariance on a cyclic group imposes precisely the displayed
relation on the source. The length of
\(\operatorname{Hom}_O(C_j,R_j)\) is \(\len_O C_j\), by cyclic
decomposition. Theorem~\ref{thm:strict113} proves the length
identity. If there were an element of order greater than \(p^j\),
a multiple of it would have order \(p^{j+1}\), contradicting
equality of the successive finite orders. Finally a nontrivial
residual value has unit difference from one. In the trivial
residual case the global cyclotomic image is open, as \(K\) is
finite over \(\QQ\), and the same invariant calculation as over
\(K_v\) gives \(D^{G_K}=p^{-h_K}O/O\).
\end{proof}

For \(K=\QQ\) and \(\eta=\alpha^{-1}\) as in
Section~\ref{sec:ray101}, this gives a finite class computation
also above \(1+v_p(r)\). For the period field \(K=K_\psi C_{l'}(\sqrt{-q},\sqrt{-d_0})\)
and \(\eta=1\) it applies
directly to the entire strict group on the period field, including
the primes dividing its degree or conductor.


\subsection{Finite odd cocycles and the strict local condition}
\label{sec:oddcocycles675}

Let \(E\) be a CM number field in a fixed algebraic closure of
\(\QQ\), stable under the chosen complex conjugation \(c\).
Fix a prime \(\ell\) and an integer \(r\ge1\). Write
\(G_E=\operatorname{Gal}(\overline{\QQ}/E)\),
\(A=(\QQ_\ell/\ZZ_\ell)(-r)\),
\(A_\nu=A[\ell^\nu]\), and let \(\kappa\) be the
\(\ell\)-adic cyclotomic character. Thus \(g\) acts by
\(\kappa(g)^{-r}\). For \(\eta\in\{+1,-1\}\), let
\(Z^\eta_\nu(E,r)\) consist of the continuous cocycles
\(z:G_E\to A_\nu\) satisfying the following conditions.
\begin{enumerate}
\item They vanish on every inertia group away from \(\ell\).
\item At every place above \(\ell\) they vanish on
\(D_v\cap\ker\kappa\), for every choice of the decomposition group
\(D_v\) in \(G_E\).
\item They satisfy \(z(cgc)=\eta z(g)\).
\end{enumerate}
Here a cocycle is the function itself. Its image in
\(H^1(G_E,A)\) is obtained by including \(A_\nu\) in \(A\) and
then passing to the class modulo coboundaries. The sign \(-1\)
is the positive-weight module called \texttt{shaOdd} in the
second source archive. Write \(\mathcal S_{\ell,r}(E)\) for the
strict group of Section~\ref{sec:critical13}, and a superscript
in parentheses for the eigensubgroup under the natural action
of \(c\) on cohomology.

\begin{theorem}[Comparison at every coefficient level for odd primes]
\leavevmode\label{thm:oddstrict675}
For \(\ell>2\), including primes dividing \([E:\QQ]\), the class
map gives
\[
 \begin{aligned}
 Z^-_\nu(E,r)&\ \simeq\
     \mathcal S_{\ell,r}(E)^{((-1)^{r+1})}[\ell^\nu],\\
 Z^+_\nu(E,r)/\delta(A)[\ell^\nu]&\ \simeq\
     \mathcal S_{\ell,r}(E)^{((-1)^r)}[\ell^\nu],
 \end{aligned}
\]
where \((\delta a)(g)=(\kappa(g)^{-r}-1)a\).
The group divided out in the second line is cyclic of order
\(\ell^\nu\). The first isomorphism in particular describes the
whole finite odd cocycle module, including its elementary divisors.
\end{theorem}
\begin{proof}
Set \(\Gamma_v=\kappa(D_v)\). This is an open subgroup of
\(\ZZ_\ell^\times\), since \(E_v/\QQ_\ell\) is finite and the
cyclotomic image over \(\QQ_\ell\) is \(\ZZ_\ell^\times\).
Put \(\Gamma'_v=\Gamma_v\cap(1+\ell\ZZ_\ell)\).
It is procyclic and torsion-free. A topological generator
\(\gamma\) acts on \(A\) by \(\gamma^{-r}\ne1\).
Multiplication by \(\gamma^{-r}-1\) on \(A\) is surjective.
Thus any continuous cocycle on \(\Gamma'_v\) agrees with a
coboundary at \(\gamma\), hence on its dense cyclic subgroup
and everywhere by continuity. The quotient \(\Gamma_v/\Gamma'_v\)
has order prime to \(\ell\). Averaging annihilates its positive
cohomology with \(\ell\)-primary coefficients. Inflation and
restriction now give \(H^1(\Gamma_v,A)=0\).

By the second local condition, the restriction of \(z\) to
\(D_v\) factors through \(\Gamma_v\). Its class with coefficients
in \(A\) is therefore zero. The first condition lets \(z\)
descend to the Galois group unramified away from \(\ell\).
Indeed the inertia groups there act trivially on its target, and
\(z\) vanishes on their closed normal span. Its class is strict.
The natural action of \(c\) on a class represented by \(z\)
is represented by \(g\mapsto(-1)^r z(cgc)\), since \(c\)
also acts on the coefficients. This gives the signs in the
two formulas.

Conversely take a strict class killed by \(\ell^\nu\) in the
eigensubgroup of sign \(\eta(-1)^r\). Represent it by a cocycle
\(y\) on the unramified-outside-\(\ell\) group. Write
\(\ell^\nu y=\delta a\) and choose \(b\in A\) with
\(\ell^\nu b=a\). Subtracting \(\delta b\) makes \(y\)
\(A_\nu\)-valued. At each \(v\mid\ell\) strictness means
\(y|_{D_v}=\delta a_v\) in \(A\), so \(y\) vanishes on
\(D_v\cap\ker\kappa\). This holds for every conjugate
decomposition group. The cocycle
\(g\mapsto(y(g)+\eta y(cgc))/2\) has all three stated properties
and represents the same class, by its assumed eigenvalue.

Finally every coboundary with values in \(A\) is unchanged by
\(g\mapsto cgc\), since the action on coefficients is cyclotomic.
For the negative sign this forces a coboundary to vanish. For the
positive sign the kernel is precisely \(\delta(A)[\ell^\nu]\).
The kernel of \(\delta:A\to Z^1(G_E,A)\) is finite, because
the global cyclotomic image is open and \(r\ge1\). A quotient
of \(\QQ_\ell/\ZZ_\ell\) by a finite subgroup is isomorphic to
\(\QQ_\ell/\ZZ_\ell\), as multiplication by the order of that
subgroup shows. This gives the stated kernel order.
\end{proof}

\begin{proposition}[The defect at two]
\leavevmode\label{prop:dyadicodd675}
Let \(s_2(E)\) be the number of dyadic places of \(E\).
For \(\ell=2\) and \(\nu\ge1\), the kernel of the class map on
\(Z^-_\nu(E,r)\) has order two. Its image modulo its intersection
with \(\mathcal S_{2,r}(E)\) has order at most \(2^{s_2(E)}\).
The image contains
\(2\bigl(\mathcal S_{2,r}(E)^{((-1)^{r+1})}[2^\nu]\bigr)\).
In particular,
\[
 \len_{\ZZ_2} Z^-_\nu(E,r)
      \le 1+s_2(E)+\len_{\ZZ_2}\mathcal S_{2,r}(E)[2^\nu].
\]
\end{proposition}
\begin{proof}
A dyadic coboundary is invariant under \(g\mapsto cgc\), so it
satisfies the odd relation exactly when it is killed by two.
The preceding computation of \(\delta(A)\) then gives one
subgroup of order two in the kernel. Its elements satisfy both
local conditions and take values in \(A[2]\), giving equality.
Now take \(\Gamma'_v=\Gamma_v\cap(1+4\ZZ_2)\). The generator
argument in the theorem still gives \(H^1(\Gamma'_v,A)=0\).
The quotient \(\Gamma_v/\Gamma'_v\) has order at most two.
It acts on the finite cyclic group \(A^{\Gamma'_v}\).
Its first cohomology is a subquotient of that cyclic group by
evaluation at a generator, and is killed by the group order.
Consequently \(H^1(\Gamma_v,A)\) has order at most two by
inflation and restriction. Localizing the image at all dyadic
places bounds its quotient by the strict part as asserted.
The construction with \(y-\delta b\) in the theorem also gives
a finite-valued representative \(y_0\) of any strict class in
the displayed eigensubgroup. The cocycle
\(g\mapsto y_0(g)-y_0(cgc)\) satisfies the odd relation
and represents twice the given class. This proves the inclusion
and the length bound.
\end{proof}

\begin{corollary}[The total of the finite odd cocycle lengths on the period fields]
\leavevmode\label{cor:finiteoddtotal675}
On the polynomial-conductor intervals of
Theorem~\ref{thm:conductor232} take \(E=E_x\). One may choose
\(\nu_{\ell,r}\ge0\) separately at each prime, weight and node.
One may also use any finite extension of \(\QQ_\ell\) with
integer ring \(\mathcal O_{\ell,r}\) and ramification index
\(e_{\ell,r}\) for the coefficient module
\((\mathcal O_{\ell,r}/\ell^{\nu_{\ell,r}})(-r)\).
With that coefficient ring understood on the left,
\[
 \sum_{\ell}\log\ell\sum_{1\le r<t(x)}
     \frac{\len_{\mathcal O_{\ell,r}}Z^-_{\nu_{\ell,r}}(E_x,r)}
          {e_{\ell,r}}
             =o_{\psi,\eps}\bigl(\exp(x^\eps/2)\bigr)
\]
uniformly on the selected intervals.
\end{corollary}
\begin{proof}
The finite coefficient module is the original one tensored with
\(\mathcal O_{\ell,r}\) over \(\ZZ_\ell\). In a finite free
integral basis all the cocycle, parity and local conditions are
coordinatewise. Thus the new solution module is the original
one tensored with this ring, compatibly with its scalar action.
A cyclic factor of length \(a\) becomes
\(\mathcal O_{\ell,r}/\ell^a\) and has length \(e_{\ell,r}a\).
It suffices to use \(\ZZ_\ell\).

For odd primes the first isomorphism of the theorem bounds the
sum by the critical-sign estimate of
Theorem~\ref{thm:critical13}. Indeed its proof gives a
polynomial bound in the field degree, the maximum weight and
the logarithm of the conductor. All are polynomial on these
fields. At two the overhead in the proposition is
\(O(t(x)(1+[E_x:\QQ]))\) after summing over \(r\).
The full dyadic lengths on these same fields are bounded by
\(O_{\psi,\eps}(t(x)\log(2x)+\exp(C_{\psi,\eps}x^{3\eps/8}))\)
by the interval selection in Theorem~\ref{thm:conductor232}
and the proof of Theorem~\ref{thm:two17}. These bounds are
\(o(\exp(x^\eps/2))\). The zero coefficient level contributes
zero. This proves the assertion uniformly in all the choices.
\end{proof}

For odd primes the unestimated sign in \eqref{eq:length12}
is the second line of Theorem~\ref{thm:oddstrict675}.
A bound for the odd cocycles alone therefore supplies the first
line. The even cocycles modulo the stated coboundaries still
enter the original length estimate.


\subsection{A cyclotomic annihilator at all coefficient primes}
\label{sec:cocycleannihilator676}

Fix an integer \(N\ge3\), put \(G_N=(\ZZ/N\ZZ)^\times\), and
let \(E\) be any CM subfield of \(\QQ(\mu_N)\). For \(r\ge1\)
put \(n=r+1\). Write \(B_n(X)\) for the Bernoulli polynomial
with generating series \(te^{Xt}/(e^t-1)\). In the integral
quotient ring
\[
  R_n=\ZZ[G_N]/([-1]-(-1)^n[1])
\]
form the element and its integer determinant
\[
 T_n=\sum_{j\in G_N}(n+1)!\,N^n B_n(j/N)[j],\qquad
 d_n=\det(T_n:R_n\longrightarrow R_n).
\]
Every \(j\) occurring as an integer denotes its representative
in \([1,N-1]\). Choose \(b_1=1+2N,\ b_2=1+2Nb_1\). They are
coprime to one another, each is congruent to one modulo \(N\),
and both exceed one.

\begin{lemma}[Moments of the integral Stickelberger coefficients]
\leavevmode\label{lem:stickmoments676}
Fix a prime \(\ell\) and one of \(b_1,b_2\) prime to \(\ell\),
denoted by \(b_0\). For \(M=N\ell^L\) let
\[
 S_j(L)=
 \sum_{\substack{1\le a<M,\ (a,M)=1\\a\equiv j\bmod N}}
             a^{\,n-1}\lfloor b_0a/M\rfloor .
\]
In \(\QQ_\ell[G_N]\) the following limit holds, also at
\(\ell=2\). In the formula for \(T_n\) on the right use its
displayed representative in the full group ring.
\[
 \lim_{L\to\infty} n b_0^{\,n-1}N(n+1)!
           \sum_{j\in G_N}S_j(L)[j]
     =(b_0^n-1)
     \begin{cases}
       T_n,&\ell\mid N,\\
       (1-\ell^{\,n-1}[\ell])T_n,&\ell\nmid N .
     \end{cases}
\]
\end{lemma}
\begin{proof}
Multiplication by \(b_0\) modulo \(M\) permutes the units in
each class \(j\bmod N\). Its representative is
\(a'=b_0a-M\lfloor b_0a/M\rfloor\). Expanding \((a')^n\),
summing, and using \(\sum (a')^n=\sum a^n\), gives
\[
 n b_0^{n-1} S_j(L)
   -\frac{b_0^n-1}{M}\sum_{\substack{a\equiv j\bmod N\\
                       1\le a<M,\ (a,M)=1}}a^n
                   \ \in\ M\ZZ .
\]
For \(v=\ell^L\), first include every \(a=j+Nh\), \(0\le h<v\).
The normalized \(n\)-th power sum is
\[
 \frac1{Nv}\sum_{h=0}^{v-1}(j+Nh)^n
   =\frac{N^{n-1}}{n+1}\,
       \frac{B_{n+1}(j/N+v)-B_{n+1}(j/N)}{v}.
\]
This follows by summing the polynomial finite-difference formula
\(B_{n+1}(y+1)-B_{n+1}(y)=(n+1)y^n\).
As \(v\) tends to zero in \(\QQ_\ell\), the quotient tends to
\(B'_{n+1}(j/N)=(n+1)B_n(j/N)\).
For \(\ell\mid N\) the condition \((a,M)=1\) is automatic.
Otherwise the excluded terms have \(a=\ell a_0\) with
\(a_0\equiv\ell^{-1}j\bmod N\), \(1\le a_0<M/\ell\).
After division by \(M\) their limit is \(\ell^{n-1}\) times
the corresponding limit at \(\ell^{-1}j\). Combining these
computations gives the group-ring formula.
\end{proof}

\begin{theorem}[One integer annihilates the odd cocycles at every prime]
\leavevmode\label{thm:uniformannihilator676}
The positive integer
\[
           J_{N,r}=2(b_1^n-1)(b_2^n-1)|d_n|
\]
annihilates \(Z^-_\nu(E,r)\) for every prime \(\ell\), every
\(\nu\ge0\) and every such \(E\). It satisfies
\(\log J_{N,r}\ll N n\log(2Nn)\) with an absolute constant.
\end{theorem}
\begin{proof}
The coefficients of \(T_n\) are integers. Indeed
\(B_n(X)=\sum_{i=0}^n\binom ni B_i X^{n-i}\) and
\((i+1)!B_i\in\ZZ\). The latter follows inductively from
\(\sum_{j=0}^{i}\binom{i+1}{j}B_j=0\) for \(i\ge1\), since
\(i!\binom{i+1}{j}\) is divisible by \((j+1)!\) for \(j<i\).
The ring \(R_n\) is free over \(\ZZ\), with basis represented
by one member of each pair \(\{j,-j\}\). Over \(\mathbb C\)
its characters are precisely those with \(\chi(-1)=(-1)^n\).
The Bernoulli formula for \(L(1-n,\chi)\) identifies the
eigenvalue of \(T_n\) at each such character as
\(-n(n+1)!N L(1-n,\chi)\), using the character as induced to
modulus \(N\). It is nonzero by the functional equation and
the Euler product at \(n>1\). Each removed Euler factor is
\(1-\chi^*(p)p^{n-1}\ne0\) when \(\chi^*(p)\ne0\).
This proves \(d_n\ne0\). On \(0\le y\le1\) the generating
series on \(|t|=1\) gives \(|B_n(y)|\le Cn!\) by Cauchy's
inequality with an absolute constant. Each eigenvalue of
\(T_n\) has modulus at most \(CN(n+1)!N^n n!\).
There are at most \(N\) eigenvalues and \(b_2\ll N^2\).
Their product gives the stated logarithmic bound.

Fix a cocycle \(z\in Z^-_\nu(E,r)\), taking \(\nu>0\). Put
\(K_M=\QQ(\mu_M)\), \(M=N\ell^L\). For \(L\ge\nu\) the
restriction of \(z\) to \(G_{K_M}\) is a homomorphism, since
that group acts trivially on \(A_\nu\). For all sufficiently
large \(L\) it is unramified everywhere. Away from \(\ell\)
this follows from the first local condition. At each \(v\mid\ell\)
the second condition makes \(z|_{D_v}\) a continuous function
of \(\kappa\) vanishing at one. Its target is finite, so this
function vanishes for \(\kappa\equiv1\bmod\ell^L\) when \(L\)
is sufficiently large. One \(L\) works for representatives of
the finitely many \(v\)'s. For their conjugates the cocycle
identity gives \(z(shs^{-1})=\kappa(s)^{-r}z(h)\) when \(h\)
acts trivially on \(A_\nu\), which proves the same vanishing.
Class field theory now gives a homomorphism
\(\phi_L:\Cl(K_M)\to A_\nu\) representing the restriction.

Write \(\sigma_a\) for the automorphism of \(K_M\) raising
roots of unity to their \(a\)-th powers. On these homomorphisms
put \(P_a\phi(x)=\phi(\sigma_a^{-1}x)\), using the natural
action on ideal classes. Set \(U_a=a^{-r}P_a\), with the
factor read modulo \(\ell^\nu\). If \(a\equiv1\bmod N\),
a lift \(s\) of \(\sigma_a\) fixes \(E\). On \(G_{K_M}\) the
identity \(z(s^{-1}hs)=\kappa(s)^r z(h)\) then gives
\(U_a\phi_L=\phi_L\) by equivariance of the Artin map. The
operators \(U_a\) commute. On the module generated by their
translates of \(\phi_L\) they therefore give an action of
\(G_N\), using surjectivity of reduction from the units
modulo \(M\). Moreover \(P_{-1}\phi_L=-\phi_L\), by the
odd cocycle relation. Thus \(U_{-1}\) acts by
\((-1)^{r+1}=(-1)^n\) on this module and the action factors
through the integral ring \(R_n\), including at two.

Choose \(b_0\) among the two specified integers prime to
\(\ell\). The classical Stickelberger theorem for the full
cyclotomic field \(K_M\) says that
\((b_0-\sigma_{b_0})\sum_a(a/M)\sigma_a^{-1}\) annihilates
\(\Cl(K_M)\), where the sum is over reduced representatives.
The hypotheses are \((b_0,M)=1\) and \(K_M=\QQ(\mu_M)\), both
of which hold here.~\cite[Ch.~6]{II-II-V-stick}
The coefficient of \(\sigma_a^{-1}\) in this integral element
is exactly \(\lfloor b_0a/M\rfloor\), by the same permutation
of representatives used above. Evaluating \(\phi_L\) on its
action on arbitrary ideal classes gives
\(\sum_j S_j(L)U_j\phi_L=0\).

By Lemma~\ref{lem:stickmoments676}, for all sufficiently large
\(L\) the coefficients of the scaled last relation agree modulo
\(\ell^\nu\) with those of the limit. Only this finite
congruence is used, so the class groups may vary with \(L\).
If \(\ell\nmid N\), the extra factor \(1-\ell^{n-1}U_\ell\)
is invertible on a module killed by \(\ell^\nu\), by a finite
geometric sum. In either case
\((b_0^n-1)T_n\phi_L=0\). In the regular representation of
\(R_n\), the image ideal \(T_nR_n\) contains \(d_n[1]\), by
the adjugate identity applied to the basis vector of \([1]\).
It follows that \((b_0^n-1)d_n\phi_L=0\).

The cocycle \((b_0^n-1)d_n z\) consequently factors through
\(\operatorname{Gal}(K_M/E)\). Conjugating an element of this
quotient by \(c\) fixes it, since \(K_M/\QQ\) is abelian. The
descended function is also odd under that conjugation. Its
values are therefore killed by two. This proves the asserted
annihilation for the one \(J_{N,r}\) independent of \(\ell,\nu,E\).
\end{proof}
\clearpage


\subsection{One integer for the full length of the odd cocycle module}
\label{sec:oddlengthinteger677}

We now pass from the exponent in
Theorem~\ref{thm:uniformannihilator676} to orders. Write
\(F_N=\QQ(\mu_N)\), \(g=\varphi(N)\) and choose primes
\(N<P_1<2N<P_2<4N\), which exist by Bertrand's postulate for
\(N\ge3\). For \(n\ge2\) put
\(W_n=(P_1^{gn}-1)(P_2^{gn}-1)\).
Keep \(d_n\) and \(J_{N,n-1}\) of the preceding subsection.
For \(t\ge2\) define the positive integer
\[
 \mathcal J_{N,t}=
   \prod_{n=2}^{t}J_{N,n-1}^{20N^3}\,W_n^{g/2}\,|d_n| .
\]
For \(t=0,1\) use the empty product. Throughout the subsection
\(E\) is any CM subfield of \(F_N\).

\begin{theorem}[Uniform orders through all weights and coefficient levels]
\leavevmode\label{thm:oddinteger677}
For each prime \(\ell\) and each choice of nonnegative integers
\(\nu_{\ell,r}\), one has
\[
 \sum_{1\le r<t}\len_{\ZZ_\ell}Z^-_{\nu_{\ell,r}}(E,r)
                    \le v_\ell(\mathcal J_{N,t}),\qquad
 \log\mathcal J_{N,t}\ll N^4t^2\log(2Nt)\quad(t\ge2).
\]
The integer works for all the subfields \(E\) at once. To include
\(r=0\) with the convention of the second archive, multiply it
by \(h_E\), the class number of \(E\). With a finite coefficient
ring of integers \(\mathcal O\) above \(\ell\), replace a length on
the left by its \(\mathcal O\)-length divided by the ramification
index. The coefficient ring may vary with the weight.
\end{theorem}
\begin{proof}
First suppose \(\ell\mid 2g\). In particular \(\ell\le N\).
The \(\ell\)-torsion in the finite-coefficient cocycle module
embeds into the space of continuous cocycles of the character
\(\bar\kappa^{-r}\) over \(\mathbb F_\ell\), unramified outside
\(\ell\). We bound the dimension of that space. Over
\(K=E(\mu_\ell)\) the action is trivial, so restriction gives
homomorphisms unramified outside the places above \(\ell\).
The kernel consists of cocycles on \(\operatorname{Gal}(K/E)\).
This finite group has order dividing \(\ell-1\), so those
cocycles are coboundaries and span a space of dimension at most
one. Kummer theory over \(K\), which contains \(\mu_\ell\),
bounds the dimension of the homomorphisms by the dimension of
the group of units of \(\mathcal O_K[1/\ell]\) modulo
\(\ell\)-th powers plus \(\log_\ell h_K\). Indeed a Kummer
element has valuation divisible by \(\ell\) away from \(\ell\).
Sending it to the resulting ideal class gives an \(\ell\)-torsion
class in the open ring, whose class number divides \(h_K\).
The kernel consists of the indicated units modulo powers.

Put \(d=[K:\QQ]\le N^2\). The unit contribution is at most
\(2d\) by the unit theorem including its roots of unity, since
there are at most \(d\) places above \(\ell\). The conductor
of \(K\) divides \(N\ell\), so its absolute discriminant is
at most \((N\ell)^d\) by the conductor--discriminant formula.
Minkowski's bound gives each ideal class an integral
representative of norm at most \(B=2^d(N\ell)^{d/2}\).
The number of ideals up to that norm is at most
\(B^2\zeta_K(2)\le B^2\zeta(2)^d\). Hence
\(\log_2 h_K\le d(3+2\log_2 N)\).
Together these estimates give dimension at most
\(1+(5+2\log_2 N)N^2\le20N^3\).
A group of exponent dividing \(\ell^\nu\) with finite
\(\ell\)-torsion is finite, by induction on \(\nu\). Each
successive layer has dimension at most the dimension of that
torsion, by multiplication by a power of \(\ell\) on the
filtration by kernels. Applying
Theorem~\ref{thm:uniformannihilator676} bounds the full length
at \(r=n-1\) by \(20N^3v_\ell(J_{N,n-1})\).

Now let \(\ell\nmid2g\). By Theorem~\ref{thm:oddstrict675} it
suffices to bound the strict eigensubgroup with eigenvalue
\((-1)^n\). Restriction embeds this group into the corresponding
group over \(F_N\). It preserves all strict local conditions
and commutes with \(c\), and corestriction after restriction
is multiplication by \([F_N:E]\), a unit at \(\ell\).

Apply the Fitting-ideal argument of Lemma~\ref{lem:fit13} to
\(F_N/\QQ\). For clarity, the set of removed primes contains
infinity and all finite primes dividing \(N\), the twist has
\(n\ge2\), and \(\ell\) is odd. The cyclotomic \(\mu\)-invariant
in the refined Coates--Sinnott theorem vanishes by
Ferrero--Washington for abelian fields. The dual of the indicated
strict group embeds in the same sign of degree-two cohomology
over \(\mathcal O_{F_N}[1/(N\ell)]\). In the sign ring
\(R_n\otimes\ZZ_\ell\) its ambient module has zeroth Fitting
ideal containing \(W_n\Theta_n\), where the character values of
\(\Theta_n\) are \(L(1-n,\chi^{-1})\) with the Euler factors
over \(N\) removed. In fact \(W_n\) kills the degree-one torsion
\(H^0(F_N,\QQ_\ell/\ZZ_\ell(n))\) by the action of Frobenius
at \(P_i\ne\ell\), whose residue degree divides \(g\).
A further omitted Euler factor at \(\ell\nmid N\) is a unit
in the \(\ell\)-adic group ring because \(n-1>0\).
These are precisely the hypotheses and the Fitting membership
in the cited lemma.~\cite[Thm.~6.11]{II-II-II-gp13}\cite{II-II-II-fw13}

The ring \(R_n\otimes\ZZ_\ell\) is a product of unramified
discrete valuation rings since \(\ell\nmid2g\). On each
factor the valuation of an element of the zeroth Fitting ideal
bounds the length of the module. Multiplying \(W_n\Theta_n\)
by the integer \(n(n+1)!N\), taking the component valuations
with residue-degree weights and adding, gives
\[
        \len_{\ZZ_\ell}\mathcal S_{\ell,n-1}(F_N)^{((-1)^n)}
                      \le v_\ell(W_n^{g/2}|d_n|).
\]
Indeed the absolute determinant of \(n(n+1)!N\Theta_n\) on the
sign ring is \(|d_n|\), by its eigenvalues as computed in
Theorem~\ref{thm:uniformannihilator676}. Inverting \(\chi\)
permutes the characters of that sign. Duality preserves
\(\ZZ_\ell\)-length. Multiplying the integer bounds for all \(n\)
proves the valuation inequality in both prime ranges.

The bound for \(\log\mathcal J_{N,t}\) follows by summing
\(\log J_{N,n-1}\ll Nn\log(2Nn)\) and
\(\log W_n\le2gn\log(4N)\). For comparison with an explicit
size used in the second archive, one also has
\(\mathcal J_{N,t}\le 2^{4(N+t+3)^{27}}\).
To see this, put \(y=N+t+3\), with \(t\ge2\). In the Cauchy
estimate from the preceding theorem one may take \(12n!\),
since \(|e^z-1|\ge 3-e>1/4\) on \(|z|=1\) and
\(|e^{uz}|<3\) for \(0\le u\le1\).
Thus each eigenvalue has modulus at most
\(12N(n+1)!N^n n!\). Using \(n,N\le y\) gives
\(\log_2|d_n|\le5y^3\).
Since \(b_1\le3N,\ b_2\le5N^2\), this gives
\(\log_2 J_{N,n-1}\le6y^3\), and
\(\log_2 W_n\le2y^3\). It follows that
\(\log_2\mathcal J_{N,t}\le130y^7\le4y^{27}\).
The empty product satisfies the same absolute bound.

At weight zero the model is the group of everywhere
unramified homomorphisms into \(\ZZ/\ell^\nu\), with no sign
condition. It is \(\operatorname{Hom}(\Cl(E),\ZZ/\ell^\nu)\)
by class field theory for the totally imaginary field \(E\).
Its length is at most \(v_\ell(h_E)\). The finite free
base-change calculation of Corollary~\ref{cor:finiteoddtotal675}
proves the assertion for every coefficient ring, including at
weight zero.
\end{proof}

\begin{corollary}[The same integer on the period interval]
\leavevmode\label{cor:oddwindow677}
On a polynomial-conductor interval \([V,V^{3/2}]\), let \(N_*\)
be its upper-endpoint cyclotomic modulus and let
\(t_*=\lfloor V^{1203/2}\rfloor\). For every node \(x\) on
this interval, every finite coefficient level, and all primes
simultaneously, \(\mathcal J_{N_*,t_*}\) bounds the positive-weight
odd cocycle lengths as in the theorem. Its logarithm is bounded
by a power of \(V\). On the family in the shorter archive the
same theorem supplies the entire odd-cocycle input, including
weight zero, with the bound \(2^{4(N+t+3)^{27}} h_E\).
For that family's fixed \(q_0,d\), writing
\(J'=h_E\mathcal J_{N,t}\), the particular contribution
\(t\,v_\ell(J')\) to the denominator estimate for the lifts
satisfies
\[
        \sum_\ell t\,v_\ell(J')\log\ell
          =t\log J'\ll_{q_0,d}x^{59}\log^5(2x).
\]
\end{corollary}
\begin{proof}
In Theorem~\ref{thm:conductor232} every \(E_x\) belongs to
\(\QQ(\mu_{N_*})\), since only the dyadic conductor changes
and each layer is contained in the top one. The weights
\(1\le r<t(x)\) are below \(t_*\), and both \(N_*,t_*\) are
polynomial in \(V\). Apply the theorem to this single \(N_*,t_*\).

For the second assertion the shorter construction uses
\(E=\QQ(\sqrt{-d},\mu_{q_1})\), where \(d\equiv7\bmod 8\)
is prime. The quadratic field has discriminant \(-d\) and
embeds in \(\QQ(\mu_d)\), so \(E\subseteq\QQ(\mu_N)\) for
\(N=q_1d\). The field is totally imaginary and abelian over
\(\QQ\). Complex conjugation fixes a totally real subfield of
index two, so \(E\) is CM. The local conditions and coefficient sign are
those specified in Section~\ref{sec:oddcocycles675}, also at
the primes dividing \(N\). Apply the theorem and multiply by
\(h_E\) for the zero weight. The result works without adjoining
any character values to the coefficient ring.

To compute its size, the number \(m\) of lifts lies between
\(\sqrt x\) and \(2\sqrt x\), and the tower exponent \(s\)
satisfies \(2^s=m^{22}\). With \(l_0\) the least power of three
at least \(27s\), the modulus satisfies \(q_1=3q_0 l_0 2^s\).
Thus \(l_0\le81s\) and \(N\ll_{q_0,d}x^{11}\log(2x)\).
The weight cutoff is five times a rank of size at most \(3x^5\).
The class number bound in the proof of the theorem gives
\(\log h_E\ll N\log(2N)\) by applying the same Minkowski estimate
to \(E\subseteq\QQ(\mu_N)\). Hence \(\log J'\ll_{q_0,d}
x^{54}\log^5(2x)\). Multiplication by \(t\le15x^5\) and prime
factorization of \(J'\) prove the asserted bound for this term.
\end{proof}

For the shorter argument this establishes the bound on the odd
cocycles used in its first-order congruence calculation. The
bounds for the denominators of the projected lifts and for their
periods involve further terms. For the original comparison
\eqref{eq:length12}, positive-weight cocycles of the other sign
still occur as described in Theorem~\ref{thm:oddstrict675}.
\clearpage

\section{Selmer groups and marked local conditions}\label{nav:selmer-tests}

\subsection{Dual Selmer classes and relations on tame inertia}
\label{sec:evenrelations712}

Let \(E\) be a CM field, let \(F=E^+\) be its maximal real
subfield, and put \(f=[F:\QQ]\). Suppose that \(E/F\) is
unramified at every finite place. We will verify this hypothesis on
both families of period fields at the end of the subsection.
Write \(G_L=\Gal(\overline{\QQ}/L)\) for each field in a fixed
algebraic closure. Choose a complex conjugation \(c\in G_F\).
Let \(\varepsilon\) be the quadratic character of \(E/F\) and
\(\kappa\) the \(\ell\)-adic cyclotomic character, for a prime
\(\ell\). The character \(\kappa_+=\kappa\varepsilon\) equals
\(\kappa\) on \(G_E\) and equals \(1\) at every real
conjugation of \(F\).

Fix an exponent \(e\ge1\) and a coefficient precision \(\nu\ge1\).
Put \(R_n=\ZZ/\ell^n\ZZ\) for \(1\le n\le\nu\). Use the two
rank-one modules
\[
       A_n=R_n(\kappa_+^{-e}),\qquad
       B_n=\operatorname{Hom}(A_n,\mu_{\ell^n})
           \simeq R_n(\kappa\kappa_+^e).
\]
In this notation \(R_n(\chi)\) is the additive group of \(R_n\)
on which \(g\) acts by multiplication by \(\chi(g)\). Choose
compatible primitive roots \(\zeta_{\ell^n}\) for the displayed
identifications. The duality pairing in the two coordinates sends
\((a,b)\) to \(\zeta_{\ell^n}^{ab}\). At level zero use the zero
modules. Let \(X\) be a finite set of primes of \(F\) away from
\(\ell\) splitting completely in \(E(\mu_{\ell^\nu})/F\). We
keep this same \(X\) through all the levels \(n\le\nu\).

At a real place or a prime in \(X\) prescribe
\(L_{n,v}=0\) in \(H^1(F_v,A_n)\). At each other prime
\(v\nmid\ell\) prescribe unramified cohomology, the kernel of
restriction to inertia. For \(v\mid\ell\) write
\(\Gamma_v=\kappa_+(G_{F_v})\), using the full \(\ell\)-adic
image at every level. Except in the following dyadic case,
take for \(L_{n,v}\) the inflation of \(H^1(\Gamma_v,A_n)\).
If \(\ell=2\) and \(-1\in\Gamma_v\), prescribe
\[
 L_{n,v}=
 \begin{cases}
  \operatorname{inf}H^1(\Gamma_v/\{1,-1\},A_n),& e\ \text{even},\\
  0,& e\ \text{odd}.
 \end{cases}
\]
The quotient acts on \(A_n\) in the first line. These subgroups
define \(S_n\subset H^1(F,A_n)\) by inverse image under
localization. Define \(D_n\subset H^1(F,B_n)\) in the same
way using \(L_{n,v}^{\perp}\) for the local Tate pairing at
each place. Write \(S_0=D_0=0\).

\begin{lemma}[Change of coefficient level]\leavevmode
\label{lem:duallocal712}
Reduction from \(n\) to \(n-1\) maps \(L_{n,v}\) onto
\(L_{n-1,v}\) at every place. Multiplication by \(\ell\)
from \(A_{n-1}\) to \(A_n\) maps \(L_{n-1,v}\) into
\(L_{n,v}\). At \(v\mid\ell\) and \(n\ge1\) one has
\[
 \frac{|L_{n,v}|}{|H^0(F_v,A_n)|}=
 \begin{cases}
    1/2,&\ell=2,\ -1\in\Gamma_v,\ e\ \text{odd},\\
    1,&\text{otherwise}.
 \end{cases}
\]
\end{lemma}
\begin{proof}
Each \(\Gamma_v\) is open in \(\ZZ_\ell^\times\). Indeed, on
the absolute Galois group of the finite local extension \(E_w\)
of \(\QQ_\ell\) the character is \(\kappa\), whose image has
finite index in \(\ZZ_\ell^\times\). For odd \(\ell\) an open
subgroup of \(\ZZ_\ell^\times\) is a product of a finite
group \(\Delta\) of order dividing \(\ell-1\) and a copy of
\(\ZZ_\ell\). This follows from the decomposition into roots
of unity and principal units, whose finite quotients have
coprime orders. Inflation is injective. Averaging over \(\Delta\)
identifies \(H^1(\Gamma_v,A_n)\) with
\(A_n^\Delta/(\gamma-1)A_n^\Delta\), for a generator \(\gamma\)
of the pro-\(\ell\) factor. Here the usual procyclic formula
follows by evaluating a cocycle at the generator. Every value
is allowed because the action factors through a finite quotient
and the coefficient group is finite of \(\ell\)-power order.
Averaging makes reduction on the invariants surjective. Since
the kernel and cokernel of an endomorphism of a finite group
have the same order, the displayed quotient has order
\(|A_n^{\Gamma_v}|\).

For \(\ell=2\), the only possible nontrivial torsion subgroup
of \(\Gamma_v\) is \(\{1,-1\}\), and the quotient by its torsion
is isomorphic to \(\ZZ_2\). Indeed, \(\ZZ_2^\times\) is the
product of \(\{1,-1\}\) and \(1+4\ZZ_2\simeq\ZZ_2\).
With trivial torsion, or with nontrivial torsion and \(e\)
even, the same procyclic calculation gives surjectivity
and the claimed ratio. For \(e\) odd and \(-1\in\Gamma_v\),
the local condition is zero. The fixed subgroup of \(A_n\) is
exactly \(A_n[2]\) since \(-1\) acts by \(-1\) and every
\(2\)-adic unit fixes \(A_n[2]\).

At a finite prime outside \(X\) away from \(\ell\), both
\(\varepsilon\) and \(\kappa\) are unramified. Hence the
unramified group is \(A_n/(\Frob_v-1)A_n\), which again
reduces surjectively and has order \(|H^0(F_v,A_n)|\).
At places with zero condition the reduction assertion is
immediate. In every case the map from \(A_{n-1}\) given by
multiplication by \(\ell\) respects the prescribed inflation
or zero subgroup.
\end{proof}

\begin{theorem}[Independent dual classes through every precision]
\leavevmode\label{thm:dualrelations712}
Put \(U=G_{E(\mu_{\ell^\nu})}\). Every \(y\in D_\nu\)
restricts to a homomorphism \(\psi_y:U\longrightarrow R_\nu\)
which is independent of the cocycle representative. Let \(W\)
be the space of their reductions in
\(\operatorname{Hom}_{\rm cont}(U,\mathbb F_\ell)\). Then
\[
 \dim_{\mathbb F_\ell}W\ \ge\
 \begin{cases}
 f+|X|-1,&\ell\ \text{odd},\\
 f+|X|-3,&\ell=2,\ \mu_4\subset E,\\
 f+|X|-4,&\ell=2\ \text{in general}.
 \end{cases}
\]
Each \(y\) is unramified at every finite prime outside
\(X\) and the primes above \(\ell\).
For \(a\in G_F\) and \(h\in U\) the homomorphism satisfies
\(\psi_y(aha^{-1})=(\kappa\kappa_+^e)(a)\psi_y(h)\) in \(R_\nu\).
\end{theorem}
\begin{proof}
Use logarithms to base \(\ell\) for the four orders
\(s_n=\log_\ell|S_n|\), \(d_n=\log_\ell|D_n|\),
\(a_n=\log_\ell|H^0(F,A_n)|\) and \(b_n=\log_\ell|H^0(F,B_n)|\).
Let \(k\) count the places in the first line of the ratio in
Lemma~\ref{lem:duallocal712}. Put \(k_0=0\) and \(k_n=k\) for
\(n\ge1\). For finite dual Galois modules, the Greenberg--Wiles
formula gives finiteness and the identity
\[
  \frac{|S_n|}{|D_n|}
  =\frac{|H^0(F,A_n)|}{|H^0(F,B_n)|}
   \prod_v\frac{|L_{n,v}|}{|H^0(F_v,A_n)|}.
\]
It applies to continuous cohomology over any number field, with
orthogonal conditions for the local cup pairing and the unramified
condition at almost all primes. Its displayed factors at real
places use ordinary invariants, also at \(\ell=2\).
At each real place \(\kappa_+=1\) on the local group, and
at each \(v\in X\) both characters are trivial modulo
\(\ell^\nu\). These places each contribute \(\ell^{-n}\).
At all remaining places the factors were computed in the lemma.
It follows, including at level zero, that
\[
 d_n-s_n=n(f+|X|)+b_n-a_n+k_n .
\]
The formula and the duality conventions used here are those
of Darmon, Diamond and Taylor.~\cite{II-II-II-gw712}

The exact coefficient sequence from \(A_{n-1}\) to \(A_n\)
to \(A_1\), with maps multiplication by \(\ell\) and reduction,
shows that the kernel of
\(H^1(F,A_{n-1})\longrightarrow H^1(F,A_n)\) has order
\(\ell^{a_1-a_n+a_{n-1}}\). Indeed, the connecting map has
domain \(H^0(F,A_1)\) and its kernel is the reduction of
\(H^0(F,A_n)\), of order \(\ell^{a_n-a_{n-1}}\).
The lemma maps \(S_{n-1}\) into \(S_n\), so
\(s_n-s_{n-1}\ge-a_1+a_n-a_{n-1}\). Also \(b_n\ge b_{n-1}\)
by the injection \(B_{n-1}\hookrightarrow B_n\). Subtracting
the two cardinality identities gives
\[
                  d_n-d_{n-1}\ge f+|X|-a_1\ge f+|X|-1 .
\]

We need to count classes which survive reduction all the way
to \(B_1\). The preimage of \(D_n\) under
\(H^1(F,B_{n-1})\to H^1(F,B_n)\) is exactly \(D_{n-1}\).
Indeed, multiplication by \(\ell\) on the dual coefficients
is adjoint under local cup products to reduction
\(A_n\to A_{n-1}\), and the latter maps every primary local
condition onto the preceding one. Therefore the kernel of
\(D_n\to H^1(F,B_1)\) is the image of \(D_{n-1}\), by the
exact sequence for \(B_{n-1}\), \(B_n\), \(B_1\). The image
in \(H^1(F,B_1)\) has dimension at least \(d_n-d_{n-1}\).
This argument keeps the classes at their full precision before
reducing them.

Finally, \(U\) is normal in \(G_F\) and acts trivially on
\(B_\nu\), so restriction kills coboundaries and gives the
homomorphisms in the statement. Expanding the cocycle at
\(aha^{-1}\) proves their equivariance. For \(n=\nu\), the
kernel of restriction on \(H^1(F,B_1)\) is inflated from
\(H^1(H,B_1)\), where \(H=\Gal(E(\mu_{\ell^\nu})/F)\).
This group \(H\) is abelian, being a subgroup of the product
of the quadratic group and the cyclotomic group. For odd
\(\ell\), complex conjugation in \(H\) acts by \(-1\) on
\(B_1\). If \(u\) is a cocycle on \(H\), evaluating it at
\(cg=gc\) gives \(-2u(g)=(g-1)u(c)\). Division by \(-2\)
exhibits a coboundary, so restriction loses no dimension.

At two the residual module is trivial, and \(H\) embeds by
\((\varepsilon,\kappa\bmod 2^\nu)\) in a product of at most
three cyclic \(2\)-groups. Any subgroup of such a product
is generated by at most three elements, by induction using
projection to one cyclic factor and its kernel. This bounds
the dimension of \(H^1(H,B_1)\) by three. If \(\mu_4\subset E\),
then \(E=F(i)\). Consequently \(\kappa_+(G_F)\subset1+4\ZZ_2\),
since its sign correction is precisely the action on \(i\).
The equally faithful pair
\((\varepsilon,\kappa_+\bmod 2^\nu)\) now takes values in
a product of at most two cyclic groups. This proves the
sharper dyadic assertion.

Outside \(X\) and \(\ell\) the unramified local groups for the
two unramified modules are exact annihilators. Their cup product
vanishes because it inflates from the procyclic unramified
quotient, of \(\ell\)-cohomological dimension one. Their orders
are the respective invariant orders as in the lemma. For a
local field of residual characteristic different from \(\ell\),
the local Euler characteristic formula and local Tate duality
give \(|H^1(F_v,B_n)|=|H^0(F_v,B_n)|\,|H^0(F_v,A_n)|\).
The two unramified groups therefore have the full orders
required by orthogonality. This proves the ramification
assertion.~\cite{II-II-II-duality13}
\end{proof}

\begin{proposition}[Reciprocity on the chosen tame generators]
\leavevmode\label{prop:evenreciprocity712}
Let \(Q\) be another finite set of primes of \(F\), disjoint
from \(X\) and from \(\ell\), splitting completely in
\(E(\mu_{\ell^\nu})/F\). Fix a prolongation of every prime
in \(X\cup Q\), so that its decomposition group lies in
\(U\). At \(v\in Q\) take a uniformizer \(\pi_v\), a lift
\(\phi_v\) of arithmetic Frobenius, and a tame inertia
element \(\tau_v\) sending a chosen \(\ell^\nu\)-th root
of \(\pi_v\) to \(\zeta_{\ell^\nu}\) times that root.

Let \(\mathcal O\) be the ring of integers of a finite extension
of \(\QQ_\ell\). Let \(z\) be a continuous cocycle of \(G_E\)
with coefficients in \((\mathcal O/\ell^\nu\mathcal O)(\kappa^{-e})\).
Suppose \(z(cgc^{-1})=z(g)\), that \(z\) is unramified outside
\(\ell\) and the primes above \(Q\), and that \(z\) vanishes on
\(\ker\kappa\) in the decomposition groups above \(\ell\).
At each chosen prime over \(X\), assume \(z\) vanishes on a
Frobenius lift. For every \(y\in D_\nu\) one then has
\[
             \sum_{v\in Q} M\,z(\tau_v)\,\psi_y(\phi_v)=0
                        \quad\text{in }\mathcal O/\ell^\nu\mathcal O .
\]
One may take \(M=2\) if \(\ell\) is odd or if every
\(\Gamma_v\) for \(v\mid2\) is torsion-free at \(\ell=2\).
In all other dyadic cases one may take \(M=4\).
\end{proposition}
\begin{proof}
First take \(\mathcal O=\ZZ_\ell\) and put \(m=\ell^\nu\).
Since \(G_F=G_E\rtimes\langle c\rangle\), the formula
\(\widetilde z(gc^j)=z(g)\) for \(g\in G_E\), \(j=0,1\),
defines a cocycle for \(A_\nu\). Indeed, multiplication
conjugates the second \(G_E\)-component by \(c^j\), which
leaves its \(z\)-value unchanged, and \(\kappa_+(c)=1\).
At any real conjugation \(c'\), its square is \(1\) and it
acts trivially on \(A_\nu\). Hence \(2\widetilde z(c')=0\).
For \(g\) in a decomposition group at \(\ell\) with
\(\kappa_+(g)=1\), one has \(g^2\in G_E\), \(\kappa(g^2)=1\),
and \(2\widetilde z(g)=z(g^2)=0\) by the local hypothesis.
It follows that \(2\widetilde z\) descends locally to a
cocycle on \(\Gamma_v\).

At a dyadic place with \(-1\in\Gamma_v\), denote this last
cocycle by \(x\). For \(e\) even, the torsion involution acts
trivially, so \(2x\) vanishes on it and descends to the
quotient. For \(e\) odd, put \(t=-1\in\Gamma_v\). Evaluating
at \(tg=gt\) gives \(2x(g)=(g-1)(-x(t))\), so \(2x\) is
a coboundary. These are precisely the two exceptional local
conditions. Away from \(\ell\) the inertia of \(F\) lies in
\(G_E\), since \(E/F\) is unramified at finite places. Thus
\(M\widetilde z\) satisfies all the primary local conditions
outside \(Q\), including at \(X\), where both inertia and
a Frobenius lift have zero value.

The global cup product of \(M\widetilde z\) with \(y\) has
sum of local invariants zero by the Brauer reciprocity law.
Outside \(Q\) each invariant is already zero by orthogonality.
At \(v\in Q\) both modules are locally trivial and \(y\) is
unramified by Theorem~\ref{thm:dualrelations712}. Identify
\(\mu_m\) with \(R_\nu\) using the fixed \(\zeta_m\). The
Kummer character of \(\pi_v\) takes the value \(1\) at
\(\tau_v\). Subtracting \(z(\tau_v)\) times that character
from \(z\) gives an unramified local character. Indeed,
wild inertia has no \(\ell\)-quotient, and the maximal
exponent-\(m\) abelian \(\ell\)-quotient of tame inertia is
cyclic of order \(m\), generated by the image of \(\tau_v\).
The unramified part pairs to zero with \(y\). Local reciprocity
pairs the Kummer class of \(\pi_v\) with an unramified character
by evaluating that character at arithmetic Frobenius. In the
invariant group \(\QQ/\ZZ\) this is \(\psi_y(\phi_v)/m\), up
to the common sign from the ordering of the two cup factors.
Thus the sum of the remaining invariants gives the asserted
identity in \(R_\nu\). The sum is taken over \(F\), so there
is one term for each \(v\), with the chosen prolongation fixing
both local coordinates.~\cite{II-II-II-duality13}

Finally \(\mathcal O\) is finite free over \(\ZZ_\ell\).
Each coordinate of \(z\) in a \(\ZZ_\ell\)-basis is an
\(R_\nu\)-valued cocycle satisfying the same hypotheses because
the action is scalar. Applying the identity to each coordinate
proves it over \(\mathcal O/\ell^\nu\mathcal O\).
\end{proof}

\begin{corollary}[Dual classes on the fields of the period intervals]
\leavevmode\label{cor:fieldrelations712}
At every node of every sufficiently distant interval of
Theorem~\ref{thm:conductor232}, and on each field used in the
shorter construction of Corollary~\ref{cor:oddwindow677},
one has \(\dim W>|X|\). This holds for every positive weight,
every prime and every precision. On the shorter family the
reciprocity identity always holds with \(M=2\), and on the
original intervals it holds with \(M\le4\).
\end{corollary}
\begin{proof}
On the original intervals
\(E=K_\psi C_J(\sqrt{-q},\sqrt{-d_0})\) is abelian over
\(\QQ\) and contains two imaginary quadratic fields of
coprime discriminants \(-q\) and \(-d_0\). At any rational prime
at least one of these fields is unramified. Inertia in
\(\Gal(E/\QQ)\) fixes that field, while complex conjugation
acts nontrivially on it. Relative inertia for \(E/F\) is the
intersection with the subgroup generated by complex conjugation.
It is therefore trivial. Moreover \(F\) contains the real
cyclotomic layer \(C_J\) of degree \(2^J\). The lower bound
for \([K_\psi C_J:\QQ]\) in Theorem~\ref{thm:conductor232} and
the fixed degree of \(K_\psi\) imply \(2^J\ge8\) at every node
as soon as \(V\) is sufficiently large. Hence even the weakest
dimension bound in Theorem~\ref{thm:dualrelations712} exceeds
\(|X|\) on the whole interval.

On the shorter family \(E=\QQ(\sqrt{-d},\mu_{q_1})\), the
integer \(d\) is a prime congruent to \(7\bmod 8\), and
\(5040\mid q_1\). For the latter divisibility, \(18480\mid q_0\)
and \(q_1\) is a multiple of \(3q_0\), with \(5040\mid3\cdot18480\).
Here the two imaginary quadratic discriminants \(-4\) and
\(-d\) are coprime, giving the same unramified conclusion.
The inclusion \(\mu_{16}\subset E\) gives \(f\ge4\). Since
\(\mu_4\subset E\), each dyadic \(\Gamma_v\) lies in the
torsion-free group \(1+4\ZZ_2\). The sharper bounds in the
theorem and proposition apply at every node of this family.
\end{proof}

For example, suppose \(k\le |X|\) homomorphisms
\(U\to R_\nu\) have already been chosen with linearly
independent reductions. The corollary gives a dual class whose
restriction has reduction outside their span. The resulting
\(k+1\) homomorphisms together map \(U\) onto \(R_\nu^{k+1}\).
Indeed their image \(H'\) surjects onto the reduction of this
module, since a proper subspace over \(\mathbb F_\ell\) has
a nonzero annihilating linear form. It follows that
\(R_\nu^{k+1}/H'\) equals its own multiple by \(\ell\), and
iteration \(\nu\) times gives zero. In particular there is a
single element on which the previous homomorphisms vanish at
full precision and the new one takes value \(1\).

The dimension bound counts available relations. To deduce the
integral estimate of Section~\ref{sec:missing}, one still has
to control their evaluations at the primes used in the
projection and compare the resulting denominators with the
actual theta periods on the same interval.
\clearpage


\subsection{Auxiliary primes for the tame values of even cocycles}
\label{sec:triangularprimes713}

Keep \(E,F,\ell,c,\kappa\) as in
Corollary~\ref{cor:fieldrelations712}. Fix \(t\ge2\) and
\(\gamma\in G_E\) such that \(u=\kappa(\gamma)\) has infinite order.
For \(0\le e<t\) define the nonzero \(\ell\)-adic integer
\[
        \Delta_e=\prod_{\substack{0\le d<t\\d\ne e}}(u^e-u^d).
\]
Fix \(\nu\ge1\), put \(m=\ell^\nu\), \(R=\ZZ/m\ZZ\), and write
\(U=G_{E(\mu_m)}\). The word \emph{even} for a cocycle \(z\) of
weight \(e\) means that the action on its coefficients is
\(\kappa^{-e}\), that \(z(chc^{-1})=z(h)\), and that \(z\)
vanishes on \(\ker\kappa\) in all decomposition groups at
\(\ell\). Any finite ramification permitted away from \(\ell\)
will be specified separately. We use the multiplier \(M\) of
Proposition~\ref{prop:evenreciprocity712}, so \(M=2\) on the
shorter family and \(M\le4\) on the original period intervals.

\begin{lemma}[Separation of the weights]\leavevmode
\label{lem:separateweights713}
Suppose given continuous homomorphisms \(\psi_i:U\to R\) with
weights \(0<e_i<t\), satisfying
\[
 \psi_i(\gamma h\gamma^{-1})=u^{e_i+1}\psi_i(h),
 \qquad \psi_i(chc^{-1})=-\psi_i(h).
\]
Let \(\psi:U\to R\) satisfy the same identities with weight \(e\).
Suppose the reductions of \(\psi\) and of all \(\psi_i\) with
\(e_i=e\) are linearly independent over \(\mathbb F_\ell\).
For each \(s\in G_E\) with \(\kappa(s)=1\) there is \(g\in U\)
such that
\[
 \psi_i(g)=0\quad\text{for all }i,\qquad
 \psi(g)=2\Delta_e,\qquad z(g)=2z(s)
\]
for every even cocycle \(z\) of weight \(e\) with coefficients
in \((\mathcal O/m\mathcal O)(\kappa^{-e})\), where
\(\mathcal O\) is any ring of integers finite over \(\ZZ_\ell\).
All three identities persist if \(g\) is replaced by \(\delta^m g\)
for any \(\delta\in G_E\) with \(\kappa(\delta)=1\).
\end{lemma}
\begin{proof}
The homomorphisms of weight \(e\), including the new one, jointly
map \(U\) onto the corresponding product of copies of \(R\).
Their reductions are independent, so the image is surjective
after reduction. The quotient of \(R^k\) by that image is then
equal to its own multiple by \(\ell\) and hence is zero. Choose
\(h\in U\) on which the old maps of weight \(e\) vanish and the
new map takes the value \(1\).

For each \(0\le d<t\), \(d\ne e\), replace the current \(h\) by
\[
       (\gamma h\gamma^{-1})^{a}h^{b},
       \qquad a\equiv u^{-1}\pmod m,\quad b\equiv-u^d\pmod m .
\]
Take ordinary integer representatives for the exponents. The
element stays in \(U\). Its value under a map of weight \(k\)
is multiplied by \(u^k-u^d\). Thus after all replacements the
result \(h'\) is killed by every old map, and
\(\psi(h')=\Delta_e\). Set
\[
         g=s(csc^{-1})h'(ch'c^{-1})^{-1}.
\]
The four factors lie in \(U\), on which all the maps are
additive. The two terms involving \(s\) cancel for the dual
maps and add for \(z\). The two terms involving \(h'\) add
for the dual maps and cancel for \(z\). This proves the
formulas. For the last assertion, \(\delta\in U\), hence
\(m\) times its value vanishes under each of the homomorphisms
in question, including \(z|_U\).
\end{proof}

\begin{theorem}[A triangular bound at distinct auxiliary primes]
\leavevmode\label{thm:auxtriangle713}
Fix any finite ordered list \((e_j,s_j)_{1\le j\le b}\), with
\(0<e_j<t\) and \(s_j\in\ker(\kappa|_{G_E})\). There exist
primes \(w_j\) of \(E\) of rational prime norm, with chosen
arithmetic Frobenius lifts \(\phi_j\) and normalized tame
generators \(\tau_j\) as in Proposition~\ref{prop:evenreciprocity712}.
Their underlying rational primes may be chosen pairwise
distinct and outside any given finite set. The primes \(v_j\)
of \(F\) below them split completely in \(E(\mu_m)/F\),
so \(Nw_j\equiv1\bmod m\). They have these two properties.

First, every even cocycle of weight \(e_j\) unramified
outside \(\ell\) satisfies \(z(\phi_j)=2z(s_j)\).
Second, allow an even cocycle of weight \(e_j\) to ramify
at all primes of \(E\) over the \(v_i\), as well as at
\(\ell\). If
\[
   z(\tau_i)=0\quad(i<j),\qquad
   z(\phi_i)=0\quad(i<j,\ e_i=e_j),
\]
then \(2M\Delta_{e_j}z(\tau_j)=0\).
Both assertions hold with the same primes for all coefficient
rings \(\mathcal O/m\mathcal O\). In particular the exponent
\(v_\ell(2M\Delta_{e_j})\) in this bound is independent of
\(\nu\).

The primes may also obey an additional condition on the pair
\((w_j,\phi_j)\) under the following hypothesis. For every
\(g\in U\), there is \(\delta\in\ker(\kappa|_{G_E})\) such
that every open neighborhood of \(\delta^m g\) contains
Frobenius lifts satisfying the condition at infinitely many
degree-one primes of \(E\).
\end{theorem}
\begin{proof}
Choose the primes and dual maps successively. At index \(j\)
put \(e=e_j\) and \(X=\{v_i:i<j,\ e_i=e\}\). By
Corollary~\ref{cor:fieldrelations712} the space of residual
restrictions of the dual classes at this \(e,X,\nu\) has
dimension greater than \(|X|\). Choose a class \(y_j\) whose
restriction \(\psi_j\) is independent modulo \(\ell\) of all
the earlier maps of weight \(e\). This also preserves the
independence required at the next occurrence of that weight.
By Theorem~\ref{thm:dualrelations712}, \(y_j\) is unramified
away from \(X\) and \(\ell\), and every map so obtained has
the two equivariances of Lemma~\ref{lem:separateweights713}.
That lemma gives \(g\in U\) which kills all the earlier maps,
takes the value \(2\Delta_e\) under the new one, and has the
required evaluations for the even cocycles at \(s_j\).

These equalities are open conditions when imposed on the
Frobenius for all the even cocycles unramified outside \(\ell\).
To see this, first use coefficients \(R\). A cocycle gives a
continuous homomorphism into the finite affine group
\(R\rtimes R^\times\) by adjoining its scalar action. Its
kernel fixes a Galois extension of bounded degree unramified
outside \(\ell\). There are finitely many such extensions by
Hermite--Minkowski finiteness for bounded degree and a fixed
ramification set. In this form, finiteness follows by bounding
the local discriminants at the permitted primes, where a
nonarchimedean local field has finitely many extensions of
bounded degree, and then applying the discriminant form of
Hermite--Minkowski. Each resulting finite Galois group has
finitely many maps to the affine group. This proves finiteness
of the cocycles themselves. All their evaluations, together
with the finitely many \(\psi_i\) and the condition of lying
in \(U\), are therefore constant on an open neighborhood of
\(g\). For a general coefficient ring take coordinates in a
\(\ZZ_\ell\)-basis of \(\mathcal O\). Each coordinate is one
of the \(R\)-valued cocycles already considered. If the
additional condition is imposed, use the same argument around
\(\delta^m g\), which has the identical prescribed evaluations.

For the unrestricted choice, take a coset of an open normal
subgroup of \(G_E\) inside the neighborhood of \(g\). The
finite Chebotarev theorem gives positive Dirichlet density
for the primes with that Frobenius conjugacy class, and a
choice of prolongation realizes the chosen representative.
The contribution of primes of \(E\) of residue degree at
least two is bounded as the real variable \(\sigma\) tends
to \(1\) from above, by \([E:\QQ]\sum_p p^{-2\sigma}\). Thus
there are infinitely many degree-one choices. In the restricted
case use the stated neighborhood hypothesis. In either case
exclude the rational primes below the previous choices, the
given exceptional primes, the divisors of \(m\), and the
rational primes ramified in \(E(\mu_m)/\QQ\). Degree one in
\(E\) implies residue degree one in \(E/F\), so the prime
below \(w_j\) splits in that quadratic extension. Its Frobenius
also fixes \(\mu_m\) by the choice of neighborhood. This proves
the full splitting assertion. Use the chosen prolongation for
\(\tau_j\) and for \(\phi_j\). This completes the induction.
\cite[Theorem~8.31]{II-II-II-chebotarev713}

For the second assertion, \(z\) is unramified at every earlier
\(v_i\) in the following sense. At the chosen \(w_i\) the
action is trivial on inertia. Wild inertia has order prime
to \(\ell\) in each finite quotient, and the pro-\(\ell\)
part of tame inertia is procyclic. Hence \(z(\tau_i)=0\)
makes the inertia value zero at \(w_i\). At its conjugate
the value also vanishes by evenness. Changing a prolongation
preserves vanishing, by the cocycle identity and the trivial
action of inertia on the coefficients. The same identity on
the whole decomposition group at a split \(v_i\) justifies
using the Frobenius lift chosen above. Thus the assumptions
give exactly the conditions at \(X\) of
Proposition~\ref{prop:evenreciprocity712}. Apply that proposition
to \(y_j\) and \(Q=\{v_i:i\ge j\}\). All terms with \(i>j\)
vanish since \(\psi_j(\phi_i)=0\) by construction. The term
at \(j\) is \(2M\Delta_{e_j}z(\tau_j)\), as claimed.
\end{proof}

\begin{corollary}[A unit determinant at the same primes]\leavevmode
\label{cor:scalarfrob713}
Let \(T:G_E\to\mathrm{GL}_r(\mathcal O')\) be a continuous
representation, where \(\mathcal O'\) is the ring of integers
of a finite extension of \(\QQ_\ell\). Suppose some element
\(\eta\in\ker\kappa\) acts modulo the maximal ideal as a scalar
\(\alpha\) of order greater than \(r\). Then every prime in the
theorem can also satisfy
\(\det(1-T(\phi_j))\in(\mathcal O')^\times\).
\end{corollary}
\begin{proof}
For given \(g\in U\), the \(r+1\) elements
\((\eta^a)^m g\), \(0\le a\le r\), have residual matrices
\(\alpha^{am}\overline T(g)\). The scalars are distinct
because the order of \(\alpha\) is prime to \(\ell\).
The polynomial \(\det(1-X\overline T(g))\) has constant
term \(1\) and degree at most \(r\), so at some such scalar
its value is nonzero. Set \(\delta=\eta^a\). Continuity makes
the determinant a unit on an open neighborhood of \(\delta^m g\).
Intersecting with any prescribed smaller neighborhood and
applying the same finite Chebotarev argument proves the
additional-condition hypothesis.
\end{proof}

The auxiliary primes here are chosen anew at each precision.
The annihilator \(2M\Delta_e\) is fixed before that choice.
A bound on the total loss also requires controlling the size
of a family of points \(s_j\) that detects all relevant
cocycles.
\clearpage


\subsection{The cost of the auxiliary primes on a period interval}
\label{sec:auxiliarycost714}

Let \(E\) be one of the period fields in
Corollary~\ref{cor:fieldrelations712} and suppose
\(E\subseteq\QQ(\mu_N)\). Fix an integer \(t\ge2\). Theorem
\ref{thm:auxtriangle713} allows us to choose one auxiliary
prime for each pair consisting of a positive weight less
than \(t\) and a test point. We now bound how many tests
are needed and how the resulting loss varies with \(\ell\).

For odd \(\ell\) let \(K=E(\mu_\ell)\), and choose
\(\gamma\in G_E\) whose cyclotomic image \(u=\kappa(\gamma)\)
generates \(C=\kappa(G_E)\) topologically. This is possible
since \(C\) is open in \(\ZZ_\ell^\times\) and is the product
of a cyclic group of order prime to \(\ell\) and a copy of
\(\ZZ_\ell\), which together are procyclic. At two take
\(K=E(i)\) and take \(\gamma\in G_K\) with \(u\) generating
\(C\cap(1+4\ZZ_2)\), the full cyclotomic image of \(G_K\).
This is an open subgroup of \(1+4\ZZ_2\). In particular
\(\gamma\) generates the whole cyclotomic image at two in the
shorter construction.

\begin{lemma}[Test points independent of precision and weight]
\leavevmode\label{lem:testpoints714}
There is a finite set \(S_\ell\subseteq\ker(\kappa|_{G_E})\)
of size at most \(2[K:\QQ]+\log_\ell h_K\), where \(h_K\)
is the ordinary class number. For any weight \(e\ge0\),
precision \(\nu\ge1\), and finite coefficient ring of integers
\(\mathcal O\) above \(\ell\), let \(z\) be a cocycle in
\((\mathcal O/\ell^\nu\mathcal O)(\kappa^{-e})\) unramified
outside \(\ell\). If \(z(\gamma)=0\) and \(z\) vanishes on
\(S_\ell\), then \(z=0\) when \(\ell\) is odd or \(i\in E\).
In the other dyadic case these vanishing conditions imply
\(z|_{G_K}=0\) and \((u^{-e}-1)z=0\). Testing also one
representative of the nonidentity coset in \(G_E/G_K\)
then detects \(z\) exactly.
\end{lemma}
\begin{proof}
Let \(P_K\) be the Galois group of the maximal pro-\(\ell\)
extension of \(K\) unramified outside the primes over \(\ell\).
The space of continuous homomorphisms
\(P_K\to\mathbb F_\ell\) has dimension at most
\(2[K:\QQ]+\log_\ell h_K\). Indeed, \(K\) contains
\(\mu_\ell\), so Kummer theory identifies this space with a
subgroup of \(K^\times/K^{\times\ell}\). Away from \(\ell\)
a representing element has valuation a multiple of \(\ell\),
because an unramified extension has integral value group.
Sending the prime-to-\(\ell\) divisor divided by \(\ell\)
to the class group of \(\mathcal O_K[1/\ell]\) gives an
\(\ell\)-torsion class. The kernel is represented by units
of \(\mathcal O_K[1/\ell]\). Indeed, a principal divisor
at this stage can be removed by an \(\ell\)-th power. The
unit theorem, including its roots of unity, bounds the unit
dimension modulo \(\ell\)-th powers by \(2[K:\QQ]\). The
open ring class group is a quotient of the ordinary class
group, giving the remaining term.

Choose elements \(t_1,\ldots,t_d\) in \(G_K\) whose images
in the maximal elementary abelian quotient of \(P_K\) form
a basis. Their images generate \(P_K\) topologically. For
otherwise the closed subgroup they generate would have proper
image in some finite \(\ell\)-group quotient, hence lie there
in a maximal subgroup of index \(\ell\). The quotient
character would contradict the basis property. For each \(i\)
choose \(h_i\in\overline{\langle\gamma\rangle}\) with
\(\kappa(h_i)=\kappa(t_i)\). The closed subgroup generated by
\(\gamma\) maps onto the closed subgroup generated by \(u\)
by compactness, and this includes \(\kappa(G_K)\) in every
case. Take \(S_\ell=\{t_i h_i^{-1}:1\le i\le d\}\).

A cocycle zero at \(\gamma\) is zero on its closed cyclic
subgroup by the cocycle identity and continuity. Hence the
tests imply \(z(t_i)=0\) for each \(i\). Over \(K\) the
homomorphism
\[
   g\longmapsto \bigl(z(g),\,\kappa(g)^{-e}\bmod\ell^\nu\bigr)
\]
into the group of affine transformations has \(\ell\)-group
image. Its translations have \(\ell\)-power order and its
linear parts are principal units modulo \(\ell^\nu\).
The homomorphism is unramified off \(\ell\), so it factors
through \(P_K\). The images of all the \(t_i\) fix zero.
They generate the image, whence \(z|_{G_K}=0\). For odd
\(\ell\) the quotient \(\Gal(K/E)\) is generated by the
image of \(\gamma\). For \(\ell=2\) and \(i\in E\) the
quotient is trivial. This proves exact detection in both
cases. In the remaining case \(K/E\) is quadratic. For
\(a\in G_K,\ g\in G_E\), evaluation at \(ag=g(g^{-1}ag)\)
gives \(\kappa(a)^{-e}z(g)=z(g)\). Take \(a=\gamma\).
The cocycle identity gives the last assertion on testing
one more coset.
\end{proof}

\begin{theorem}[One integral bound for the tame tests on an interval]
\leavevmode\label{thm:trianglebudget714}
For each \(\ell\) use the above choice of \(u\) and put
\[
 h_\ell=\max_{1\le a<t}v_\ell(u^a-1),\qquad
 b_\ell=(t+2)h_\ell,\qquad
 c_\ell=(2+t|S_\ell|)\,b_\ell .
\]
Here \(v_\ell(\ell)=1\). In Theorem~\ref{thm:auxtriangle713}
take one prime for every pair \((e,s)\) with \(0<e<t\),
\(s\in S_\ell\), in increasing order of weight and any fixed
order within a weight. The resulting exponent needed at
each tame generator is at most \(b_\ell\). For all \(0\le e<t\)
and all \(1\le a<t\), respectively, the following divisibilities
hold in \(\ZZ_\ell\), as follows.
\[
          \Delta_e\mid\ell^{b_\ell},\qquad
          4(u^a-1)^2\mid\ell^{b_\ell}.
\]
The primes may vary with \(\nu\) while the test points and all
the displayed bounds are fixed. If \(Y\) is any integer with
\(N+t<Y\), then \(c_\ell\le Y^6 h_\ell\), and \(h_\ell=0\)
for \(\ell>\max(N,t)\).

In particular suppose a common interval has
\(E\subseteq\QQ(\mu_{N_*})\) and \(2\le t\le t_*\) throughout.
Put \(Y_*=N_*+t_*+1\), \(g_*=\varphi(N_*)\), and choose
primes \(N_*<P_1<2N_*<P_2<4N_*\) for \(N_*\ge3\).
The positive integer
\[
       H_*=\prod_{a=1}^{t_*-1}
              (P_1^{2g_*a}-1)(P_2^{2g_*a}-1)
\]
satisfies, at every node,
\[
  c_\ell\le Y_*^6 v_\ell(H_*),\qquad
  \sum_{\ell} c_\ell\log\ell
       \le Y_*^6\log H_*
       \ll Y_*^6 N_*t_*^2\log(4N_*).
\]
\end{theorem}
\begin{proof}
Each difference \(u^e-u^d\), \(d\ne e\), is a unit times
\(u^{|e-d|}-1\), up to sign. Thus every product \(\Delta_e\)
has valuation at most \((t-1)h_\ell\). At odd primes \(2M\)
is a unit. At two we chose \(u\in1+4\ZZ_2\), so \(h_2\ge2\)
and \(v_2(2M)\le3\le2h_2\). Consequently
\(v_\ell(2M\Delta_e)\le b_\ell\), and
\(v_\ell(4(u^a-1)^2)\le3h_\ell\le b_\ell\).
This proves both divisibilities and the tame bound. There are
\((t-1)|S_\ell|\) primes and \(|S_\ell|\) kernel test points.
The additional two units in \(2+t|S_\ell|\) allow for testing
\(\gamma\) and the extra coset when it occurs. If \(\ell=2\)
and that extra coset is omitted, the remaining positive-weight
cocycle is killed by \(\ell^{h_2}\), by
Lemma~\ref{lem:testpoints714}.

For \(\ell>\max(N,t)\), the cyclotomic \(\ell\)-power fields
are disjoint from \(\QQ(\mu_N)\). At every finite level this
follows from the Chinese remainder isomorphism of Galois groups
at the relatively prime moduli. Thus \(u\bmod\ell\) generates
\(\mathbb F_\ell^\times\), of order \(\ell-1\ge t\). All the
factors \(u^a-1\) are units. For the remaining primes the
field \(K\) lies in a cyclotomic field of modulus \(L\le Y^2\),
using the modulus \(N\ell\) at odd primes and \(4N\) at two.
Also \(d=[K:\QQ]\le Y^2\), by the degree bound for this
cyclotomic field. The absolute discriminant of \(K\) is at
most \(L^d\), by the cyclotomic discriminant formula and
the relative discriminant formula in a tower. Minkowski's
bound gives an integral representative of each ideal class
of norm at most \(B=2^d L^{d/2}\). Counting ideals by their
inverse squared norm bounds their number up to \(B\) by
\(B^2\zeta_K(2)\le B^2\zeta(2)^d\). Therefore
\(\log_2 h_K\le d(3+2\log_2Y)\), and the lemma gives
\(|S_\ell|\le(5+2\log_2Y)Y^2\le Y^4\) in this prime range.
Since \(t+2\le Y\) and \(2+tY^4\le Y^5\), the cost bound
follows. When \(h_\ell=0\) the cost is already zero.

To get a single integer at all nodes, the index of
\(\kappa(G_E)\) in \(\ZZ_\ell^\times\) divides
\([E:\QQ]\), and this degree divides \(g_*\). The first
divisibility follows by taking the preimage of \(\kappa(G_E)\)
in \(G_\QQ\) and using the index formula and surjectivity of
\(\kappa|_{G_\QQ}\). For \(P\) one of \(P_1,P_2\) different
from \(\ell\), it follows that \(P^{2g_*}\) is a limit of
powers of \(u\). At two its square exponent also guarantees
\(P^{2g_*}\equiv1\bmod4\), so it belongs to the group generated
by our chosen \(u\) there as well. For integer powers \(u^j\),
including negative powers, \(u^a-1\) divides \(u^{aj}-1\).
The principal ideal it generates is closed in \(\ZZ_\ell\).
Passing to the limit gives \(u^a-1\mid P^{2g_*a}-1\).
The primes \(P_1,P_2\) exist by Bertrand's postulate. Taking
valuations proves \(h_\ell\le v_\ell(H_*)\), and prime
factorization gives the sum bound. The asserted estimate
on its logarithm follows by summing
\(4g_*a\log(4N_*)\) over \(a<t_*\).
\end{proof}

\begin{corollary}[Comparison with the two actual level groups]
\leavevmode\label{cor:actualmarks714}
On the shorter family, write \(r\) for the prescribed rank,
\(t_0=5r\), and \(p_0,p_1\) for its two level primes. At a node
put \(N=q_1d\), \(f=[E^+:\QQ]\) and
\[
       G_0=\prod_{i=0}^1
                |\mathrm{GL}_{rt_0}(\mathbb F_{p_i})|^f .
\]
For each \(2\le t\le t_0\) and integer \(Y>N+t\), all the
tame tests above have \(c_\ell\le Y^6 v_\ell(G_0)\). In
particular their total weighted cost is at most
\[
             Y^6\log G_0 \le Y^6 f (rt_0)^2\log(p_0p_1).
\]
The detection in Lemma~\ref{lem:testpoints714} is exact at
every coefficient prime of this family.
\end{corollary}
\begin{proof}
At a chosen place over \(p_i\ne\ell\), arithmetic Frobenius
has cyclotomic value \(p_i\), because \(p_i\) splits completely
in \(E\) by Section~\ref{sec:levelmass678}. Here \(u\) generates
the whole cyclotomic image, including at two because \(\sqrt{-1}\)
belongs to \(E\). Hence \(u^a-1\) divides \(p_i^a-1\) by the
closed-ideal argument just used. Counting ordered bases gives
\[
    |\mathrm{GL}_d(\mathbb F_p)|
        =p^{d(d-1)/2}\prod_{a=1}^d(p^a-1) \le p^{d^2}.
\]
For \(1\le a<t\le t_0\), the factor at that \(a\) occurs in
the group of dimension \(rt_0\). The two level primes are
distinct, so there is always a choice different from \(\ell\).
This proves \(h_\ell\le v_\ell(G_0)\) and both estimates.
\end{proof}

On the original intervals, Theorem~\ref{thm:conductor232} and
Corollary~\ref{cor:oddwindow677} provide the single modulus
\(N_*\) and weight cutoff \(t_*\), polynomial in \(V\), that
the theorem requires. On the shorter family \(q_1=3q_0l_0 2^s\),
where \(l_0\) is the least power of three at least \(27s\).
For fixed \(q_0,d\) these moduli divide the modulus taken at
the largest \(s\) on the interval. The bounds on \(N\) and
\(t\) in Corollary~\ref{cor:oddwindow677} are polynomial there
as well. Thus the tame cost and the previously bounded odd
cocycle contribution are available on the same nodes and
intervals. The estimate on \(c_\ell\) counts the valuation
loss in the cohomological tests. Its use in an integral
projection still requires the congruence argument for the
Hecke action on the theta lifts.

For the geometric Frobenius convention in that action, apply
Theorem~\ref{thm:auxtriangle713} at \(s^{-1}\) for each test
point \(s\), and use the inverse of its arithmetic Frobenius.
All these elements act trivially on the finite twisted
coefficients. Inversion negates the value of a cocycle there,
so the test at geometric Frobenius gives \(2z(s)\). It changes
the sign of the diagonal dual value and leaves its valuation
unchanged. Testing \(\gamma^{-1}\) is equivalent to testing
\(\gamma\) by the cocycle identity.
\clearpage


\subsection{Scalar cocycles after the paired inertia and Frobenius tests}
\label{sec:pairedscalar717}

Fix a field \(E\), a prime \(\ell\), an exponent \(0\le e<t\)
and finite precision \(\nu\ge1\) as in
Theorem~\ref{thm:trianglebudget714}. Write
\(A=\mathcal O/\ell^\nu\mathcal O\), where \(\mathcal O\) is
the ring of integers of a finite extension of \(\QQ_\ell\).
The coefficient character is \(\chi=\kappa^{-e}\).
At a prime above \(\ell\), consider the following elements
of the absolute Galois group. For \(e>0\) take the kernel of
the full cyclotomic character in its decomposition group.
For \(e=0\) take its inertia group. Write \(L_e\) for the
union of these sets over all places above \(\ell\) and all
their prolongations. This set is stable under \(g\mapsto cgc^{-1}\).

\begin{lemma}[The local crystalline evaluation equations]
\leavevmode\label{lem:cryszero717}
At a place \(v\mid\ell\), each integral crystalline cocycle
of \(\mathcal O(-e)\) vanishes on \(L_e\) in its chosen local
group. The same evaluation equations hold for a finite
coefficient cocycle satisfying all integral linear relations
among the evaluations of the local crystalline cocycles.
For a global cocycle, the equations at one prolongation of
each \(v\) imply the equations at all prolongations.
\end{lemma}
\begin{proof}
Write \(K=E_v\). The Bloch--Kato computation gives
\(H^1_f(K,\QQ_\ell(-e))=0\) for \(e>0\). At \(e=0\) this
group consists of the unramified homomorphisms. In particular
the computation is over the actual completion \(K\), with
arbitrary absolute ramification. Scalar extension to the
fraction field of \(\mathcal O\) preserves these descriptions
by taking coordinates in a \(\QQ_\ell\)-basis.
For \(e>0\) the integral cocycle is therefore of the form
\(g\mapsto(\chi(g)-1)a\) over the fraction field and has
value zero when \(\kappa(g)=1\). For \(e=0\) it is zero
on inertia. The integral lattice embeds in its fraction
space, giving these equalities integrally.
~\cite[Corollary~3.8.4 and Example~3.9]{II-II-II-bktate717}

Each such single evaluation with coefficient one is itself
an integral linear relation. Finally the cocycle identity gives
\[
 z(hgh^{-1})=\chi(h)z(g)+(1-\chi(g))z(h).
\]
The elements being tested have \(\chi(g)=1\), proving the
assertion for the conjugate prolongations as well.
\end{proof}

Use the primes of Theorem~\ref{thm:trianglebudget714} and put
\(I_e=\{j:e_j=e\}\), so \(I_0\) is empty. For any auxiliary
element \(a\) put \(a'=(cac^{-1})^{-1}\).
Write \(C=\ell^{b_\ell}\). Let \(Z\) be an \(\mathcal O\)-module
of continuous \(A(\chi)\)-valued cocycles unramified outside
\(\ell\) and the pairs \(w_j,\bar w_j\). Write \(Z_L\) for
those vanishing on \(L_e\), and let \(Z_f\) consist of the
members of \(Z\) satisfying the local crystalline relations
in the lemma at every \(v\mid\ell\). All lengths below are
over \(\mathcal O\). These modules are finite. Indeed the
bounded-degree and finite-ramification-set argument in the
proof of Theorem~\ref{thm:auxtriangle713} applies to the
affine group of \(A(\chi)\).

\begin{theorem}[Loss of length after the paired tests]
\leavevmode\label{thm:scalarcost717}
For \(j\in I_e\), let \(z\in Z_L\) satisfy
\(z(\tau_i)=z(\tau'_i)\) for all \(i<j\) and
\(z(\phi_i)=z(\phi'_i)\) for all \(i<j\) in \(I_e\).
Then
\[
              C\bigl(z(\tau_j)-z(\tau'_j)\bigr)=0 .
\]
Now let \(S\subseteq Z_L\) be any submodule whose members
vanish at every \(\tau_j,\tau'_j\) and agree at
\(\phi_j,\phi'_j\) for \(j\in I_e\).
Write \(\overline S\) for its image in
\(H^1(G_E,A(\chi))\). Let \(O_e\) denote the strict odd
cocycles \(Z^-_\nu(E,e)\) with coefficients \(\mathcal O\)
if \(e>0\). For \(e=0\) put \(O_0=\operatorname{Hom}
(\operatorname{Cl}(E),A)\). Then
\[
       \len(Z/Z_L)+\len(C\overline S)
           \ \le\ \len(Z/Z_f)+\len O_e .
\]
The same inequality holds with \(CS\) in place of
\(C\overline S\) if \(e=0\), or if every \(z\in S\) satisfies
\((u^{-e}-1)z(\gamma)=0\), with \(u,\gamma\) as in
Lemma~\ref{lem:testpoints714}.
\end{theorem}
\begin{proof}
Write \(z^c(g)=z(cgc^{-1})\). This is a cocycle for the same
character since \(\kappa(cgc^{-1})=\kappa(g)\).
At every \(\tau_i\) and \(\phi_i\) the action on \(A\) is
trivial. The identity for inverses then gives
\[
         z(a)-z(a')=(z+z^c)(a)
           \qquad(a=\tau_i,\phi_i).
\]
For \(j\in I_e\) one has \(e>0\). Thus \(z+z^c\) is even,
has the allowed ramification, and vanishes on \(L_e\).
The earlier equalities in the theorem are exactly the vanishing
conditions of Theorem~\ref{thm:auxtriangle713}. Its multiplier
\(2M\Delta_e\) divides \(C\) by
Theorem~\ref{thm:trianglebudget714}. This proves the first assertion.

Each member of \(S\) is unramified away from \(\ell\).
At either member of an auxiliary pair inertia acts trivially
on the coefficients, so the cocycle on inertia is a homomorphism.
Wild inertia and the prime-to-\(\ell\) factors of tame inertia
map to zero. A generator of the remaining procyclic factor
and its conjugate inverse have both been tested. This proves
unramifiedness at the pair. In weight zero the condition on
\(L_0\) now gives an everywhere unramified homomorphism. Class
field theory for the totally imaginary field \(E\) embeds
\(S\) into \(O_0\), and its coboundaries are zero.

Suppose \(e>0\). Sending \(z\) to \(z^c-z\) maps \(S\)
to \(O_e\), by the ramification and local conditions just
checked. It sends every coboundary to zero because \(\chi\)
is invariant under conjugation by \(c\). Thus it induces
a map \(\overline S\to O_e\). An element in its kernel has
a representative \(z\in S\) with \(z^c=z\) as functions.
For each \(s\in S_\ell\), the reading prime for \((e,s)\)
therefore gives
\[
       4z(s)=2(z+z^c)(s)=(z+z^c)(\phi_j)=0 .
\]
Put \(D=u^{-e}-1\). To handle an arbitrary representative
define the cocycle
\[
            z_*(g)=D z(g)-(\chi(g)-1)z(\gamma).
\]
It is unramified outside \(\ell\), has \(z_*(\gamma)=0\),
and \(4z_*\) vanishes at every \(s\in S_\ell\), since
\(\chi(s)=1\). Lemma~\ref{lem:testpoints714} now yields
\(4D z_*=0\) in all cases. At the places of exact detection
one already has \(4z_*=0\). In cohomology
\([z_*]=D[z]\), so \(4D^2[z]=0\).
If \(Dz(\gamma)=0\) holds initially, one may apply the same
test lemma directly to \(4Dz\), giving \(4D^2z=0\) as a
function. Since \(4D^2\) divides \(C\) in \(\mathcal O\)
by Theorem~\ref{thm:trianglebudget714}, the corresponding
kernel is killed by \(C\). Multiplication by \(C\) consequently
factors through the image of the map to \(O_e\), and the
length of its image is at most \(\len O_e\). Finally
Lemma~\ref{lem:cryszero717} gives \(Z_f\subseteq Z_L\),
which proves the displayed bound.
\end{proof}

\begin{corollary}[The scalar costs on the period fields]
\leavevmode\label{cor:scalarsum717}
For weights \(0\le e<t\) allow different coefficient rings
\(\mathcal O_e\), with ramification indices \(a_e\), and
different precisions. Use a common set of auxiliary primes
at the largest precision. With \(E\subseteq\QQ(\mu_N)\),
the sum over \(e\) of
\[
 \frac{\len(Z/Z_L)+\len(C\overline S)-\len(Z/Z_f)}{a_e}
\]
is at most \(v_\ell(h_E\mathcal J_{N,t})\).
The cocycle version with \(CS\) holds under the corresponding
normalization hypotheses of the theorem. When \(N,t\) vary on
the period interval, the bounds of
Corollary~\ref{cor:oddwindow677} apply.
\end{corollary}
\begin{proof}
For precision \(\nu'<\nu\) the map of coefficient modules
\(a\bmod\ell^{\nu'}\mapsto\ell^{\nu-\nu'}a\bmod\ell^\nu\)
is an equivariant injection. Hence all evaluations and the
triangular annihilation at precision \(\nu\) hold at
\(\nu'\). The primes can simultaneously satisfy the determinant
condition of Corollary~\ref{cor:markeddet715} for its prescribed
representation. Apply Theorem~\ref{thm:oddinteger677} to
the bounds of Theorem~\ref{thm:scalarcost717} and sum.
\end{proof}

The estimate singles out the local crystalline colength
\(\len(Z/Z_f)\) in the congruence calculation. Its contribution
for the actual families of representations, and the corresponding
bound for their higher-rank summands, still enter the integral
period comparison.
\clearpage

\subsection{Crystalline cocycles at the spherical coefficient primes}
\label{sec:flrelations724}

The character determinants of Theorem~\ref{thm:charrows719} contain
the index \(l_e\) of a local condition inside a module defined by relations
among matrices. At coefficient primes in a common Fontaine--Laffaille
range, this index is zero at every auxiliary level. We first construct
the local extension that realizes a solution of the matrix relations.

\begin{lemma}[A lattice for the matrix cocycle]\leavevmode
\label{lem:graphrealization724}
Let \(G\) be a profinite group, let \(\mathcal O\) be the ring of integers
of a finite extension \(L/\QQ_\ell\), and let \(u_x\) be a finite family
of continuous representations on free \(\mathcal O\)-lattices \(U_x\).
Set \(D(g)=(u_x(g)-1)_x\). Suppose \(z:G\to T/\ell^\nu T\) is
a continuous cocycle, where \(T\) is a \(G\)-stable lattice. Assume that
every finite \(\mathcal O\)-linear relation among the \(D(g)\) holds among
the corresponding \(z(g)\). The extension defined by \(z\) is the
pushout of the reduction of an integral, rationally split extension
of \(\mathcal O\) by a lattice \(Y\). Rationally, \(Y\) is a direct sum
of copies of the \(U_x\otimes L\).
\end{lemma}
\begin{proof}
Let \(W=\bigoplus_x\operatorname{End}_{\mathcal O}(U_x)\), with \(g\)
acting on each matrix by multiplying it on the left by \(u_x(g)\).
Reading columns gives the asserted direct sum description of \(W\).
Let \(P\subset W\) be the \(\mathcal O\)-span of the \(D(g)\). The identity
\(hD(g)=D(hg)-D(h)\) makes \(P\) stable. The relations define a linear
map \(b:P\to T/\ell^\nu T\) sending \(D(g)\) to \(z(g)\). The cocycle
identity applied to \(hg\) makes \(b\) equivariant. By the elementary
divisor theorem there is \(s\) such that \(P\cap\ell^sW\subset\ell^\nu P\).
Indeed in a basis in which \(P\) has elementary divisors \(d_i\), one
can take \(\ell^s\) divisible by every \(\ell^\nu d_i\). Thus \(b\)
extends by zero on \(\ell^sW\) to \(Y=P+\ell^sW\). Its reduction on
\(Y/\ell^\nu Y\) is again equivariant.

On \(Y\oplus\mathcal O v\) put \(g v=v+D(g)\). This action is continuous
by the lattice topology and the cocycle identity. With \(I=(1_{U_x})_x\)
one has \(D(g)=gI-I\). The vector \(v-I\) therefore gives a section
over \(L\), since \(Y\otimes L=W\otimes L\). Reduction and pushout along
\(b\) give exactly \(g v=v+z(g)\). In particular the torsion extension
is a quotient of the direct sum of \((Y\oplus\mathcal O v)/\ell^\nu\)
and \(T/\ell^\nu T\), by the graph of the negative of \(b\).
\end{proof}

\begin{theorem}[Lifting the local cocycle at full precision]\leavevmode
\label{thm:fllift724}
Suppose \(K/\QQ_\ell\) is finite unramified, \(\ell>2\), and \(G=G_K\).
Suppose \(T\otimes L\) and all \(U_x\otimes L\) are crystalline. Assume
that their de Rham filtration jumps at all embeddings, together with \(0\),
lie in one interval of integers \([A,B]\) with \(B-A\le\ell-2\).
Use the convention that the cyclotomic character \(\kappa\) has
filtration jump \(-1\). Every cocycle in Lemma~\ref{lem:graphrealization724}
is the reduction of an integral cocycle defining a crystalline extension
of the trivial representation by \(T\otimes L\).
\end{theorem}
\begin{proof}
Write \(j=-A\). Twist the extensions by \(\kappa^{-j}\). Their jumps
now lie in \([0,B-A]\), and the quotient is the character of jump \(j\).
Use the covariant integral Fontaine--Laffaille functor for this interval.
It is exact and fully faithful, its essential image on torsion objects
is closed under subquotients, and every stable lattice in a crystalline
representation with these jumps is in its image. In the filtered
description a Tate character \(\kappa^{-j}\) has a generator in degree
\(j\) fixed by its \(j\)-th divided Frobenius. Morphisms are strict for
the filtration. These are the integral and torsion comparison theorems,
including their versions with coefficients.
~\cite{II-II-II-fl724,II-II-II-ghls724}

Write \(S=W(k_K)\otimes_{\ZZ_\ell}\mathcal O\), where \(k_K\) is the
residue field of \(K\), and \(S_\nu=S/\ell^\nu S\). The Frobenius on
these rings acts on the Witt vectors and fixes the coefficients. One
may enlarge \(L\) to contain all embeddings of \(K\). For descent,
\(\mathcal O\) is a direct summand of its finite ring extension as an
\(\mathcal O\)-module, since \(1\) is primitive. After restricting
coefficients, project a lifted cocycle back to \(T\) using a projection
sending \(1\) to \(1\). On rational extensions this is a pushout in the
category of crystalline representations. It preserves crystallinity
because this category is closed under direct sums and quotients.

By the preceding lemma the rational representation behind the graph
lattice is a sum of crystalline representations in the given interval.
Its reduction and that of \(T\) therefore belong to the Fontaine--Laffaille
image after the twist. Closure under quotients puts the pushout in that
same image. By full faithfulness and exactness the resulting filtered
module is an extension of the rank-one module \(S_\nu v_0\), of pure
degree \(j\), by the reduction \(M_\nu=M/\ell^\nu M\) of the filtered
lattice \(M\) for \(T(\kappa^{-j})\). More explicitly, kernels and images
of morphisms in the torsion category are carried to their kernels and
images on representations, so the two ends have the specified
identifications.

Strictness gives a lift \(v\) of \(v_0\) in filtration \(j\). The module
splits over \(S_\nu\) with this choice. In filtration degrees \(i\le j\)
one adds \(S_\nu v\) to \(\mathrm{Fil}^i M_\nu\), and in degrees \(i>j\)
one has precisely \(\mathrm{Fil}^i M_\nu\). Thus the only new Frobenius
datum is \(\phi_j(v)=v+a\), with \(a\in M_\nu\). Choose any
\(\widehat a\in M\) reducing to \(a\). On \(M\oplus S\widehat v\) impose
the same split filtration and set
\[
       \phi_i(\widehat v)=\ell^{j-i}(\widehat v+\widehat a)
                         \qquad(0\le i\le j),
\]
extending semilinearly. These maps satisfy the divided Frobenius
relations. Their images generate the module because the images from
\(M\) generate \(M\) and \(\phi_j(\widehat v)\) generates the quotient.
This is an integral Fontaine--Laffaille extension, and its reduction
with both ends identified is the given extension. Apply the functor
and undo the Tate twist. A section of the quotient gives an integral
cocycle. If its reduction differs from \(z\) by a coboundary, change
the section by a lift of the defining element of \(T/\ell^\nu T\).
The resulting cocycle reduces to \(z\) at every group element.
\end{proof}

\begin{corollary}[Simultaneous vanishing of the scalar colengths]\leavevmode
\label{cor:scalarfl724}
At a node on the period fields of Corollary~\ref{cor:fieldrelations712}
fix the definite groups, weights and splitting characters in
Theorem~\ref{thm:splitattachment722} and Corollary~\ref{cor:normalizedtransfer722}.
Let \(\ell>2\) be unramified in \(E/\QQ\), with hyperspecial level above
\(\ell\), and prime to the finite conductors of the splitting characters.
All auxiliary levels used here are away from \(\ell\).
For a group of rank \(n\) with algebraic weight \(a=(a_{\tau,i})\),
let \(h_\alpha\) be the largest absolute filtration jump of its
normalizing character \(\alpha\). Take the maximum \(H\), over the
groups under consideration, of
\[
             t-1,\qquad n+\max_{\tau,i}|a_{\tau,i}|+h_\alpha .
\]
If \(2H\le\ell-2\), then \(l_e=0\) for all \(0\le e<t\). The conclusion
applies to every auxiliary level and every precision. On any one finite
interval the maximum of these bounds over its nodes gives one threshold
for the same fields on which the periods are taken.
\end{corollary}
\begin{proof}
We check the local hypotheses for the actual automorphic systems.
The base change used in Theorem~\ref{thm:splitattachment722} is
spherical at a hyperspecial place, including an inert place.
Its local parameter is a direct sum of the parameters of its cuspidal
norm twists. Since a spherical parameter is trivial on inertia and
has zero monodromy, both properties hold for each of these summands.
Each normalized cuspidal factor
\(\pi_{ij}=\rho_i|\det|^{s}\), \(s=j+(a_i-n)/2\) in the notation of
that theorem (here \(a_i\) denotes the rank of \(\rho_i\)), is regular
algebraic in its own rank. It is also essentially conjugate self-dual.
Indeed the unitary cuspidal center \(\rho_i\) of its discrete string
satisfies \(\rho_i^c=\rho_i^\vee\) by uniqueness of its unitary
cuspidal support. The ratio needed to identify
\(\pi_{ij}^\vee\) with a character twist of \(\pi_{ij}^c\)
is \(|\det|_E^{-2s}\). This is the pullback of \(|\cdot|_F^{-2s}\)
by the norm, with value \(1\) at \(-1\) at each real place.

The spherical case of local-global compatibility for regular algebraic
essentially conjugate self-dual cuspidal representations now gives
crystallinity for each factor at every \(w\mid\ell\).
The representations so obtained agree with those in
Theorem~\ref{thm:splitattachment722} by their good Frobenius
polynomials and semisimplicity. In rank one the same assertion is the
crystallinity of the local avatar of an algebraic character with
unramified smooth part. In particular this applies to \(\alpha\).
The Hodge weight formula for the full direct sum in rank \(n\) has
entries \(i-n-a_{\tau,i}\) in Guerberoff's convention. In the opposite
sign convention the absolute values are the same. This formula follows
by joining the shifted infinitesimal characters of the cuspidal
summands, as in the proof of his Theorem~2.3. It depends only on the
algebraic infinity type. Tensoring by \(\alpha\) shifts the weights at
each embedding by its single filtration jump. The jumps of all the
normalized systems hence lie in \([-H,H]\).
~\cite[Thms.~0.1, 2.2 and proof of Thm.~2.3]{II-II-II-guerberoff724}

The coefficient lattice for weight \(e\) is \(\mathcal O(\kappa^{-e})\)
and has jump \(e\in[0,H]\). Every cocycle in \(\mathcal Z_e\) restricts
at each \(w\mid\ell\) to one obeying all the local matrix relations,
since local relations are among the global relations defining
\(\mathcal Z_e\). Apply Theorem~\ref{thm:fllift724} with \(K=E_w\).
The resulting integral crystalline lift shows that its reduction
satisfies every evaluation relation used in Lemma~\ref{lem:cryszero717}.
This proves the asserted equality of the two subgroups defining \(l_e\).
Auxiliary level away from \(\ell\) changes neither the infinity type nor
the hyperspecial local component condition. On a fixed interval there
are finitely many fields, nodes and weights, so the common maximum has
the stated effect at primes satisfying the indicated local conditions.
\end{proof}

\begin{proposition}[The main summand at a spherical coefficient prime]
\leavevmode\label{prop:mainselmer724}
Under the same prime and weight assumptions, let \(T\) in
Proposition~\ref{prop:blockrows719} be a lattice in a direct summand of
one of the normalized representations. Suppose \(T\) is unramified at
the auxiliary pairs. Let \(S\) contain all old places of ramification
away from \(\ell\). Then \(\mathcal V_T\) injects into the subgroup of
\(H^1(E,(T\otimes L)/T)\) unramified outside \(S\cup\{w\mid\ell\}\)
and lying at each \(w\mid\ell\) in the image of
\(H^1_f(E_w,T\otimes L)\). In particular, its length \(\lambda_T\) is
bounded by the length of that Selmer group whenever the latter is finite.
\end{proposition}
\begin{proof}
A continuous cocycle into a discrete torsion module has finite image by
compactness, so it takes values in \(\ell^{-\nu}T/T\) for some \(\nu\).
Multiplication by \(\ell^\nu\) identifies this with \(T/\ell^\nu T\).
Its local matrix relations are preserved by this identification. At
\(w\mid\ell\) the rational summand is crystalline with jumps in the same
interval, since subobjects of crystalline representations are crystalline
and the de Rham filtration is exact. Lift the resulting local cocycle
by Theorem~\ref{thm:fllift724} and divide by \(\ell^\nu\). This gives a
class in \(H^1_f(E_w,T\otimes L)\) with the prescribed image.

Away from \(S\), the auxiliary pairs and \(\ell\), all matrices are
trivial on inertia. The corresponding single evaluation equation
gives zero for the cocycle on inertia. At an auxiliary place the
target is unramified, so the restriction to inertia is a homomorphism
to an \(\ell\)-group. It kills wild inertia and factors through the
pro-\(\ell\) tame quotient. The chosen tame generator has value zero
by the definition of \(\mathcal V_T\). Finally a coboundary coming
from \(b\in(T\otimes L)/T\) takes value \((y-1)b\) at the element
\(\eta\) in Section~\ref{sec:relationmatrices719}. The same definition
requires this value to vanish and \(y-1\) is a unit. Thus \(b=0\),
which proves injectivity.
\end{proof}

The local conclusion in Corollary~\ref{cor:scalarfl724} is uniform
in the precision. For primes below the weight threshold, or ramified
in the fields, the torsion comparison used here gives no such bound.
A bound for the global Selmer length of the main summand and the
integral comparison for the projected periods on these same intervals
are also still required by Section~\ref{sec:missing}.
\clearpage

\subsection{One anticyclotomic twist for simultaneous integral tests}
\label{sec:rayinterpolation726}

At a new twist prime, the local contribution of a main summand depends
on the valuation of a root of unity minus one. The order of this root
also affects the values of integral determinants. We construct a single
character on the period field for which all its prescribed powers can
be used in finitely many such polynomial tests. The order needed for
the construction can be chosen before the polynomials are given.
The auxiliary characters and the use of cyclotomic radii follow the
arguments in the supplied archive, in \texttt{AuxiliaryPrimes.lean}
and \texttt{RootsOfUnity.lean} in its Selmer directory.
Throughout this subsection \(v(\ell)=1\) on \(\overline{\QQ}_\ell\).

\begin{lemma}[Independent ray characters on the period field]\leavevmode
\label{lem:raycoordinates726}
Let \(E\) be one of the CM fields of Corollary~\ref{cor:fieldrelations712},
let \(c\) be complex conjugation, and let \(S\) be any finite set of
finite places containing its conjugates and all places above \(\ell\).
Fix \(d,N\ge1\). There exist \(d\) places \(w_i\) of distinct rational
prime norms outside \(S\), all different from their conjugates, and
continuous additive characters
\(\lambda_i:G_E\to\ZZ/\ell^N\ZZ\). They have the following properties.
Each satisfies \(\lambda_i^c=-\lambda_i\), is trivial on the decomposition
groups in \(S\), and is unramified outside \(w_i,w_i^c\). At \(w_i\)
its value on a suitable tame generator is \(1\). At \(w_i^c\) the
conjugate generator has value \(-1\). It is trivial on the decomposition
groups at \(w_j,w_j^c\) for \(j\ne i\). Arithmetic Frobenius lifts at
the \(w_i\) can be chosen on which all the characters vanish.
\end{lemma}
\begin{proof}
Enlarge \(S\) to a finite conjugation-stable set \(T\) whose prime ideal
classes generate the class group of \(E\). This is possible by factoring
ideals in finitely many generating classes. The ring \(\mathcal O_{E,T}\)
then has trivial class group. The group of its units is finitely generated,
including its torsion. Form the extension of \(E\) by \(\mu_{\ell^N}\)
and the \(\ell^N\)-th roots of finitely many generators of these units.
It is a finite Galois extension. Enlarge it by the fields of the ray
characters already constructed, taking their conjugates as well.
Chebotarev supplies completely splitting primes \(w\) of \(E\) outside
all the places used so far. One may choose rational prime norm since
the higher residue degree prime ideals of a number field have Dirichlet
density zero. Avoid all rational primes ramified in \(E\) and all norms
previously chosen. An inert place over \(E^+\) would have square norm,
so \(w\ne w^c\). Write \(q=Nw\). Then \(\ell^N\mid q-1\), and every
\(T\)-unit reduces to an \(\ell^N\)-th power at \(w\).
~\cite{II-II-II-chebotarev713}

Choose a surjection from the cyclic residue group \(k_w^\times\) onto
\(\ZZ/\ell^N\ZZ\). Via reduction at \(w\) this gives a character on
\[
        D_T=\prod_{v\in T} E_v^\times
                    \ \times\!\prod_{v\notin T}\mathcal O_v^\times
                         \subset\mathbb A_{E,f}^\times .
\]
It kills \(E^\times\cap D_T=\mathcal O_{E,T}^\times\). The equality
\(\mathbb A_{E,f}^\times=E^\times D_T\) follows by taking the ideal of
an idele outside \(T\) and using the triviality of the \(T\)-class group.
The local character therefore extends to the finite ideles, kills the
principal ideles, and is continuous because \(D_T\) is open. Extend it
trivially at infinity. By global reciprocity it defines a Galois
character \(\xi\). Put \(\lambda=\xi-\xi^c\). The character \(\xi^c\)
is unramified at \(w\), so \(\lambda\) is still onto on its tame
inertia. The local components at \(T\) are trivial, and at all other
places except \(w,w^c\) they are unramified. Before repeating the
construction put \(w,w^c\) in \(T\). Earlier characters vanish at the
new decomposition groups by the splitting imposed on the new primes,
and the new character vanishes at the earlier groups by the choice of
\(T\). At \(w_i\) any Frobenius lift may be multiplied by a power of
the chosen tame generator to make its value under \(\lambda_i\) zero.
All other \(\lambda_j\) vanish on that decomposition group.
\end{proof}

\begin{lemma}[A uniform range of cyclotomic orders]\leavevmode
\label{lem:gausscircle726}
Fix a positive real radius in valuation coordinates \(\rho\), a number
\(B\ge0\), \(0<\delta<1\), and a lower exponent \(h_*\ge1\) with
\(r_{h_*}\le\rho\), where \(r_h=1/\varphi(\ell^h)\). Choose \(N\ge h_*\)
large enough that
\[
                   \delta\sum_{h=h_*}^{N}\frac{1}{h+1}> B .
\]
For every nonzero polynomial \(g(X)=\sum b_i X^i\) with integral
algebraic \(\ell\)-adic coefficients and
\(v_\rho(g)=\min_i(v(b_i)+i\rho)\le B\), there is a primitive
\(\ell^h\)-th root \(\zeta\) with \(h_*\le h\le N\) such that
\[
                   v(g(\zeta-1))\le v_{r_h}(g)+\delta .
\]
\end{lemma}
\begin{proof}
Factor \(g\) over \(\overline{\QQ}_\ell\) and count roots with multiplicity,
putting \(v(0)=+\infty\). For \(u\ge0\) the Gauss valuation satisfies
\[
 v_u(g)=v_0(g)+\sum_{\alpha:\,v(\alpha)>0}\min(v(\alpha),u).
\]
Indeed it is multiplicative. For a rational \(u\), substitute a scalar
of valuation \(u\) times \(X\), divide each factor by a coefficient
of least valuation, and reduce modulo the maximal ideal. The product
of nonzero reduced polynomials is nonzero. This proves multiplicativity
for rational \(u\), and continuity of the finite minima gives it for
real \(u\). Applying this to each linear factor proves the formula,
also for polynomials with nonintegral roots.

Let \(R_h\) be the multiset of roots with \(v(\alpha)=r_h\).
The \(r_h\) are positive and distinct. Since \(v_0(g)\ge0\),
\(\sum_{h=h_*}^{N}|R_h|r_h\le v_\rho(g)\), so there exists an \(h\)
for which \(|R_h|r_h(h+1)<\delta\). For \(\alpha\in R_h\) the
polynomial \(F_\alpha(Y)=(1+\alpha+Y)^{\ell^h}-1\) is integral and
monic. Its roots are \(\xi-1-\alpha\) for all \(\ell^h\)-th roots
\(\xi\). They are integral. The coefficient of \(Y\) has valuation
\(h\), as \(1+\alpha\) is a unit. The same Gauss formula at \(u=1\)
shows that
\[
   \sum_{\xi\ {\rm primitive\ of\ order}\ \ell^h}
           \min\bigl(v(\xi-1-\alpha),1\bigr)
                         \le v_1(F_\alpha)\le h+1.
\]
There are \(1/r_h\) primitive roots. Averaging this inequality over
them and summing over \(R_h\) gives one \(\zeta\) for which the sum of
the truncated valuations over \(R_h\) is less than \(\delta<1\).
Every term consequently equals its untruncated valuation.
The identity \(v(\zeta-1)=r_h\) follows from the cyclotomic polynomial
at \(1\), since all the primitive roots are conjugate and their
differences from \(1\) have product of absolute valuation \(1\).
For a root outside \(R_h\), the valuation of \(\zeta-1-\alpha\)
equals \(\min(r_h,v(\alpha))\) by the ultrametric inequality.
Subtracting the Gauss valuation of the product now gives
\[
   v(g(\zeta-1))-v_{r_h}(g)
       =\sum_{\alpha\in R_h}\bigl(v(\zeta-1-\alpha)-r_h\bigr)<\delta .
\]
This includes a constant \(g\), for which the sum is empty.
\end{proof}

\begin{theorem}[One choice of twist for all powers and polynomials]
\leavevmode\label{thm:allraypowers726}
Fix \(\ell,E,S,d\) as in Lemma~\ref{lem:raycoordinates726}, an integer
\(m\ge1\), and real numbers \(\rho,\varepsilon,\eta>0\) and \(B\ge0\).
One can choose \(N,w_i,\lambda_i\) before the following tests are
specified. A test consists of an integral polynomial \(P_a\) in \(d\)
variables, an integer \(s_a\) with \(1\le |s_a|\le m\), and a witness
\(x_a\in\overline{\QQ}_\ell^d\) with \(v(x_{a,i})\ge\rho\).
Allow any finite list with \(\sum_a v(P_a(x_a))\le B\).
There is one finite-order anticyclotomic character \(\vartheta\) trivial
on \(S\), with tame conductor supported on the selected pairs. Write
\(\gamma_i\) for the selected tame generator at \(w_i\) and \(\zeta_i
=\vartheta(\gamma_i)\). Then
\[
 v\bigl(P_a((\zeta_i^{s_a}-1)_i)\bigr)
                      \le v(P_a(x_a))+\varepsilon
                       \quad\text{for every }a.
\]
The same estimate holds for the sum of these valuations, with
\(\varepsilon\) added only once. For every \(1\le |s|\le m\),
\(\vartheta^s\) has nontrivial inertia at all \(2d\) selected places and
\[
             \sum_{w\in\bigcup_{i=1}^d\{w_i,w_i^c\}} v(\vartheta^s(\gamma_w)-1)
                                \le \eta ,
\]
where the sum ranges over the \(2d\) distinct places.
\end{theorem}
\begin{proof}
Put \(L=\lfloor\log_\ell m\rfloor\), \(\rho'=\rho/m\), and choose
\(h_*>L\) with \(r_{h_*}\le\rho'\) and \(2d\ell^Lr_{h_*}\le\eta\).
Take \(\delta=\min(1,\varepsilon)/(d+1)\) and apply
Lemma~\ref{lem:gausscircle726} with \(B+1,\rho',h_*,\delta\) to
choose \(N\) using divergence of the harmonic series. Now construct the ray
characters of this order. Up to this point the polynomial degrees,
their number and their witnesses have played no part.

For a test of power \(s>0\), substitute \((1+X_i)^s-1\) in each
coordinate. For \(s=-b<0\) use \((1+X_i)^{-b}-1\) and multiply by
\(\prod_i(1+X_i)^{b\deg_{X_i}P_a}\). This gives, in both cases, an
integral polynomial \(Q_a\). The extra factors are units whenever
\(1+X_i\) is a unit. There are witnesses \(y_a\) for these polynomials
with \(v(y_{a,i})\ge\rho'\) and \(v(Q_a(y_a))=v(P_a(x_a))\).
For \(b=|s|\), solve \((1+Y)^b=1+x\) if \(s>0\), and
\((1+Y)^b=(1+x)^{-1}\) if \(s<0\). All roots in \(Y\) are integral
because the equation is monic integral. Their product has valuation
\(v(x)\). Thus some root has valuation at least \(v(x)/b\). If \(x=0\)
take \(Y=0\).

Replace the witness coordinates successively by common values. At
coordinate \(i\), let \(g_a(X)\) be \(Q_a\) specialized to the current
values of its other variables. All coefficients are integral and the
unreplaced coordinate of each witness still has valuation at least
\(\rho'\). For \(g=\prod_a g_a\), multiplicativity gives
\(v_{\rho'}(g)\le\sum_a v(Q_a(\text{current witness}))\).
Inductively this last sum is at most \(B+(i-1)\delta<B+1\).
Choose a primitive root of an allowed order by the preceding lemma.
For each factor \(v(g_a(\zeta_i-1))-v_{r_h}(g_a)\ge0\) by the
ultrametric inequality, and the sum of these excesses is at most
\(\delta\). Since \(v_{r_h}(g_a)\le v_{\rho'}(g_a)\) and the latter
is at most its previous witness valuation, each increase and the
increase of the sum are bounded by \(\delta\). After \(d\) coordinates
the total allowance is \(d\delta\le\varepsilon\).

Finally take \(\vartheta(g)=\prod_i\zeta_i^{\lambda_i(g)}\), well
defined since each order divides \(\ell^N\). At a chosen tame generator
its value is \(\zeta_i\), or \(\zeta_i^{-1}\) at the conjugate. For
\(0<|s|\le m\) its value there has order \(\ell^{h_i-v_\ell(s)}>1\).
The valuation of its difference from \(1\) is therefore
\(\ell^{v_\ell(s)}r_{h_i}\). Summing proves the ramification bound.
\end{proof}

\begin{corollary}[Selmer conditions at the new twist places]\leavevmode
\label{cor:raytamagawa726}
Suppose \(T_1,\ldots,T_b\) are fixed Galois lattices on \(E\).
Include all their ramified places outside \(\ell\) in \(S\).
For every \(T_j\) of rank \(r_j\), and every \(0<|s|\le m\), imposing
the finite local conditions at the new pairs on the divisible
coefficient Selmer group for \(T_j\otimes\vartheta^s\) costs at most
\(r_j\eta\) in normalized length. The character may satisfy the
polynomial conclusions of Theorem~\ref{thm:allraypowers726} at the same
time.
\end{corollary}
\begin{proof}
Extend coefficients to contain the character values. At each new place
inertia acts by the single nontrivial scalar
\(u=\vartheta^s(\gamma_w)\). Proposition~\ref{prop:maintamagawa725}
applies with determinant \((1-u)^{r_j}\). Thus the length divided by the
absolute ramification index of the coefficient ring is at most
\(r_j v(1-u)\) locally. The quotient of the group with relaxed conditions
by the group with finite conditions injects into the direct sum of the
local quotients. The sum in Theorem~\ref{thm:allraypowers726} proves the
claim. Triviality on \(S\) ensures that all local representations at
the coefficient places are unchanged by this twist.
\end{proof}

This result controls the local cost of adding twist primes on a fixed
period field. To bound the global Selmer length by specializing minors,
one needs witnesses with a quantitative bound \(B\) for the actual
relation matrices, and compatible witnesses for the ratios of the
Hecke values entering the periods.
\clearpage

\subsection{Integral minors for the Selmer groups of the twists}
\label{sec:selmerpolys727}

We describe the polynomials needed for the main summand at a fixed
precision. In one case a witness can be constructed at the untwisted
point, with no enlargement of the group from relaxing ramification.
This gives a bound for one twist and all its powers in terms of a
group at the original level.

Fix a CM period field \(E\), a coefficient ring \(\mathcal O\) with
fraction field \(L\), a finite conjugation-stable set \(S\) containing
the coefficient places, and a rank-\(r\) lattice \(T\) unramified
outside \(S\). Specify linear evaluation equations on cocycles, with
coefficients in \(\mathcal O\). Each equation involves a finite sum
of coordinates of values at specified elements of \(G_E\). For the
application include \(z(\eta)=0\) at the marked element of
Section~\ref{sec:relationmatrices719} and all equations satisfied by
local integral crystalline cocycles at the coefficient places.
The list of equations itself may be infinite. Extend the coefficients
of the equations when extending scalars. We use normalized length,
ordinary length divided by the absolute ramification index of the
coefficient ring, and the valuation \(v(\ell)=1\).

\begin{proposition}[A finite set of sharp integral minors]\leavevmode
\label{prop:sharpminors727}
Let \(W\) be a finite set of additional places disjoint from \(S\),
and let \(\lambda_i:G_E\to\ZZ/\ell^N\ZZ\), \(1\le i\le d\), be
characters unramified outside \(S\cup W\). Fix \(M\ge1\).
There is a finite matrix \(D(X)\) over
\(\mathcal O[X_1,\ldots,X_d]\), of full column rank after every
specialization in characteristic zero, with the following property.
For any extension of coefficient rings \(\mathcal O'/\mathcal O\) and
any \(\ell^N\)-th roots \(\zeta_i\in\mathcal O'\), let
\(\chi=\prod_i\zeta_i^{\lambda_i}\). The minimum of the valuations
of the maximal column-size minors of \(D((\zeta_i-1)_i)\) equals the
normalized length of the cocycles into \((T\otimes\mathcal O'/\ell^M)(\chi)\)
unramified outside \(S\cup W\) and satisfying the given equations.
Here and below the tensor with \(\mathcal O'/\ell^M\) is over \(\mathcal O\).
\end{proposition}
\begin{proof}
Let \(G\) be the Galois group of the maximal extension unramified outside
\(S\cup W\), and \(H\) the common kernel of \(T\bmod\ell^M\) and the
\(\lambda_i\) in \(G\). It is open and normal. Put
\[
              \Gamma=G/\overline{[H,H]H^{\ell^M}}.
\]
This is a finite group. Indeed the fixed field \(K\) of \(H\) is a
number field. The abelian quotient of \(H\) of exponent \(\ell^M\) is
unramified outside the finite set over \(S\cup W\). By local reciprocity
its characters there kill \((K_v^\times)^{\ell^M}\), which contains
an open subgroup of the local units. At \(v\nmid\ell\) this follows
on sufficiently deep principal units by Hensel's lemma, and at
\(v\mid\ell\) by the logarithm on deep units. Choose a modulus with
these local exponents. Global class field theory makes the abelian
quotient a quotient of the finite ray class group of that modulus.
The possible residue characters at infinity are trivial since \(K\)
is totally imaginary. This proves finiteness of \(\Gamma\).

The action of each of the indicated twists modulo \(\ell^M\) is trivial
on \(H\). A cocycle restricts on \(H\) to a continuous homomorphism
killing commutators and \(\ell^M\)-th powers, and hence their closure.
The cocycle consequently factors through \(\Gamma\). Conversely
inflation from \(\Gamma\) gives all these cocycles. Take \(r\) variables
for the value \(Z_g\) at each \(g\in\Gamma\). Lift the entries of the
fixed matrices \(T(g)\bmod\ell^M\) to \(\mathcal O\), writing
\(\widetilde T(g)\) for a choice of lift. Let \([b]\in[0,\ell^N-1]\)
be the integer representative of \(b\in\ZZ/\ell^N\ZZ\). Use the rows
of the polynomial linear equations
\[
 Z_{gh}-Z_g-\widetilde T(g)
                 \prod_i(1+X_i)^{[\lambda_i(g)]}\,Z_h=0
                       \qquad(g,h\in\Gamma).
\]
At the specified roots of unity these are exactly the cocycle equations
over \(\mathcal O'/\ell^M\). Evaluation equations map to constant rows
by projecting their indexing group elements to \(\Gamma\). The span of
these rows in the finite free \(\mathcal O\)-module of all rows is
finitely generated. Include any finite generating list. This constructs
a matrix \(A(X)\) with \(b=r|\Gamma|\) columns. Adjoin the rows of
\(\ell^M I_b\) to form \(D(X)\).

At a specialization the desired module is the kernel of \(A\) on
\((\mathcal O'/\ell^M)^b\). By diagonalizing \(A\) over the valuation
ring its normalized length is \(\sum_{j=1}^b\min(v(a_j),M)\) for the
diagonal entries, filling with zeros when necessary and putting
\(v(0)=+\infty\). This is also the length of the quotient of the
integral row lattice by the row conditions, that is, of
\((\mathcal O')^b/\operatorname{rowspan}D\). Its zeroth Fitting ideal
is generated by the column-size minors of \(D\), giving the claimed
minimum. The adjoined rows alone have nonzero determinant.
\end{proof}

\begin{lemma}[Detection at a single precision]\leavevmode
\label{lem:precisionlength727}
For any fixed twist over \(\mathcal O'\), consider cocycles with values
in \((T\otimes L')/(T\otimes\mathcal O')\), where \(L'\) is the fraction
field of \(\mathcal O'\), subject to the same equations and outside
ramification condition. If the normalized length at precision
\(\ell^M\) in Proposition~\ref{prop:sharpminors727} is strictly less
than \(M\), it equals the normalized length with these divisible
coefficients.
\end{lemma}
\begin{proof}
The subgroup killed by \(\ell^M\) is identified with the finite
precision cocycles by multiplying their values by \(\ell^M\) on
\(\ell^{-M}(T\otimes\mathcal O')/(T\otimes\mathcal O')\).
Every individual cocycle has bounded order, by compactness of the
Galois group and discreteness of the target. If the module contained
an element not killed by \(\ell^M\), the cyclic submodule it generates
would have an \(\ell^M\)-torsion subgroup of normalized length \(M\),
by the ideal classification in a discrete valuation ring. This would
contradict the strict inequality.
\end{proof}

\begin{theorem}[An auxiliary pair preserving the original Selmer bound]
\leavevmode\label{thm:preserveselmer727}
Assume \(T\) has a marked element \(\eta\) in inertia at a place
of \(S\) away from \(\ell\), acting by a scalar \(y\) whose reduction
is different from \(1\). Assume also that
\(T^c\otimes L\simeq T^\vee\otimes L(\kappa^a)\) for some integer \(a\).
Fix \(M,m\ge1\) and \(\varepsilon,\eta_0>0\).
Let \(b_M\) be the normalized length at precision \(\ell^M\) for the
untwisted lattice over \(\mathcal O\), with ramification permitted
only in \(S\) and the chosen evaluation equations. There exist a single
pair \(w,w^c\) and a character \(\vartheta\) of the type constructed in
Theorem~\ref{thm:allraypowers726}, trivial on all decomposition groups
in \(S\), for which the following hold for every \(0<|s|\le m\).
\begin{enumerate}
\item The normalized length of the group of cocycles for
\(T\otimes\vartheta^s\), with ramification permitted at the pair and
subject to the same equations, is at most \(b_M+\varepsilon\) at
precision \(\ell^M\). If \(b_M+\varepsilon<M\), the same bound holds
with divisible coefficients.
\item The two normalized local costs of imposing the finite conditions
at \(w,w^c\) have sum at most \(r\eta_0\).
\end{enumerate}
If the untwisted group with the specified equations modulo the maximal
ideal of \(\mathcal O\) is zero, the corresponding groups with divisible
coefficients are zero for every \(\ell\)-power twist of the form above
on the selected pair, for all its powers.
\end{theorem}
\begin{proof}
Choose the order bound \(\ell^N\) using the proof of
Theorem~\ref{thm:allraypowers726} with \(d=1\), any fixed
\(\rho>0\), and \(B=2m b_M\). The untwisted group is finite by the
finite quotient argument above with \(W\) empty. We refine the
prime choice in Lemma~\ref{lem:raycoordinates726}. Write \(K_0\)
for the splitting field over \(E\) of the roots of the selected
\(S'\)-units of that proof, with \(\mu_{\ell^N}\) included. Here \(S'\)
denotes its enlarged generating set of places. Since \(\eta\) is in
inertia away from \(\ell\), it fixes \(\mu_{\ell^N}\). On each chosen
radical it acts by multiplication by an \(\ell^N\)-th root of unity.
Thus \(g=\eta^{\ell^N}\) fixes \(K_0\). Apply Chebotarev in the
compositum of \(K_0\) with the field of the residual representation
of \(T\). Choose a prime of rational prime norm with Frobenius the
conjugacy class of \(g\), outside \(S'\) and outside the finite
ramification of this compositum. The new prime splits in \(K_0\),
as required for the ray character. Its residue cardinality \(q\) is
\(1\bmod\ell^N\). The residual Frobenius of \(T\) there is the scalar
\(\bar y^{\ell^N}\ne1\), since the multiplicative group of the
coefficient residue field has order prime to \(\ell\). At the conjugate
place the polarization identifies its eigenvalues with the inverses
times \(q^a\), so they too differ from \(1\) in the residue field.
The equality of characteristic polynomials used here holds over \(L\)
and hence holds on reduction for any invariant lattice. In particular
\(T(\phi_v)-qI\) is invertible over \(\mathcal O\) at both new places.

For the untwisted lattice, allowing ramification at these places does
not add any cocycles. Indeed inertia acts trivially on the target of
an untwisted cocycle there. Its restriction factors through pro-\(\ell\)
tame inertia, and at a tame generator the Frobenius identity is
\((T(\phi_v)-qI)z(\gamma_v)=0\). Invertibility implies that this
value is zero at every precision, also after any coefficient extension.
At \(X=0\), the matrix of Proposition~\ref{prop:sharpminors727} on the
enlarged set of places thus has minimum minor valuation exactly
\(b_M\). Select one minor \(P\) attaining that minimum. For each
of the \(2m\) signed powers use this same polynomial and its witness
\(x=0\). The polynomial argument of Theorem~\ref{thm:allraypowers726}
uses only the prescribed order bound and the independent ray
character at the chosen pair. It applies to the refined choice of
prime and gives \(v(P(\zeta^s-1))\le b_M+\varepsilon\) simultaneously.
The actual length is at most that valuation. Lemma~\ref{lem:precisionlength727}
and Corollary~\ref{cor:raytamagawa726} prove the first two assertions.

For the last assertion the residual normalized cocycle space remains
zero at the enlarged set of places by the same tame calculation.
Every finite \(\ell\)-power root of unity reduces to \(1\).
Thus after any such twist the residual equations are exactly the
untwisted residual equations. Extending the residue field preserves
their zero kernel, by linear algebra on the finite quotient, or by
taking coordinates in a residue field basis. A nonzero cocycle with
divisible coefficients has finite order. If the smallest power of a
uniformizer killing it is the \(j\)-th power, multiplying by the
\((j-1)\)-st power gives a nonzero cocycle killed by the uniformizer.
That multiple would
give a nonzero solution of the residual equations. This proves
vanishing.
\end{proof}

For a main summand with the squarefree marked type and polarization
of Section~\ref{sec:markedmatrix715}, the scalar hypothesis holds
for every \(\ell\) at one of its two original level primes. At coefficient
primes covered by Corollary~\ref{cor:scalarfl724}, every member of
\(\mathcal V_T\) for the corresponding twisted family satisfies the
local crystalline equations, provided the actual family has the
hyperspecial coefficient level and the common weight interval of that
corollary. This follows from Theorem~\ref{thm:fllift724} at each
coefficient place, with the twist trivial there. Taking the equations
at \(\eta\) and at the coefficient places thus gives a finite minor
bound for these cocycles whenever the untwisted test gives
\(b_M+\varepsilon<M\). For a large untwisted group the construction
above preserves its bound. A quantitative reduction of such groups
under twisting is still needed for the common interval estimate.
\clearpage

\subsection{The primary and dual terms for a polarized marked summand}
\label{sec:polarizedlength728}

Let \(T\) be a lattice over a discrete valuation ring \(\mathcal O\)
finite over \(\ZZ_\ell\), with uniformizer \(\pi\), fraction field \(L\)
and \(v_\pi(\pi)=1\). Write \(V=T\otimes L\), \(G=G_E\),
\(r=\operatorname{rank}T\), \(A_a=\mathcal O/\pi^a\), and
\(W_a=T/\pi^aT\). We use the \(\mathcal O\)-linear Tate dual
\(W_a^*=\operatorname{Hom}_{\mathcal O}(W_a,L/\mathcal O)(1)\).
It identifies with the usual character dual with the Tate twist
after choosing a generator of the inverse different of
\(\mathcal O/\ZZ_\ell\) and using the trace to
\(\QQ_\ell/\ZZ_\ell\). Lengths and local orthogonal complements in the
statements below are independent of this choice.

\begin{theorem}[The global congruence term at every precision]\leavevmode
\label{thm:selftwist728}
Suppose the \(\mathcal O\)-span of \(T(G)\) is \(M_r(\mathcal O)\)
and \(V^c\simeq V^\vee(1-t)\), with \(t\ge1\). Let \(I_t\subset
\mathcal O\) be the ideal generated by
\[
                  (1-\kappa(g)^t)\operatorname{tr}T(g),
                                \qquad g\in G .
\]
Then \(\operatorname{Hom}_{A_a[G]}(W_a,(W_a^*)^c)\) is cyclic and has
ordinary \(\mathcal O\)-length \(\min(a,v_\pi(I_t))\), where the
valuation of the zero ideal is \(+\infty\). The ideal may be computed
using good Frobenius elements on any set of primes of Dirichlet density
one. It is unchanged by every finite-order anticyclotomic twist after
extension to a ring containing the character values, with the original
ideal extended to that ring.
\end{theorem}
\begin{proof}
Every \(G\)-stable full lattice in \(V\) is homothetic to \(T\).
Indeed it is stable under every integral matrix, since those matrices
are spanned by the action of \(G\). The coordinate idempotents split
it into fractional ideals on the coordinate axes, and the off-diagonal
matrix units identify those ideals. The same argument applies after a
scalar character twist by units. The rational polarization identifies
\((V^\vee)^c(1)\) with \(V(t)\), so the corresponding integral lattices
can be identified equivariantly after a scalar rescaling. With a choice
of \(\pi^{-a}\) for the finite coefficient dual this identifies the
homomorphism group in the theorem with
\(\operatorname{Hom}_{A_a[G]}(W_a,W_a(t))\).

The residual representation is absolutely irreducible by the full matrix
span. Homomorphisms between it and its residual twist have residue
dimension at most one, by Schur's lemma. In the finite-level
homomorphism group the subgroup killed by \(\pi\) is exactly this
residual homomorphism space, by dividing all entries by \(\pi^{a-1}\).
The structure theorem over a discrete valuation ring therefore makes
the group cyclic. A homomorphism of order ideal \((\pi^s)\), \(s>0\),
has entries divisible by \(\pi^{a-s}\). After that division its
reduction to precision \(s\) has at least one unit entry. Its residual
map between the irreducible modules is nonzero, so it is invertible
modulo \(\pi\), hence modulo \(\pi^s\). Conversely an isomorphism at
precision \(s\le a\), multiplied by \(\pi^{a-s}\), supplies a
homomorphism of precisely that order at precision \(a\).

For completeness, at precision \(s\) equality of traces at every
\(g\) suffices for an isomorphism. Write \(B=\mathcal O/\pi^s\),
\(u=T\bmod\pi^s\) and \(u'=\kappa^t u\). Both maps \(B[G]\to M_r(B)\)
are onto, since twisting multiplies each spanning matrix by a unit.
The trace pairing on \(M_r(B)\) is perfect, as one sees from the matrix
units. Equality of traces on \(G\) gives equality on \(B[G]\), hence
equality of the two kernels by testing each element against every
right multiplier. We obtain a \(B\)-algebra automorphism of \(M_r(B)\)
carrying \(u(g)\) to \(u'(g)\). To see that it is inner, let \(F_{ij}\)
be the images of the standard matrix units. The projections \(F_{ii}\)
split \(B^r\) into \(r\) isomorphic projective summands, each free of
rank one since \(B\) is local. Choose a basis of the first summand and
transport it to the others by \(F_{i1}\). In this basis all \(F_{ij}\)
are the standard units. Thus the two actions are isomorphic.
Trace equality is precisely \(I_t\subset(\pi^s)\). This proves the
length formula. Reduction modulo any \(\pi^s\) makes the trace test
a finite union of Frobenius conjugacy conditions, proving the
density-one assertion by Chebotarev. A character twist multiplies
each displayed generator of \(I_t\) by a unit. Anticyclotomicity
preserves the polarization with the same \(t\).
\end{proof}

\begin{corollary}[Uniform size on the two marked level primes]\leavevmode
\label{cor:congcost728}
Suppose the main summand is as in Theorem~\ref{thm:markedmatrix715}
with \(1\le r\le r_0\), \(t\le t_0\), on the shorter period field.
For each \(\ell\) pick \(p_i\ne\ell\) with a marked place of the
prescribed primitive type, or its indicated rank-one type.
Then the supremum over \(a\) of the normalized length in
Theorem~\ref{thm:selftwist728} is at most \(v_\ell(p_i^{rt}-1)\).
For \(G_0=\prod_i |\mathrm{GL}_{r_0t_0}(\mathbb F_{p_i})|^f\) of
Corollary~\ref{cor:actualmarks714}, the sum of these suprema with
weights \(\log\ell\) is at most \((\log G_0)/f\).
\end{corollary}
\begin{proof}
By the theorem on the marked local matrices, \(g=\phi_w^r\) acts by
a unit scalar and \(\kappa(g)=p_i^r\). Any isomorphism modulo
\(\pi^s\) from \(T\) to \(T(t)\) intertwines that scalar with
\(p_i^{rt}\) times itself, implying \(p_i^{rt}\equiv1\bmod\pi^s\).
Use the isomorphism criterion in the proof of the preceding theorem.
The integer \(p_i^{rt}-1\) divides
\(\lvert\mathrm{GL}_{r_0t_0}(\mathbb F_{p_i})\rvert\), since
\(rt\le r_0t_0\). Thus even when the choice of \(i\) depends on
\(\ell\), the bound follows by taking valuations of the product of
the two orders. The same scalar argument applies to any unramified
twist at the chosen place, and Theorem~\ref{thm:selftwist728} covers
arbitrary anticyclotomic twists.
\end{proof}

\begin{lemma}[Integral equations for a Bloch--Kato local condition]
\leavevmode\label{lem:localevaluation728}
Let \(K/\QQ_\ell\) be a finite extension at which \(V\) is de Rham.
Let \(Z_f\) be the integral cocycles \(G_K\to T\) whose rational
class lies in \(H^1_f(K,V)\). At precision \(\pi^a\), the solutions to
all coordinate evaluation relations on \(Z_f\) are exactly the
reductions of \(Z_f\). Write \(\mathcal L_K\) for their image in
\(H^1(K,W_a)\). Then
\[
   \operatorname{length}_{\mathcal O}\mathcal L_K
       -\operatorname{length}_{\mathcal O}H^0(K,W_a)
        = a\,\dim_L\bigl(D_{\rm dR}(V)/\mathrm{Fil}^0\bigr).
\]
\end{lemma}
\begin{proof}
The rational space of finite-condition raw cocycles has dimension
\(r+\dim_L(D_{\rm dR}(V)/\mathrm{Fil}^0)\). The Bloch--Kato local
fundamental exact sequence gives
\(\dim H^1_f=\dim H^0+\dim D_{\rm dR}/\mathrm{Fil}^0\);
the coboundaries have dimension \(r-\dim H^0\).
~\cite[Cor.~3.8.4]{II-II-II-bktate717}
Integral-valued cocycles form a full lattice in this space. In fact
rational cocycles have bounded denominators by compactness, and a
finite set of coordinate evaluations bounds the lattice from above, since coordinate evaluations separate the rational vectors and span
their dual. Moreover the reductions of all evaluation functionals
separate \(Z_f/\pi Z_f\). A cocycle all of whose values are divisible
by \(\pi\) can itself be divided by \(\pi\) in \(Z_f\).
By finite-dimensional duality followed by Nakayama's lemma the
evaluations therefore span \(Z_f^\vee\) integrally. Assigning coordinate
values in \(A_a\) and obeying every integral relation is consequently
exactly the same as assigning a linear map \(Z_f^\vee\to A_a\),
which is an element of \(Z_f/\pi^aZ_f\). This proves the reduction
assertion. Every raw coboundary at precision \(a\) lifts integrally.
Their length is \(ar-\operatorname{length}H^0(K,W_a)\).
Subtract it from the length of \(Z_f/\pi^aZ_f\) to obtain the formula.
\end{proof}

\begin{theorem}[Comparison with the orthogonal Selmer group]\leavevmode
\label{thm:balancedcount728}
Assume the hypotheses of Theorem~\ref{thm:selftwist728}, that \(V\) is
de Rham above \(\ell\), and that an inertia element away from \(\ell\)
acts by \(yI_r\) with \(y-1\) a unit. At the coefficient places use
the conditions of Lemma~\ref{lem:localevaluation728}. At a finite set
\(S_0\) away from \(\ell\), including all other ramified places, allow
arbitrary local classes. At any additional unramified places \(v\)
where \(W_a\) and \(\mu_{\ell^k}\) are trivial as local modules for
some \(\ell^k\) killing \(A_a\), choose a normalized tame generator
\(\tau_v\), an arithmetic Frobenius lift \(\phi_v\), and impose
\(z(\tau_v)=s_v z(\phi_v)\), \(s_v\in A_a\). Use the unramified
condition elsewhere. Denote these conditions by \(\mathcal L\) and
their local Tate annihilators on \(W_a^*\) by \(\mathcal L^\perp\).
Then
\[
  \operatorname{length} H^1_{\mathcal L}(E,W_a)
  \le \operatorname{length} H^1_{\mathcal L^\perp}(E,W_a^*)
         +\sum_{w\in S_0}\operatorname{length}H^0(E_w,W_a^*) .
\]
At an original marked place with no eigenvalue \(1\) on full tame
inertia, the summand on the right is at most
\(v_\pi(\det(1-T(\sigma_w)))\). All lengths in this assertion are
ordinary \(\mathcal O\)-lengths.
\end{theorem}
\begin{proof}
The equality of local orders in the Greenberg--Wiles formula, obtained
by taking orders in Poitou--Tate duality, gives the difference of the
two Selmer lengths as
\[
   \operatorname{length}H^0(E,W_a)
   -\operatorname{length}H^0(E,W_a^*)
   +\sum_v\bigl(\operatorname{length}\mathcal L_v
                         -\operatorname{length}H^0(E_v,W_a)\bigr).
\]
At an infinite place \(\mathcal L_v\) has order one. Both global
invariant groups are zero since the marked element acts by scalars
\(y,y^{-1}\) on the two modules. At a finite unramified place with
unramified condition, the local contribution is zero by the kernel
and cokernel formula for Frobenius on a finite module. At an extra
place with the indicated linear constraint, the cocycle is a
homomorphism and its two local coordinates are independently arbitrary in \(W_a\).
Indeed the tame relation is \((q_v-1)z(\tau_v)=0\), automatic
under the hypothesis. The graph of multiplication by \(s_v\) has
length \(ar=\operatorname{length}H^0(E_v,W_a)\).
At the places of \(S_0\), the finite local Euler characteristic
formula gives precisely \(\operatorname{length}H^0(E_w,W_a^*)\).
~\cite{II-II-II-duality13}

Let \(f=[E^+:\QQ]\). After coefficient extension to split all embeddings,
the dimension in Lemma~\ref{lem:localevaluation728} counts negative
de Rham filtration jumps at the embeddings belonging to the place.
Extension of scalars does not change that dimension over the
coefficient field. Under the polarization, a jump \(u\) at an
embedding \(\tau\) is paired with \(t-1-u\) at its conjugate
embedding. Since \(t-1\ge0\), each of the \(fr\) such pairs contributes
at most one negative jump. Thus the total positive contribution
above \(\ell\) is at most \(afr\). The \(f\) complex places each
contribute \(-ar\), giving the inequality. At the specified marked
places, the dual invariant group is contained in the kernel of
\(T(\sigma_w)^{-\mathsf t}-1\) on the finite coefficient dual. Over
\(L\) this matrix is invertible, and its determinant has the same
valuation as \(\det(1-T(\sigma_w))\). Elementary divisors bound its
kernel length as claimed.
\end{proof}

The term in Theorem~\ref{thm:selftwist728} is the size of a space of
equivariant homomorphisms. In Theorem~\ref{thm:balancedcount728} the
remaining term is a dual Selmer group of global cohomology classes.
To use the first bound in the second inequality one needs a proof
that the relevant dual classes inject into that homomorphism space
after choosing the auxiliary local conditions.
\clearpage

\subsection{Reducing the Selmer groups by local conditions at split pairs}
\label{sec:pairedselmer729}

We now construct the local conditions required at the end of
Section~\ref{sec:polarizedlength728}. The construction works for
finitely many polarized lattices at once. It is important that the
number of primes can be bounded before specifying the orders of
the ray characters. We give the argument for the prime selection
and the global duality pairing in detail. The commutator argument
and the descent by lengths are drawn from the supplied archive's
\texttt{Mirror.lean} and \texttt{Descent.lean}.

\begin{lemma}[Varying a linear condition over a finite valuation ring]
\leavevmode\label{lem:slopecount729}
Let \(C\) be a finite module over \(\mathcal O\), with two linear
maps \(u,v:C\to(\mathcal O/\pi^a)^b\). Outside at most \(b\)
residues of \(d\) in the residue field, one has
\[
 \operatorname{length}\ker(d u-v)\le
       \operatorname{length}\bigl(\ker u\cap v^{-1}(u(C))\bigr).
\]
\end{lemma}
\begin{proof}
Let \(C'\) be the sum in \(C\) of all \(\ker(xu-v)\) for \(x\in\mathcal O\).
Then \(v(C')\subset u(C')\). Choose finitely many generators of \(C'\),
and lift the columns of their \(u\)- and \(v\)-values to integral
matrices \(U,V\). The length of the image generated by \(U\) in
\((\mathcal O/\pi^a)^b\) is \(ab\) minus the minimum \(k\) of the
valuations of the \(b\)-rowed minors of \([U\ \pi^a I_b]\), by elementary
divisors. Choose a minor attaining \(k\). The corresponding determinant
of \([XU-V\ \pi^a I_b]\) is a polynomial of degree at most \(b\)
whose leading coefficient is the selected determinant, by multilinearity
in the columns. Its least coefficient valuation is therefore at
most \(k\). Divide by that common uniformizer power and reduce. This
gives a nonzero polynomial of degree at most \(b\). Outside its residue
roots the image length of \(d u-v\) on \(C'\) is at least that of \(u\)
on \(C'\). Subtract from the length of \(C'\), observing that
\(\ker(d u-v)\subset C'\) and that \(\ker(u|_{C'})\subset
\ker u\cap v^{-1}(u(C))\).
\end{proof}

\begin{theorem}[Paired conditions with bounded primary length]\leavevmode
\label{thm:tiedlength729}
Over the same field \(E\), take finitely many lattices \(T_i\) satisfying
Theorem~\ref{thm:balancedcount728}, of ranks \(r_i\) and polarizing
parameters \(t_i\ge1\). For each \(i\) require an element \(\eta_i\)
of inertia away from \(\ell\) with scalar value \(y_i\) such that
\(y_i^2-1\) is a unit. Fix \(S\), closed under conjugation and containing
all ramification of the lattices, the marked places, and the coefficient
places. Use the finite local conditions of Lemma~\ref{lem:localevaluation728}
above \(\ell\), unrestricted conditions at the other places in \(S\),
and unramified conditions elsewhere.

Fix a precision \(a_0\ge1\) in powers of \(\pi\). Enlarge the residue
field by an unramified coefficient extension until its order exceeds
\(1+2a_0\max_i r_i\). There is a bound \(d_0\), depending on these data
and \(a_0\), with the following property for every \(N\) such that
\(\pi^{a_0}\mid\ell^N\) in \(\mathcal O\). One can find at most \(d_0\)
places \(J\) of distinct rational prime norms, outside \(S\) and
disjoint from their conjugates, with a first place \(w_0\in J\).
Each carries an anticyclotomic additive ray character
\(\lambda_j:G_E\to\ZZ/\ell^N\ZZ\), trivial at the coefficient places,
of tame conductor supported on \(j,j^c\), with a chosen tame generator
\(\tau_j\) mapping to \(1\). At each later place \(j\) the chosen
arithmetic Frobenius \(\phi_j\) satisfies
\[
          \lambda_{w_0}(\phi_j)=1,\qquad
          \lambda_k(\phi_j)=0\quad(k\in J,\ k\ne w_0).
\]
The first Frobenius may be prescribed in any given neighborhood of
an inertia element at an old place away from \(\ell\).
There are slopes \(s_{ij}\in\mathcal O\), zero at the first place.
At a later place impose the two conditions
\[
 z(\tau_j)=s_{ij}z(\phi_j),\qquad
 z((\tau_j^c)^{-1})=s_{ij}z((\phi_j^c)^{-1}),
 \qquad g^c=cgc^{-1}.
\]
A zero slope means the unramified condition at both places. Whenever
\(s_{ij}\ne0\), the decomposition groups at the pair act trivially on
\(W_{i,a_0}=T_i/\pi^{a_0}T_i\) and on \(\mu_{\ell^N}\).
For the resulting Selmer group at each \(1\le a\le a_0\) one has
\[
 \operatorname{length}\operatorname{Sel}_{J}(W_{i,a})
 \le v_\pi(I_{t_i})
       +\sum_{\substack{v\in S\\v\nmid\ell}}
           \sup_{1\le b\le a_0}\operatorname{length}H^0(E_v,W_{i,b}^*).
\]
A member of the Selmer group is represented throughout by the unique
cocycle vanishing at \(\eta_i\).
\end{theorem}
\begin{proof}
The unique normalization follows since \(y_i-1\) is a unit. The local
conditions are conditions on cohomology classes, including the two
displayed equations when the local actions are trivial.

Choose \(w_0,\lambda_{w_0}\) by the idele construction of
Lemma~\ref{lem:raycoordinates726}. To prescribe its Frobenius near
an old inertia element, invert the coefficient places and choose the
class-group generating primes away from that old place and its conjugate.
The Kummer extension of the resulting units is unramified there,
as verified below. The inertia element fixes the extension, so
Chebotarev in a compositum with the finite quotient specifying the
neighborhood gives the required prime. Multiplication of its Frobenius
by tame inertia at the new prime preserves this finite quotient and can
make the new character zero on that Frobenius. Avoid the finite set
of primes ramified in the compositum when making the choice.
Impose the unramified condition there for every \(i\). Its addition
does not change any of the original Selmer groups. The sum \(D\) of
their ordinary lengths over \(i\) and \(1\le a\le a_0\) is finite
by Proposition~\ref{prop:sharpminors727}, or its proof with a
uniformizer power. The local finite conditions are given by integral
equations by Lemma~\ref{lem:localevaluation728}. We shall take
\(d_0=1+D\).

Fix one lattice for the moment and omit \(i\). Let \(H\) be the closed
subgroup generated by \([g,\eta]\) and \([g,\eta^c]\), \(g\in G_E\),
with \([g,h]=ghg^{-1}h^{-1}\). Normality follows because a conjugate of
\([g,\eta]\) by \(b\) is \([bg,\eta][b,\eta]^{-1}\), and the same
holds for \(\eta^c\). It is stable under \(c\). The actions of \(T\)
at \(\eta,\eta^c\) are the scalars \(y,y^{-1}\) by polarization and
\(\kappa(\eta)=1\). Hence \(H\) acts trivially on \(T\), its Tate dual
and all the abelian characters. For a normalized cocycle \(z\) into
\(W_a\),
\[
                z([g,\eta])=(1-y)z(g).
\]
Indeed \(g\eta g^{-1}=[g,\eta]\eta\) and evaluating both sides gives
the formula. In particular, if some value of \(z\) has a unit
coordinate, the span of \(z(H)\) is all of \(W_a\). That span is
stable under \(T(G)\) by the cocycle conjugation identity on \(H\).
It contains a vector with a unit coordinate and is therefore the
whole module by the full matrix span.

We record which new primes can be used. For every \(h\in H\) and
every neighborhood of \(\tau_{w_0}h\) in \(G_E\), there is a new
prime \(w\) with arithmetic Frobenius lift in this neighborhood,
with \(T\bmod\pi^{a_0}\) and \(\mu_{\ell^N}\) locally trivial.
The old character values at its Frobenius are \(1\) for \(w_0\)
and \(0\) for the others. It supplies a ray character as in the
statement that vanishes on the decomposition groups at all earlier
places except possibly the first pair.
Here is the prime construction. In the idele argument of
Lemma~\ref{lem:raycoordinates726}, take the set of inverted primes
to include all coefficient places and all earlier nonfirst pairs.
Require it to avoid the two places of the first pair and the marked
places carrying \(\eta,\eta^c\). One may still make the open class
group trivial. Every ideal class has a representative prime outside
a specified finite set by Chebotarev in the Hilbert class field.
Adjoin such primes for generators and their conjugates. The resulting
Kummer extension of the units is unramified at the omitted places, since \(\ell\) is invertible there and extracting \(\ell^N\)-th roots of
units gives finite étale algebras over their valuation rings.
Consequently \(\tau_{w_0}h\) fixes this entire extension, since
\(h\) is generated by commutators with inertia at omitted places.
Its image under every old character on \(H\) is zero, and its image at
\(\tau_{w_0}\) is the stated one. Both \(T\) and \(\kappa\) are
trivial on \(\tau_{w_0}h\). Impose all the specified conditions
in one finite Galois quotient, including a coset inside the given
neighborhood, \(T\bmod\pi^{a_0}\) at this element and its conjugate,
and the Kummer and character fields. Chebotarev gives a new
degree-one prime with a prolongation realizing the exact group element
in this finite quotient. Avoid the ramified primes of the quotient.
The idele construction at the new prime then applies. Multiplying the
new Frobenius lift by a power of its tame generator makes the new
character zero there, without changing its image in that quotient.
This proves the assertion. In particular every step in the construction
is available for arbitrarily large \(N\).

For an available \(w\), let \(C\) be the primary Selmer cocycles with
ramification relaxed at \(w,w^c\) but with all previous conditions.
Take \(u(z)=(z(\tau_w),z((\tau_w^c)^{-1}))\) and
\(v(z)=(z(\phi_w),z((\phi_w^c)^{-1}))\). These are linear maps to
\(W_a^2\). Their common domain is finite, and \(\ker u\) is exactly
the old Selmer group, since inertia acts trivially and its homomorphisms
factor through the cyclic pro-\(\ell\) tame quotient. If some old
class has \(v(z)\notin u(C)\), Lemma~\ref{lem:slopecount729} permits
a nonzero residue \(d\) with a strict decrease for this precision
and a weak decrease at all precisions up to \(a_0\). In fact the
union of their bad residue sets has at most \(2r a_0\) members.
Set \(s_{iw}=d^{-1}\) using any unit lift of the chosen residue.
Give all other lattices slope zero, retaining exactly their old
groups. Append this place whenever a failure \(v(z)\notin u(C)\)
occurs for some lattice, precision and available place. Each step
lowers the total length by an integer at least one. After at most
\(D\) such steps, every old class at every precision has its two
Frobenius values in the image of the relaxed pair of tame evaluations
at every available new place.

It remains to bound the lengths at this stopping point. Fix \(i,a\).
If the group is nonzero, let \(\pi^b\), \(1\le b\le a\), be its
exponent. Multiplying by \(\pi^{a-b}\) identifies the group at
precision \(b\) with the subgroup at precision \(a\) killed by
\(\pi^b\), hence with the entire group. This follows directly on raw
cocycles from the fixed linear equations, including those above
\(\ell\), and the normalization. In precision \(b\) choose a class
\(z\) of full order. Some coordinate of one of its values is a unit,
so \(z(H)\) spans \(W_b\).

Let \(y'\) be a cocycle representing any class in the dual Selmer
group at this precision, normalized by \(y'(\eta)=0\). Write
\(\langle\, ,\,\rangle:W_b\times W_b^*\to
\pi^{-b}\mathcal O/\mathcal O(1)\) for evaluation, identifying the
root-of-unity factor using compatible primitive roots when computing
local invariants. On \(H\times H\) put
\[
 B(h,k)=\langle z(h),y'(k)\rangle
                          -\langle z(k^c),y'(h^c)\rangle .
\]
We first prove \(B(h,h)=0\). At the first pair both classes are
unramified, since the unramified conditions at a place away from
\(\ell\) are exact annihilators for unramified modules. As inertia
there acts trivially, both cocycles vanish on its tame inertia,
including its conjugate. Thus the diagonal expression takes the same
value at \(\tau_{w_0}h\) and at \(h\). Finite precision makes it
locally constant. Choose an available \(w\) whose Frobenius \(\phi\)
has that same value. The stopping property supplies \(z'\) in the
relaxed group with
\[
        z'(\tau)=z(\phi),\qquad
        z'((\tau^c)^{-1})=z((\phi^c)^{-1})=-z(\phi^c).
\]
For the reciprocity calculation choose \(\tau\) normalized on a
uniformizer radical using the fixed primitive roots. This can be done
by choosing at each new prime the residue logarithm in
the ray character so that this generator has character value \(1\).
Then \((\tau^c)^{-1}\) is the corresponding normalized generator at
the conjugate place, since \(c\) inverts the roots of unity. An
arithmetic Frobenius there is \(\phi^c\).

Apply Brauer reciprocity to \(z'\) and \(y'\). All local pairings
away from the new pair vanish by the old orthogonal conditions.
At the new pair \(y'\) is unramified. The local invariant of the
cup product is evaluation of its Frobenius value on the tame value
of \(z'\), up to the same sign at both primes. This follows by
subtracting the unramified part and using the Kummer character of a
uniformizer as in Proposition~\ref{prop:evenreciprocity712}, coordinate
by coordinate. Thus the two remaining terms give \(B(\phi,\phi)=0\).
For ramified coefficients the numerical invariant uses a trace to
\(\QQ_\ell/\ZZ_\ell\). Reciprocity for every scalar multiple of \(z'\)
then gives vanishing in the coefficient-valued pairing
by trace duality with the inverse different. This proves \(B(h,h)=0\).

The restrictions of both cocycles to \(H\) are homomorphisms.
Expansion at \(hk\) consequently gives \(B(h,k)+B(k,h)=0\).
If \(h\) is conjugated by \(\eta\) in the first argument, \(B(h,k)\)
is multiplied by \(y\) through the actions of
\(\eta\) on \(W_b\) and \(\eta^c\) on \(W_b^*\), both equal to \(y\).
The same substitution in the second argument multiplies \(B(k,h)\)
by \(y^{-1}\). Using the antisymmetry twice gives
\((y^2-1)B(h,k)=0\), whence \(B=0\).

There is now a unique linear map \(F:W_b\to(W_b^*)^c\) taking
\(z(h)\) to \(y'(h^c)\). Indeed any linear relation among the
\(z(h)\) gives the same relation among the \(y'(h^c)\), after pairing
with every \(z(k^c)\), by \(B=0\). These vectors span \(W_b\) and
evaluation is perfect. Conjugating \(h\) by elements of \(G_E\)
shows that \(F\) is equivariant. It depends linearly on the
normalized dual cocycle. If \(F=0\), the latter vanishes on \(H\),
and the commutator formula for the dual with scalar \(y^{-1}\)
forces the normalized cocycle to vanish everywhere. This embeds
the dual Selmer group into the homomorphism group computed in
Theorem~\ref{thm:selftwist728}. Its length is at most \(v_\pi(I_t)\).
Theorem~\ref{thm:balancedcount728} now bounds the primary group
as stated. At auxiliary pairs of zero slope the unramified condition
contributes zero even when Frobenius acts nontrivially. This completes
the proof at all precisions and for all the lattices.
\end{proof}

On the marked field of Corollary~\ref{cor:actualmarks714}, suppose
the only ramified places of the untwisted lattices away from \(\ell\)
are the places of the two original levels, with the marked types of
Section~\ref{sec:markedmatrix715}. Take \(S\) to consist of these and
the coefficient places. The results on the two determinants then
bound the right side of Theorem~\ref{thm:tiedlength729}, normalized
by \(v_\pi(\ell)\), by
\[
                (2r_i+1)v_\ell(G_0).
\]
In fact Corollary~\ref{cor:congcost728} bounds the global term, and
the last estimate of Section~\ref{sec:tamerows725} bounds the sum
of the local dual invariant lengths by the same determinant
valuations used there. This applies at the spherical coefficient
primes of Corollary~\ref{cor:scalarfl724} for the indicated main
summands, with their prescribed polarization. The slopes just chosen
give extra local equations for the original lattice. The passage from
these equations to specializations by genuine finite-order characters
requires an estimate for dropping the equations after a deformation.
\clearpage

\subsection{Short Selmer groups for genuine twists of the main summands}
\label{sec:smalltwisted730}

The local equations at the split pairs of
Theorem~\ref{thm:tiedlength729} apply first to the untwisted
representation. Here we remove those equations and obtain a bound for
an actual twist, using the same character for all the summands and all
its prescribed powers. Write \(v(\ell)=1\). All lengths in this
subsection are divided by the absolute ramification index of the
coefficient ring.

\begin{lemma}[Tame inertia along a line of characters]\leavevmode
\label{lem:linetame730}
Use the prime pairs of Theorem~\ref{thm:tiedlength729} at precision
\(\ell^M\) for one lattice \(T\). Write \(s_j\) for its slopes.
Let \(0<\rho\le1\), \(N\rho\ge M\), and let \(y\) be an integral
algebraic \(\ell\)-adic number with \(v(y)\ge\rho\). Set
\[
          x_{w_0}=y,\qquad x_j=s_j y\quad(j\ne w_0),\qquad
          \psi_x(g)=\prod_j(1+x_j)^{[\lambda_j(g)]}.
\]
Here the brackets mean representatives modulo \(\ell^N\) as in
Proposition~\ref{prop:sharpminors727}. Modulo \(\ell^M\) this is
a character. For any cocycle \(z\) of the resulting action on
\(T/\ell^M\), after coefficient extension to contain \(y\), at each
later pair the two quantities
\(z(\tau)-s_j z(\phi)\) are killed by
\[
                          (1+y)T(\phi)-I .
\]
Use \((\tau,\phi)=(\tau_j,\phi_j)\) and
\(((\tau_j^c)^{-1},(\phi_j^c)^{-1})\), respectively. At \(y=0\),
the two tame values at the first pair are killed by \(T(\phi)-I\)
for the corresponding Frobenius elements.
\end{lemma}
\begin{proof}
If \(v(b)\ge\rho\), the sum \(\sum_{h=0}^{\ell-1}(1+b)^h\) has
valuation at least \(\min(1,(\ell-1)v(b))\ge\rho\).
This follows by expanding it as
\(\sum_{d=1}^{\ell}\binom{\ell}{d}b^{d-1}\).
Also \((1+b)^\ell-1\) has valuation at least \(\rho\). Factoring
the longer geometric sum into \(N\) sums of length \(\ell\) gives
\[
        v\left(\sum_{h=0}^{\ell^N-1}(1+b)^h\right)\ge N\rho\ge M .
\]
In particular \((1+b)^{\ell^N}\equiv1\bmod\ell^M\).
Each power in \(\psi_x\) thus defines a character modulo \(\ell^M\).
At a selected twist place the underlying \(T\) is unramified and the
linear inertia action of the twist is scalar. The affine action
\(u\mapsto z(g)+T(g)\psi_x(g)u\) is trivial on wild inertia,
since there the action on the lattice is trivial and a continuous
homomorphism from a pro-residue-characteristic group to an
\(\ell\)-primary group vanishes. On the chosen tame generator its
\(\ell^N\)-th power is the identity because the translation part is
multiplied by the geometric sum just estimated, and the scalar part
has \(\ell^N\)-th power one. The norm is \(1\bmod\ell^N\), so the
tame Frobenius relation makes the affine actions of \(\phi\) and
\(\tau\) commute, also for the conjugate inverses.

At a later place those two linear actions are
\(D=(1+y)T(\phi)\) and \((1+s_jy)I\), since only the first ray
character takes value \(1\) on \(\phi\). At the conjugate inverses
the ray-character values are identical by anticyclotomicity.
Commutation of affine actions gives
\((D-I)z(\tau)=s_jy z(\phi)\). When \(s_j\ne0\), one has
\(T(\phi)=I\bmod\ell^M\). When \(s_j=0\) the right side vanishes.
Both cases give the asserted equation. At \(y=0\) all character
values are one and the same affine relation at the first place gives
\((T(\phi)-I)z(\tau)=0\).
\end{proof}

\begin{theorem}[Simultaneous twists with bounded divisible Selmer groups]
\leavevmode\label{thm:twistedselmer730}
Use the lattices \(T_i\), field \(E\), set \(S\) and integral
coefficient-place evaluation equations of
Theorem~\ref{thm:tiedlength729}. Suppose there is an inertia element
\(g_*\) at a place of \(S\) away from \(\ell\) such that
\(\det(1-T_i(g))\ne0\) for \(g=g_*,(g_*^c)^{-1}\) and every \(i\).
Define
\[
 \begin{split}
 B_i={}& v(I_{t_i})
       +\sum_{\substack{w\in S\\ w\nmid\ell}}
                 \sup_{b\ge1}\operatorname{len}_{\rm norm}
                                     H^0(E_w,W_{i,b}^*)\\
       &+\sum_{g\ {\rm in}\ (g_*,(g_*^c)^{-1})} v(\det(1-T_i(g))).
 \end{split}
\]
Assume each \(B_i\) is finite. Given \(m\ge1\) and
\(\varepsilon,\eta_0>0\), there exists a finite-order anticyclotomic
character \(\vartheta\), trivial on all decomposition groups above
\(\ell\), ramified at a finite set of new split pairs, such that for
every \(i\) and \(0<|s|\le m\) the following holds.
The cocycles for \(T_i\otimes\vartheta^s\) with divisible coefficients,
normalized to vanish at \(\eta_i\), obeying the coefficient-place
equations and unramified outside \(S\) and the new pairs, have
normalized length at most \(B_i+\varepsilon\).
The sum of the normalized valuations of \(\vartheta^s(\gamma)-1\)
at normalized tame generators of the new places is at most \(\eta_0\).
All conclusions hold after further finite coefficient extensions.
\end{theorem}
\begin{proof}
Choose \(M>\max_i B_i+\varepsilon+1\) integral. Write
\(a_0=M v_\pi(\ell)\) and make the unramified coefficient extension
of Theorem~\ref{thm:tiedlength729}. This leaves the normalized bounds
unchanged. For the invariant groups it follows by flat base change
in the kernels of the action equations, with \(\pi\) still a
uniformizer. The trace ideal and determinants extend by scalars.
That theorem bounds the number of prime pairs by an integer \(d_0\)
depending on the untwisted data and \(a_0\), independently of \(N\).
Choose a rational \(0<\rho\le1\) so small that
\(2(\max_i r_i)d_0\rho<\min(1,\varepsilon/2)\).

We choose \(N\) before constructing any new primes. Require
\(N\rho\ge M\), and choose it also at least as large as the
cyclotomic exponent in Theorem~\ref{thm:allraypowers726} for at most
\(d_0\) variables, witness coordinate bound \(\rho\), signed powers
up to \(m\), tolerance \(\varepsilon/2\), tame bound \(\eta_0\), and
sum-of-values bound \(2m\sum_i(B_i+1)\). Its proof gives the same
exponent for every smaller number of variables by using the bounds
for \(d_0\) throughout. Now choose the pairs, ray characters and
slopes by Theorem~\ref{thm:tiedlength729}. Prescribe the first
Frobenius sufficiently close to \(g_*\) for the two determinant
valuations in \(B_i\) to agree at both selected Frobenius elements,
for all \(i\). This is an open condition because the finitely many
specified determinants are nonzero.

For a fixed \(T_i\), form the polynomial matrix of
Proposition~\ref{prop:sharpminors727}, allowing ramification at all
the pairs, retaining the normalization and the fixed coefficient-place
equations. Make a second matrix with the equations of
Theorem~\ref{thm:tiedlength729} at the later pairs only.
Both matrices represent the indicated cocycles also at the points of
Lemma~\ref{lem:linetame730}. Indeed, the actions there are genuine
modulo \(\ell^M\) and trivial on the same group \(H\) used to form the
finite quotient in Proposition~\ref{prop:sharpminors727}. Its proof
including the elementary divisor calculation applies verbatim.

At the origin the length for the second matrix is at most \(B_i\).
With all the equations it has the bound from
Theorem~\ref{thm:tiedlength729}. Dropping the equations at the first
pair adds at most the sum of the two normalized determinant valuations
there, by Lemma~\ref{lem:linetame730} at \(y=0\). Indeed the quotient
injects through its two tame values into the two matrix kernels on
the finite coefficient modules, each of length bounded by its
determinant valuation. Select a column-size minor of the second matrix
attaining its length at the origin. Restrict this polynomial to the
line \(X_{w_0}=Y,\ X_j=s_{ij}Y\) and call it \(P^0(Y)\).
It has \(v(P^0(0))\le B_i\). Define
\[
      P^1(Y)=\prod_{j\ne w_0}
             \ \prod_{\phi=\phi_j,(\phi_j^c)^{-1}}
                           \det((1+Y)T_i(\phi)-I).
\]
Each factor has degree \(r_i\) with unit leading coefficient. Hence
\(v_\rho(P^1)\le 2r_i d_0\rho\), and
\(v_\rho(P^0)\le v(P^0(0))\). There is a single algebraic \(y\)
with \(v(y)=\rho\) at which both polynomials attain their Gauss
valuations. To see this, write \(Y=\varpi_\rho U\) with
\(v(\varpi_\rho)=\rho\), divide each polynomial by a coefficient of
least valuation, and avoid the finitely many roots of the two resulting
nonzero residue polynomials and zero in
\(\overline{\mathbb F}_\ell\). A unit lift of the selected residue
gives the required \(y\).

At this point on the line, forgetting the later equations increases
the length by at most \(v(P^1(y))\). The quotient injects by their
left sides into the kernels of the matrices in
Lemma~\ref{lem:linetame730}, and each kernel has normalized length
at most its determinant valuation. The group still satisfying the
equations has length at most \(v(P^0(y))\) by its chosen minor.
Choose a column-size minor \(P_i\) of the first matrix attaining its
length at this point \(x_i\). We have proved
\[
         v(x_{i,j})\ge\rho,\qquad
         v(P_i(x_i))\le B_i+2r_i d_0\rho .
\]
These are witnesses for the actual common prime set. For each signed
power use the same witness and polynomial for that lattice.

Apply the polynomial specialization argument of
Theorem~\ref{thm:allraypowers726}. It depends on the inertia
coordinates and their orders and does not require the earlier
characters to vanish at the later Frobenius elements. The characters
constructed here have exactly the required inertia coordinates, since
each ramifies only at its own pair and takes value \(1\) at its
chosen tame generator. We obtain roots \(\zeta_j\) and one character
\(\vartheta=\prod_j\zeta_j^{\lambda_j}\). The first matrix at the
point \((\zeta_j^s-1)_j\) therefore has normalized kernel length at
most \(B_i+\varepsilon<M\) for the cocycle equations modulo \(\ell^M\).
Here kernel length refers to the equations before adjoining
\(\ell^M I\), as in Proposition~\ref{prop:sharpminors727}.
Lemma~\ref{lem:precisionlength727} gives the divisible-coefficient
assertion. The same minor works over every further finite coefficient
extension, since its normalized valuation is unchanged.
The inertia estimate is the last calculation of
Theorem~\ref{thm:allraypowers726}.
\end{proof}

\begin{corollary}[The higher-rank part of the trace bound at spherical primes]
\leavevmode\label{cor:maincost730}
On the shorter family of Corollary~\ref{cor:actualmarks714}, take
finitely many untwisted marked main summands as above, of rank \(r_i\)
and with the specified squarefree or rank-one local types at the two
old level primes. Assume their ramification away from \(\ell\) is
confined to those levels. Take a spherical coefficient prime in
the common range of Corollary~\ref{cor:scalarfl724}. At this node
one twist and its first \(m\) powers and inverses give, for each main
summand, the estimates in the notation of Theorem~\ref{thm:ribettrace720}
\[
    l_e=0\quad(0\le e<t_i),\qquad
    \frac{\lambda_{T_i\otimes\vartheta^s}}{v_\pi(\ell)}
             \le (4r_i+1)v_\ell(G_0)+\varepsilon .
\]
Here the trace families have the hyperspecial level and fixed
infinity type of Corollary~\ref{cor:scalarfl724} and contain the
selected twisted main summand. The estimate is independent of the
later auxiliary primes and their precision. For each \(\ell\) one
may take \(\varepsilon=2^{-\ell}\).
\end{corollary}
\begin{proof}
At one of the old level primes different from \(\ell\), take for
\(g_*\) a full tame generator. At any place of the prescribed type
the inertia spectrum excludes \(1\). The same holds at the
conjugate inverse, by polarization and trivial cyclotomic inertia.
Section~\ref{sec:markedmatrix715} bounds each determinant valuation
by \(r_i v_\ell(G_0)\). All other old local invariant contributions
in \(B_i\) have total normalized length at most
\(2r_i v_\ell(G_0)\), as in Theorem~\ref{thm:tiedlength729}.
The global congruence term has normalized length at most
\(v_\ell(G_0)\) by Corollary~\ref{cor:congcost728}. The marked scalar
and matrix hypotheses follow from Theorem~\ref{thm:markedmatrix715}.
Together these verify the hypotheses and give the asserted uniform
upper bound for \(B_i\).

For the trace cocycles in \(\mathcal V_{T_i\otimes\vartheta^s}\),
the normalization is the same, because the twist is unramified at
the original marked inertia element. Its local components are trivial
above \(\ell\), so the fixed crystalline evaluation equations at these
places are precisely those of the untwisted summand. They hold for
the trace cocycles by the lattice lifting argument of
Theorem~\ref{thm:fllift724}, after choosing some finite precision
for each such cocycle and scaling back to divisible coefficients.
At the later auxiliary places of the trace construction the target is
unramified and the cocycles vanish on the chosen tame generators, so they are unramified there. Thus the trace cocycles belong
to the group bounded by Theorem~\ref{thm:twistedselmer730}.
The scalar colengths vanish by Corollary~\ref{cor:scalarfl724}.
Finally the added errors are summable with weights \(\log\ell\),
since \(\sum_{n\ge2}2^{-n}\log n<\infty\).
\end{proof}

This treats the Selmer length for the chosen twists in the indicated
coefficient range. To obtain the integral period bound of
Section~\ref{sec:missing} the twist must also control the ratios of
the Hecke values in those same periods. At coefficient primes outside
the common Fontaine--Laffaille range one still needs the quantitative
comparison with the crystalline local equations.
\clearpage

\section{Critical values and Selmer lengths on a common interval}\label{nav:selmer-critical}

\subsection{Critical Hecke values at the twists used for Selmer groups}
\label{sec:samecritical731}

Fix a CM field \(E\) of Corollary~\ref{cor:fieldrelations712} and
write \(c\) for its conjugation. Use arithmetic Frobenius for
ideal characters. Fix a complex embedding of the algebraic numbers
and an embedding \(\iota:\overline{\QQ}\hookrightarrow
\overline{\QQ}_{\ell}\). Write \(v(\ell)=1\). A character \(u\) in
this subsection is a unitary algebraic Hecke character of \(E\) with
\(u^c=u^{-1}\). Its conductor \(\mathfrak f_0\) is nontrivial,
prime to \(\ell\), and invariant under \(c\). At the integer \(z\ge1\)
we require that the two ideal exponents of
\(\alpha(\mathfrak a)=u(\mathfrak a)(N\mathfrak a)^{-z}\)
at each conjugate pair of embeddings are, in some order,
\(k_\sigma-z,-k_\sigma-z\), with \(k_\sigma\ge z\).
These are precisely the exponents obtained from the norm-one
characters of weight \(-k_\sigma\) and their inverses by pulling
back along \(x\mapsto x/\bar x\) and multiplying by \(N^{-z}\).
In particular \(\alpha\) is critical at zero. In the notation
\(\beta-\alpha'\) for a critical infinity type, the exponents on
the negative side have \(\alpha'\ge1\), and the others have
\(\beta\ge0\).

For a set \(J\) of split primes containing at most one prime of each
conjugate pair, disjoint from \(\mathfrak f_0\ell\), put
\(\mathfrak f_J=\mathfrak f_0\prod_{j\in J}jj^c\).
The notation \(L_{\mathfrak f}(z,u)\) omits all primes dividing
\(\mathfrak f\). Our ratio for a finite-order anticyclotomic
character \(\psi\) of conductor dividing \(\prod_{j\in J}jj^c\) is
\[
  R_J(u,z,\psi)=
    \frac{L_{\mathfrak f_J}(z,u\psi)}
         {L_{\mathfrak f_0}(z,u)}.
\]
For use with a Galois character we identify its finite-order values
with complex roots of unity using \(\iota\). We first construct
integral coefficients on the ray class groups whose evaluations give
these ratios up to one constant.

\begin{lemma}[Keeping additional places split in the ray characters]
\leavevmode\label{lem:fixedsplit731}
In Theorem~\ref{thm:tiedlength729} fix any finite conjugation-stable
set \(V\) of primes disjoint from \(S\). The ray characters can all
be taken trivial on the decomposition groups at \(V\), with the
same bound \(d_0\) on the number of pairs.
\end{lemma}
\begin{proof}
At each use of the idele construction in that theorem include \(V\)
among the inverted primes \(T\), and choose all the new pairs
outside \(V\). The extra local factors \(E_v^\times\) of \(D_T\)
are killed by the defining residue character. The Kummer generators
are still units at the distinguished old inertia place in the first
step. In each later step they are units at the first pair and at
the old marked pair used in that step. Extracting their
\(\ell^N\)-th roots is unramified at those places. Hence the
prescribed Frobenius elements still fix the Kummer extension in
the proof of Theorem~\ref{thm:tiedlength729}. Chebotarev permits the
same choices with the extra finite exclusions. The primary Selmer
conditions used to count the steps involve the original \(S\).
Inverting primes in the idele construction adds no allowed
ramification to those Selmer groups. Their initial lengths and the
bound \(d_0\) are consequently the same.
\end{proof}

\begin{proposition}[An integral distribution at all selected tame levels]
\leavevmode\label{prop:cmraymeasure731}
Suppose \(\ell\nmid \lvert D_E\rvert\). Fix an integral prime ideal
\(\mathfrak c\) away from \(\mathfrak f_0\ell\) such that
\(A=N\mathfrak c\,\alpha(\mathfrak c)\ne1\). For any of the
collections \((\lambda_j)_{j\in J}\) constructed above with
\(\mathfrak c,\mathfrak c^c\) split in their character fields, put
\(H_J=(\ZZ/\ell^N\ZZ)^J\). There is a constant \(C\ne0\) independent
of \(J,N\) and integral elements
\(\Theta_J\in\overline{\ZZ}_\ell[H_J]\) such that, for every
character \(\psi\) of \(H_J\),
\[
      \psi(\Theta_J)=\iota\!\left(C L_{\mathfrak f_J}(z,u\psi)\right).
      \tag{*}\label{eq:raycritical731}
\]
On the right, \(\psi\) is pulled back to ideals through the
\(\lambda_j\). The argument of \(\iota\) in this formula is
algebraic. In particular \(\Theta_\varnothing=C L_{\mathfrak f_0}(z,u)\)
after applying \(\iota\) to the product, and this value is nonzero.
\end{proposition}
\begin{proof}
We use the partial ray-class form of the Kings--Sprang integrality
theorem. We spell out its hypotheses because a denominator depending
on the auxiliary level would affect the choice of \(N\).
For a critical ideal character \(\chi\) of nonunit conductor
\(\mathfrak f\), let \(L_{\mathfrak f}(\chi,0,b)\) denote the partial
value for a ray class \(b\), including the values of \(\chi\) in
the sum over ideals. For the CM type determined by the negative
exponents, fix the CM abelian variety with Betti lattice
\(\mathcal O_E\) and bases of the eigendifferentials on it and its
dual. The theorem supplies the period factor \(D\) from these
bases and the infinity type. For each integral representative
\(\mathfrak b\) of \(b\), prime to \(\mathfrak f\mathfrak c\ell\),
it gives integrality at \(\iota\) of
\[
 D\left(N\mathfrak c\,\chi(\mathfrak c)
           L_{\mathfrak f}(\chi,0,b)
          -L_{\mathfrak f}(\chi,0,b[\mathfrak c])\right).
 \tag{1}\label{eq:partialcm731}
\]
Such representatives exist by the ray class theorem. The only
localizations in the partial-value formula beyond its base ring
are at the modulus and at primes dividing
\(N(\mathfrak b\mathfrak c)\). All are units at the specified
coefficient place. This is Corollary~4.13 in the cited version
of the paper. \cite{II-II-kingssprang731}

For clarity, the CM data for \(D\) can be fixed before varying
\(\mathfrak f\). The existence theorem of complex multiplication
defines the CM variety for the indicated lattice and CM type over
a number field after choosing a polarization. It has good reduction
at the selected place after a finite extension. Indeed semistable
reduction gives a toric part whose rational character group, if
nonzero, would carry a unital action of \(E\). Its dimension is at
most \([E:\QQ]/2\), too small for a nonzero \(E\)-module. The action
comes from extension of endomorphisms to the Néron model. Good
reduction follows. After also containing the embeddings of \(E\),
the eigendifferentials split over the local integer ring since
\(\ell\nmid D_E\). Each eigensummand for the CM type, and each
dual eigensummand of the appropriate type, is free of rank one.
Choose bases there once. For each new level extend the ring to
define the required prime-to-\(\ell\) torsion sections. The old
bases remain bases. Thus the nonzero complex scalar \(D\) used in
(1) is identical at every level. This statement concerns a fixed
coefficient place. Its scalar need not be bounded as \(E\) varies.

Apply (1) with \(\mathfrak f=\mathfrak f_J\). To ensure an exact
conductor, let \(\psi_0\) be primitive of order \(\ell^N\) in every
coordinate, and take \(\chi=\alpha\psi_0\). At each new prime
only its own coordinate ramifies, with tame conductor exponent one.
At every old prime \(\psi_0\) is unramified. The conductor is
therefore exactly \(\mathfrak f_J\), including for \(J=\varnothing\)
with \(\psi_0=1\). The infinity type is unchanged. On a ray class
\(b\) a partial value of \(\alpha\psi_0\) is the partial value of
\(\alpha\) times \(\psi_0(b)\), first by the convergent series and
then by continuation. Removing the common root of unity
\(\psi_0(b[\mathfrak c])\) from (1) shows that
\[
 d_b=D\left(A L_{\mathfrak f_J}(\alpha,0,b)
              -L_{\mathfrak f_J}(\alpha,0,b[\mathfrak c])\right)
\]
is algebraic and integral at \(\iota\). Map each ray class \(b\)
to \(\lambda(b)=(\lambda_j(b))_j\) by reciprocity. Set
\(\Theta_J=\sum_b\iota(d_b)[\lambda(b)]\). Evaluating and reindexing
the second sum gives
\(D(A-\psi(\mathfrak c)^{-1})L_{\mathfrak f_J}(\alpha\psi,0)\).
The finite local character at \(\mathfrak c\) is trivial by
construction. Since \(\alpha=uN^{-z}\), this proves (*) with
\(C=D(A-1)\). The proof involves no division by the size of a ray
class group.

It remains to verify nonvanishing and the interpretation of the
absolute value. For real \(s>1\), conjugate unramified Euler factors
in \(L_{\mathfrak f_J}(s,u\psi)\) have positive product. At a
conjugation-fixed prime the value is \(1\) or \(-1\), so its factor
is positive as well. The product converges absolutely. At \(s=1\)
the classical analytic continuation of Hecke \(L\)-functions makes
both \(L_{\mathfrak f_J}(s,u\psi)\) and
\(L_{\mathfrak f_J}(s,(u\psi)^2)\) holomorphic there, since their
nonzero angular infinity types rule out a norm character. For
\(s>1\) the product of \(\zeta_{E,\mathfrak f_J}(s)^3\),
\(\lvert L_{\mathfrak f_J}(s,u\psi)\rvert^4\), and
\(\lvert L_{\mathfrak f_J}(s,(u\psi)^2)\rvert\) is at least one.
Its logarithm is a sum over prime powers with numerators
\(3+4\cos t+\cos(2t)=2(1+\cos t)^2\). A zero of the first
Hecke value at one would force that product to tend to zero, since
the zeta function has a simple pole. The value at one is thus
positive by continuity. In particular the ratios \(R_J\) agree
with the ratios of absolute values at every integer in question.
\end{proof}

\begin{lemma}[The Euler factors along the first ray character]
\leavevmode\label{lem:criticalwitness731}
Suppose an old inertia element \(g_*\) away from \(\ell\) has
\(u_\ell(g_*)\ne1\), where \(u_\ell\) is the avatar with the
prescribed ideal values at arithmetic Frobenius. Write
\[
    b=v(\Theta_\varnothing),\qquad
    d=v(u_\ell(g_*)-1),\qquad \beta=b+2d .
\]
On a set \(J\) from Theorem~\ref{thm:tiedlength729}, choose the
first Frobenius sufficiently close to \(g_*\) modulo \(\ell^M\)
for \(u_\ell\). Represent \(\Theta_J\) by the polynomial
\[
   P_J(X)=\sum_{h\in H_J}(\Theta_J)_h
                       \prod_{j\in J}(1+X_j)^{[h_j]},
       \qquad 0\le[h_j]<\ell^N .
\]
If \(0<\rho\le1\) is rational and
\(\beta+2\lvert J\rvert\rho<M\le N\rho\), this polynomial has a
witness \(x\) with \(v(x_j)\ge\rho\) and
\(v(P_J(x))\le\beta+2(\lvert J\rvert-1)\rho\).
\end{lemma}
\begin{proof}
Write \(j_0\) for the first prime. At augmentation, omitting
\(j_0,j_0^c\) multiplies \(\Theta_\varnothing\) by
\((1-a_0q_0^{-z})(1-a_0^{-1}q_0^{-z})\), with
\(a_0=u_\ell(\phi_{j_0})\) and \(q_0=Nj_0\).
Since \(q_0\equiv1\bmod\ell^N\), these two factors agree modulo
\(\ell^M\) with \(1-u_\ell(g_*)^{\pm1}\). Their valuations are
both \(d<M\). Thus \(P_{\{j_0\}}(0)\) has valuation \(\beta\).

Put \(Z=1+Y\) on the \(j_0\)-coordinate and put all other \(X_j=0\).
For \(Z^{\ell^N}=1\), the value of the corresponding character at
arithmetic Frobenius at every later \(j\) is \(Z\), by the choice
of \(\lambda_{j_0}\). Removing each pair \(j,j^c\) therefore
contributes
\[
   (1-a_j q_j^{-z}Z)(1-a_j^{-1}q_j^{-z}Z^{-1}),
   \qquad a_j=u_\ell(\phi_j),\quad q_j=Nj .
\]
This Euler deletion follows from (*) because the scalar there is
fixed. Multiply each displayed expression by \(Z\), and take the
product \(B(Y)\) of the resulting quadratics over the later primes.
Each has unit leading coefficient, since \(a_j,q_j\) are units.
For \(h=\lvert J\rvert-1\), the polynomial
\[
 Z^h P_J(Y,0,\ldots,0)-P_{\{j_0\}}(Y)B(Y)
\]
vanishes at all the distinct roots of \(Z^{\ell^N}-1\).
Division by this monic polynomial has integral quotient and zero
remainder.

Use the Gauss valuation at radius \(v(Y)=\rho\). The second term
has valuation at most \(\beta+2h\rho\), by the constant coefficient
of \(P_{\{j_0\}}\), the leading coefficient of \(B\), and
multiplicativity. The modulus \(Z^{\ell^N}-1\) has Gauss valuation
at least \(N\rho\). Indeed it factors as \(Y\) times the product of
the \(N\) successive geometric sums of length \(\ell\), each of
valuation at least \(\rho\) by the binomial calculation in
Lemma~\ref{lem:linetame730}, also valid for the Gauss valuation.
The congruence and \(v_\rho(Z^h)=0\) now give
\(v_\rho(P_J(Y,0,\ldots,0))\le\beta+2h\rho\).
Choosing \(Y\) of valuation \(\rho\) avoiding the finitely many
roots of the reduced polynomial on the residue circle attains this
valuation, as in Theorem~\ref{thm:twistedselmer730}.
\end{proof}

\begin{theorem}[One twist for Selmer lengths and critical Hecke values]
\leavevmode\label{thm:commoncritical731}
In Theorem~\ref{thm:twistedselmer730} suppose \(\ell\nmid D_E\).
Fix finitely many \((u,z)\) as above, each with finite conductor
prime to \(\ell\), with \(u_\ell(g_*)\ne1\) at the same old
inertia element as in that theorem. For every \(m\ge1\) and
\(\varepsilon,\eta_0>0\), its bounds \(B_i+\varepsilon\) for the
Selmer groups and \(\eta_0\) at the new inertia groups hold for
one twist \(\vartheta\) for which, at every \(0<|s|\le m\),
\[
    \max\{0,v(\iota(R_J(u,z,\vartheta^s)))\}
          \le 2v(u_\ell(g_*)-1)+\varepsilon .
    \tag{2}\label{eq:samebounds731}
\]
Here \(J\) is the actual auxiliary prime set of that twist.
For the ramified critical pairs on the shorter fields of
Corollary~\ref{cor:maincost730}, take a full tame generator of the
chosen old place \(w_*\) away from \(\ell\). The right hand side is
then at most \(2v_\ell(Nw_*-1)+\varepsilon\).
\end{theorem}
\begin{proof}
Choose one smoothing prime \(\mathfrak c\) outside \(S\) and all old
conductors. For each character prescribe Frobenius sufficiently close
to \(g_*\) for the character \(u_\ell\kappa^{1-z}\), where
\(\kappa\) is cyclotomic. Its value at \(g_*\) is \(u_\ell(g_*)\)
since the inertia is away from \(\ell\). Continuity and
Chebotarev in a finite common quotient therefore give
\(A\ne1\) for every pair at the same \(\mathfrak c\). Exclude the
conjugates of the old places too. Use
Lemma~\ref{lem:fixedsplit731} at
\(V=\{\mathfrak c,\mathfrak c^c\}\).
The constants \(b,\beta\) for all pairs are now fixed finite
numbers, before the order or number of the new characters is chosen.

Take an integer \(M\) exceeding every \(B_i,\beta\) by more than
\(1+\varepsilon\). At this precision let \(d_0\) be the bound of
Theorem~\ref{thm:tiedlength729}. Choose \(0<\rho\le1\) rational
such that \(2d_0\max(1,\max_i r_i)\rho<\min(1,\varepsilon/2)\).
Choose \(N\) with \(N\rho\ge M\), large enough for the polynomial
specialization in Theorem~\ref{thm:allraypowers726} at tolerance
\(\varepsilon/2\), tame bound \(\eta_0\), and total witness bound
\[
     2m\left(\sum_i(B_i+1)+\sum_{(u,z)}(\beta+1)\right).
\]
Construct the actual pairs with the first Frobenius close enough
for every determinant and every \(u_\ell\). The minor for \(T_i\)
constructed in Theorem~\ref{thm:twistedselmer730} has witness
valuation at most \(B_i+2r_id_0\rho\). For every \((u,z)\),
Lemma~\ref{lem:criticalwitness731} provides on these very pairs a
polynomial of witness valuation at most \(\beta+2d_0\rho\).
Include each polynomial with every signed exponent \(s\) in a single
application of the specialization theorem. Its proof uses only the
orders and the independent inertia coordinates of the \(\lambda_j\)
at this step. It yields one choice of roots \(\zeta_j\), hence
\(\vartheta=\prod_j\zeta_j^{\lambda_j}\), with the required Selmer
and tame bounds. Each prescribed power ramifies at every new pair.
Moreover the characters \(\lambda_j\) were differences
\(\xi_j-\xi_j^c\) of idele characters, so the product is pulled
back along \(x\mapsto x/\bar x\) from a finite-order character of
the norm-one torus.

At the selected roots the Hecke polynomial is
\(\vartheta^s(\Theta_J)\), of valuation at most
\(\beta+\varepsilon\). Equation (*) at \(J\) and at the empty level
shows that its quotient by \(\Theta_\varnothing\) is precisely
\(\iota R_J\). Subtracting \(b\) proves (2). At a chosen old place
the given tame character is nontrivial and its order divides
\(Nw_*-1\). If \(\xi\ne1\) is such a root, the identity
\(\prod_{\omega^{Nw_*-1}=1,\ \omega\ne1}(1-\omega)=Nw_*-1\)
and integrality of every factor give
\(v(1-\xi)\le v_\ell(Nw_*-1)\). This proves the final assertion.
\end{proof}

In (2) the normalization of the CM periods cancels individually
for each field and each critical pair. The choice
\(\varepsilon=2^{-\ell}\) is thus available throughout the stated
coefficient range. The comparison in Section~\ref{sec:missing}
still requires the bound for the full integral projected theta
period with its prescribed local tests and the local cohomology
loss at coefficient primes outside the common weight range.
\clearpage

\subsection{The local cohomology loss on a fixed interval of weights}
\label{sec:localfinite732}

Let \(K/\QQ_\ell\) be any finite extension, \(L/\QQ_\ell\) a finite
coefficient field, and \(\mathcal O\) its integer ring. Fix a
\(G_K\)-stable free \(\mathcal O\)-lattice \(T\), and write
\(V=T[1/\ell]\), \(A=V/T\). Allow a fixed finite Galois extension
\(K'/K\). Assume that \(V\) is crystalline over \(K'\) with all its
filtration jumps in an integer interval \(I\) containing zero.
The interval bounds the jumps at every embedding, including all
coefficient embeddings after restriction of scalars. We obtain a
bound independent of the size of an auxiliary family of crystalline
lattices and the precision of their torsion quotients.

A finite \(\ell\)-primary \(G_{K'}\)-module is called torsion
crystalline in \(I\) if it is a quotient of a lattice in a
crystalline \(\QQ_\ell\)-representation with jumps in \(I\).
This property is preserved by submodules, quotients and finite
direct sums. For submodules, take their preimage in a lattice
surjecting onto the given finite module. The preimage has finite
index in that lattice and spans the same rational representation.

Write \(T_n=T/\ell^nT\), \(\mathcal O_n=\mathcal O/\ell^n\mathcal O\).
For a raw continuous cocycle \(a:G_K\to T_n\) let \(U_n(a)\)
be the extension on \(T_n\oplus\mathcal O_n\) with action
\[
            g(v,t)=(gv+t\,a(g),t).
\]
Write \(C_n\) for the set of these cocycles for which \(U_n(a)\)
restricted to \(G_{K'}\), considered over \(\ZZ_\ell\), is torsion
crystalline in \(I\). Send \(a\) into \(Z^1(K,A)\) by dividing
its values by \(\ell^n\) modulo \(T\). Finally let \(Z_f(K,T)\)
consist of the integral raw cocycles whose rational classes lie
in the Bloch--Kato group \(H^1_f(K,V)\).

\begin{theorem}[Finite error and a criterion using two precisions]
\leavevmode\label{thm:localfinite732}
The images of \(C_n\) form an increasing union \(C\), an
\(\mathcal O\)-submodule of \(Z^1(K,A)\), and \(C[\ell^n]=C_n\)
under the specified identification. Its maximal divisible subgroup
is the image \(C_f\) of \(Z_f(K,T)[1/\ell]\) in \(Z^1(K,A)\).
In particular \(F=C/C_f\) is finite. Let
\(R=\operatorname{rank}_{\ZZ_\ell}Z_f(K,T)\). For \(n\ge1\),
\[
 \log_\ell|C_{n+1}|-\log_\ell|C_n|=R
       \quad\Longleftrightarrow\quad \ell^n F=0 .
\]
At any such \(n\), \(\log_\ell|F|=\log_\ell|C_n|-nR\).
All the groups here concern the fixed tuple \((K,K',T,I)\).
\end{theorem}
\begin{proof}
Addition of cocycles uses the direct sum of two extensions, pullback
by the diagonal of their quotients, and pushout by addition on the
two kernels. Multiplication by \(a_0\in\mathcal O\) is pushout by
\(a_0\) on the kernel. Each operation uses subquotients of sums of the
original modules and of \(T_n,\mathcal O_n\), so preserves the torsion
crystalline condition. To pass from \(n\) to \(b\ge n\), pull back the
extension along \(\mathcal O_b\to\mathcal O_n\) and push out along
\(T_n\to T_b,\ x\mapsto\ell^{b-n}x\). The resulting cocycle is
\(\ell^{b-n}a\). This proves the increasing-union assertion.
Conversely suppose a member coming from precision \(b\) is killed
by \(\ell^n\) with \(b\ge n\). Its upper right entries at precision
\(b\) lie in \(\ell^{b-n}T_b\). In \(U_b\) restrict the first
summand to \(\ell^{b-n}T_b\), leaving its second summand intact.
In this submodule the elements \((0,t)\) with \(t\in\ell^n\mathcal O_b\)
are fixed and form a submodule. Quotienting by it gives exactly
\(U_n\) for the original cocycle valued in \(A[\ell^n]\).
This proves the identity \(C[\ell^n]=C_n\), including the direction
in which the representing precision decreases.

In particular \(C[\ell]\) is finite. Indeed there are finitely many
raw cocycles into \(T_1\). Each gives a homomorphism to the finite
affine group with translation group \(T_1\) and linear group the
image of \(G_K\) on \(T_1\). Local fields have finitely many extensions
of bounded degree inside a fixed algebraic closure. Through each
resulting finite quotient there are only finitely many maps.
The Pontryagin dual of the discrete \(\ell\)-primary group \(C\) is
therefore a finitely generated \(\ZZ_\ell\)-module. In detail, its
quotient modulo \(\ell\) is dual to \(C[\ell]\), and lifts of
generators generate the compact dual by topological Nakayama.
One can see the last assertion by passing to the quotient by the
closed span of those lifts. It equals its own \(\ell\)-multiple and
has no nonzero finite \(\ell\)-group quotient, hence is zero.
The structure theorem over \(\ZZ_\ell\) now gives a finite group
modulo the maximal divisible subgroup of \(C\).

To identify that subgroup take the Tate module
\(\varprojlim_n C[\ell^n]\) under multiplication by \(\ell\).
In upper right coordinates the transition maps are reductions
\(T_{n+1}\to T_n\). A compatible family thus gives a continuous
integral \(z:G_K\to T\), with \(U_n(z\bmod\ell^n)\) torsion
crystalline in \(I\) at \(K'\) for every \(n\). The lattice on
\(T\oplus\mathcal O\) defined by \(z\) has these reductions.
The crystalline limit theorem of Liu, with the proof of Gao and
Liu, makes its rational representation crystalline over \(K'\).
To apply the theorem, forget to \(\ZZ_\ell\), twist the whole lattice
so the interval starts at zero in the Hodge--Tate convention used
there, and use the fixed width of the interval at all \(n\).
The underlying rank is finite. \cite{II-II-liulimit732,II-II-gaoliulimit732}

Over \(K'\) both ends of the extension are crystalline. Exactness
of the crystalline period functor on crystalline representations
shows that a lift of the quotient generator exists in the invariant
periods of the extension. Equivalently, the rational class of \(z\)
maps to zero over \(K'\) with coefficients in
\(B_{\mathrm{cris}}\otimes_{\QQ_\ell}V\). This vanishing descends
to \(K\). After subtracting a coboundary, a cocycle zero over
\(G_{K'}\) factors through \(\operatorname{Gal}(K'/K)\) with values
in the \(G_{K'}\)-invariants. Averaging over this finite group in a
\(\QQ_\ell\)-vector space gives a coboundary. Hence \(z\in Z_f(K,T)\).
Conversely such a \(z\) splits the extension over
\(B_{\mathrm{cris}}\) already over \(K\). Over \(K'\) the two
ends have full crystalline periods. The splitting and these periods
give full crystalline periods for the middle term. Its jumps remain
in \(I\) by exactness of the de Rham filtration. Its integral
reductions consequently belong to \(C_n\). For coefficients of
degree greater than one, the extension over \(\QQ_\ell\) has one
quotient basis vector for each \(\QQ_\ell\)-basis vector of \(L\).
Splitting the vector belonging to \(1\) is equivalent to splitting
them all by \(L\)-linearity.

We have identified the Tate module with \(Z_f(K,T)\), with its maps
to \(C[\ell^n]\) given by \(z\mapsto z/\ell^n\bmod T\).
By the structure theorem the images of these projections are exactly
the \(\ell^n\)-torsion of the maximal divisible subgroup. Their union
is \(C_f\) as claimed. The kernel of
\(Z_f(K,T)[1/\ell]\to Z^1(K,A)\) is \(Z_f(K,T)\) itself, since a
rational finite cocycle taking integral values lies in that lattice.
Write the elementary divisors of \(F\) over \(\ZZ_\ell\) as
\(\ell^{a_1},\ldots,\ell^{a_b}\). The same structure theorem gives
\[
      \log_\ell|C_n|=nR+\sum_{j=1}^b\min(n,a_j).
\]
Taking first differences proves the assertion about two precisions.
\end{proof}

\begin{proposition}[Independence of the coefficient extension]
\leavevmode\label{prop:coeffdefect732}
For a finite extension of coefficient rings \(\mathcal O'/\mathcal O\)
the finite quotient for \(T'=T\otimes_\mathcal O\mathcal O'\) is
\(F\otimes_\mathcal O\mathcal O'\). In particular
\[
      \Delta(K,K',T,I)=
       \frac{\operatorname{length}_{\mathcal O}F}{v_\pi(\ell)}
\]
is invariant under finite coefficient extension, where
\(\pi\) is a uniformizer and \(v_\pi(\pi)=1\).
\end{proposition}
\begin{proof}
Choose an \(\mathcal O\)-basis of \(\mathcal O'\). For an allowed
cocycle over \(\mathcal O'_n\) first pull back its extension to the
copy of \(\mathcal O_n\) generated by \(1\) in the quotient. This
is an injection because \(\mathcal O\cap\ell^n\mathcal O'
=\ell^n\mathcal O\). Push out the kernel along any coordinate
functional \(T'_n\to T_n\) in the chosen basis. The resulting
extensions belong to the original torsion crystalline class by
the subquotient property. Thus each coordinate belongs to \(C_n\).
Conversely scalar extending any allowed extension gives, on
restriction to \(\ZZ_\ell\), a finite direct sum of copies of the
original one, since \(\mathcal O'\) is free over \(\mathcal O\).
Multiplying by the basis scalars and adding the extensions proves
\(C'_n=C_n\otimes_\mathcal O\mathcal O'\).
Taking the increasing union gives the claim for \(C\). The same
basis argument on rational cocycles split over \(B_{\mathrm{cris}}\)
gives the equality for \(C_f\), hence for the quotient by flatness.
Both its length and \(v_\pi(\ell)\) are multiplied by the relative
ramification index. This proves the formula.
\end{proof}

\begin{corollary}[A bound uniform in the auxiliary levels on a fixed interval]
\leavevmode\label{cor:fixedinterval732}
Consider the matrix families of Theorem~\ref{thm:ribettrace720}
on an actual period field \(E\). At every \(w\mid\ell\) suppose
all local members of these families and the target become crystalline
over the same fixed finite Galois extension \(K'_w/E_w\) with jumps in \(I_w\), where
\(0\in I_w\). These choices may depend on the field and the node
and are made independently of auxiliary levels away from \(\ell\).
At arbitrary precision one has
\[
    \frac{l_e}{v_\pi(\ell)}
         \le \sum_{w\mid\ell}
            \Delta(E_w,K'_w,\mathcal O(-e),I_w), \qquad 0\le e<t .
 \tag{3}\label{eq:scalarloss732}
\]
For the higher-rank lattices satisfying
Theorem~\ref{thm:twistedselmer730}, the length of the trace
cocycles for each of its selected twists satisfies
\[
    \frac{\lambda_{T_i\otimes\vartheta^s}}{v_\pi(\ell)}
        \le B_i+\varepsilon+
                 \sum_{w\mid\ell}\Delta(E_w,K'_w,T_i,I_w).
 \tag{4}\label{eq:mainloss732}
\]
When the data are spherical as in
Corollary~\ref{cor:scalarfl724}, one can take \(K'_w=E_w\) without
its numerical condition \(2H\le\ell-2\). The finite bounds then
apply on the same fields and the same nodes for weights of any
width. On any one finite interval and any fixed finite set of
targets, they are uniform over the auxiliary families and precision.
\end{corollary}
\begin{proof}
The graph construction of Lemma~\ref{lem:graphrealization724} needs
only the matrix relations. At each \(w\) its rational lattice for
the lifted extension is a sum of the given representations and
the trivial quotient. It is crystalline over \(K'_w\) in \(I_w\).
After reduction and pushout onto the target, the resulting
extension is torsion crystalline there, being a subquotient of a
finite direct sum with the target lattice. The local raw cocycle
thus belongs to \(C_n\). Map its image in \(C[\ell^n]\) to \(F\).
The kernel consists exactly of reductions of \(Z_f(E_w,T)\) at
this precision, by the Tate module description in
Theorem~\ref{thm:localfinite732}. For precision \(\pi^j\) dividing
\(\ell^n\), the same pullback and pushout embed the values into
\(A[\ell^n]\). A rational finite cocycle whose image is killed by
\(\pi^j\) becomes integral when multiplied by \(\pi^j\), giving the
same kernel description at that precision. By Lemma~\ref{lem:localevaluation728}
these are exactly the integral local evaluation equations.
The quotient measuring \(l_e\) therefore injects into the direct
sum of the indicated finite groups, proving (3).

A divisible-coefficient raw trace cocycle takes values at finite
precision by compactness. The same maps to \(F\) are compatible
with increasing the precision. At all \(w\mid\ell\) the selected
twists are trivial, so the local groups for the main target are
those for \(T_i\), after adjoining coefficient values of the
characters. Proposition~\ref{prop:coeffdefect732} applies.
The kernel of all the maps satisfies the coefficient-place
equations. The earlier unramified and normalization arguments in
Corollary~\ref{cor:maincost730} apply also here. The kernel lies
in the group of Theorem~\ref{thm:twistedselmer730}, proving (4).
In the spherical case the local-global compatibility and weight
calculation in the proof of Corollary~\ref{cor:scalarfl724}
precede its use of the Fontaine--Laffaille width and give
crystallinity and the bounded interval directly. On a finite
interval take the maximum over its finitely many relevant groups
for each fixed set of coefficient primes. For primes in the
common Fontaine--Laffaille range the original trace loss bound
from Corollary~\ref{cor:scalarfl724} applies.
\end{proof}

\begin{proposition}[A nonzero defect for wide intervals]
\leavevmode\label{prop:logexample732}
For \(\ell>2\) and any finite unramified \(K/\QQ_\ell\), let
\(T=\ZZ_\ell\) with trivial action and take \(K'=K\).
Set \(a=(\ell-1)\ell^{m-1}\), \(m\ge1\).
If \(I\) contains \([-a,0]\), then \(F\) has an element of order
\(\ell^m\).
\end{proposition}
\begin{proof}
In the crystalline sum \(\QQ_\ell(\kappa^a)\oplus\QQ_\ell\) with
basis \(e_1,e_2\), use the lattice spanned by \(e_1\) and
\(e_2+\ell^{-m}e_1\). It is stable since every \(\ell\)-adic unit
has \(a\)-th power congruent to \(1\bmod\ell^m\). Modulo \(\ell^m\)
both diagonal characters are trivial. The upper right entry is the
homomorphism
\[
          g\longmapsto
          \frac{\kappa(g)^a-1}{\ell^m}\pmod{\ell^m}.
\]
These matrices give a member of \(C_m\), since the displayed lattice
has jumps \(0,-a\). Inertia over \(K\) maps onto
\(\ZZ_\ell^\times\) under \(\kappa\), as the cyclotomic tower is
totally ramified and \(K\) is unramified. At an inertia element
with \(\kappa(g)=1+\ell\), the upper right entry is a unit.
Indeed the binomial theorem gives
\(v_\ell((1+\ell)^{\ell-1}-1)=1\), and raising a principal
unit of depth at least one to the \(\ell\)-th power increases its
depth by exactly one for odd \(\ell\). Thus the cocycle in
\(A[\ell^m]\) has an inertia value of order \(\ell^m\).
Every member of \(C_f\) for the trivial target is unramified by
Lemma~\ref{lem:cryszero717}. Its class in \(F\) has the asserted
order.
\end{proof}

The numerical loss required in Section~\ref{sec:missing} concerns
growth as the period intervals enlarge. The finite group in (3)
is independent of precision, but its size can grow with the
weight interval, as the last proposition shows. A quantitative
estimate at the weights of the theta chain and the full integral
projected-period comparison are still required.
\clearpage

\subsection{The critical-value comparison at primes dividing the discriminant or old conductor}
\label{sec:allcritical733}

In constructing a twist we may multiply every partial Hecke value
at every auxiliary level by a single nonzero scalar. The constant
may depend on the old conductor, the infinity type and the chosen
coefficient place. Its valuation disappears on dividing by the
untwisted value. We use this freedom to extend
Proposition~\ref{prop:cmraymeasure731} to the coefficient places
dividing the discriminant of \(E\) or one of the two old levels.
The argument uses the integral Eisenstein--Kronecker class before
its projection to a CM eigenspace.

\begin{lemma}[A denominator for the moment at a torsion section]
\leavevmode\label{lem:nonetalemoment733}
Let \(B\) be an abelian scheme over a local integer ring
\(R\subset\overline{\QQ}\) of residue characteristic \(\ell\).
Let \(x\) be a section killed by the positive integer \(n\).
Use the Poincaré bundle with connection on \(B\) and the formal
completion of the universal vector extension of its dual at zero.
Write \(\widehat{\mathcal P}^{\natural}\) for its completed
coefficient sheaf, and \(H\) for the cotangent module of this
universal extension at zero. For \(b\ge0\) the characteristic-zero
moment of order \(b\) at \(x\), multiplied by \(n^b\), defines an
integral map
\[
     x^*\widehat{\mathcal P}^{\natural}
              \longrightarrow {\rm TSym}^{b}H.
\]
Here \({\rm TSym}^b\) denotes symmetric tensors, with integral
divided powers. The map commutes with automorphisms fixing \(x\).
\end{lemma}
\begin{proof}
For the endomorphism \([n]\) of \(B\) the functoriality of the
Poincaré bundle gives, on every finite infinitesimal neighborhood,
an integral map \([n]_\#:\mathcal P^{\natural,(j)}
\to[n]^*\mathcal P^{\natural,(j)}\). This is Theorem~2.8 of
Kings--Sprang. Pulling back at \(x\) gives a map to the fiber at
zero, since \([n]x=0\). Compose this integral map with the integral
moment \({\rm mom}^b_0\) at zero of Definition~2.6. On the generic
fiber the definition of the torsion splitting in Corollary~2.9
and Definition~2.10 gives the commutative identity
\[
   {\rm mom}^b_0\circ[n]_{\#,x}
        ={\rm TSym}^b(d[n]^*)\circ{\rm mom}^b_x
        = n^b\,{\rm mom}^b_x .
\]
Indeed the torsion splitting first applies \([n]_{\#,x}\) and
then inverts \([n]_{\#,0}\) over the fraction field. The moment
at zero is functorial under the induced map on the formal group.
Multiplication by \(n\) acts as \(n\) on its cotangent module.
The composite on the left is defined over \(R\) even when \(n\)
is divisible by \(\ell\). All its maps are functorial for
automorphisms fixing \(x\). \cite{II-II-kingssprang731}
\end{proof}

\begin{theorem}[Integral partial values with fixed old ramification]
\leavevmode\label{thm:badcmmeasure733}
In the notation of Section~\ref{sec:samecritical731}, allow any
coefficient prime \(\ell\). Assume that the conductor
\(\mathfrak f_0\) of \(u\) contains places of two distinct rational
residue characteristics. For every \(J\) keep all new primes away
from \(\ell\) and the old conductor. Take the smoothing prime
\(\mathfrak c\) of Proposition~\ref{prop:cmraymeasure731}, with
\(N\mathfrak c\) prime to \(\ell N\mathfrak f_J\). There is a
single nonzero scalar \(D'\), independent of \(J,N\) and of the
ray classes, such that the numbers
\[
 D'\left(A L_{\mathfrak f_J}(\alpha,0,b_0)
                 -L_{\mathfrak f_J}(\alpha,0,b_0[\mathfrak c])\right),
       \qquad A=N\mathfrak c\,\alpha(\mathfrak c),
\]
are algebraic and integral at \(\iota\) for every ray class \(b_0\).
Here \(\alpha=uN^{-z}\) as before. In particular (*) in
Section~\ref{sec:samecritical731} holds with
\(C=D'(A-1)\) on the same systems of split ray characters.
\end{theorem}
\begin{proof}
Put \(\mathfrak f=\mathfrak f_J\) and choose a representative
\(\mathfrak b\in b_0\) prime to \(\ell\mathfrak f\mathfrak c\),
with \(N\mathfrak b\) prime to \(\ell N(\mathfrak f\mathfrak c)\).
For example the ray class theorem gives such a prime by excluding
the finitely many primes over the latter rational primes.
Let \(\mathcal A\) be the CM abelian variety for the fixed type
and lattice \(\mathcal O_E\). For a fractional ideal
\(\mathfrak a\) let \(\mathcal A(\mathfrak a)\) be its Serre tensor
by \(\mathfrak a\), with complex period lattice \(\mathfrak a\Omega\)
in the compatible uniformizations of Kings--Sprang, Section~4.2.
Use \(B=\mathcal A(\mathfrak f\mathfrak b^{-1})\). The point
\(x=\Omega\bmod\mathfrak f\mathfrak b^{-1}\Omega\) has annihilator
ideal \(\mathfrak f\). In fact \(a x=0\) for \(a\in\mathcal O_E\)
means \(a\mathfrak b\subseteq\mathfrak f\), which is equivalent
to \(a\in\mathfrak f\). After a finite extension, use good models
at the coefficient place as in Proposition~\ref{prop:cmraymeasure731}
and define the torsion sections there. The ideal isogenies extend
over the local integer rings by the Néron mapping property.

Write \(D_{\mathfrak c}=B[\mathfrak c]\). It is finite étale, its
order is \(N\mathfrak c\), and its isogeny and dual isogeny have
degree prime to \(\ell\). The section \(x\) is disjoint from
\(D_{\mathfrak c}\) also in the special fiber. Indeed
\(\mathfrak f_0\) has a prime divisor of residue characteristic
\(p\ne\ell\), and the \(p\)-primary component of \(x\) is nonzero
by its exact annihilator. It stays nonzero under specialization on
finite étale \(p\)-power torsion. No point of \(D_{\mathfrak c}\)
has such a component. The subgroup
\(\Gamma=\{e\in\mathcal O_E^\times:e\equiv1\bmod\mathfrak f\}\)
fixes \(x\) and preserves the trace-zero function on \(D_{\mathfrak c}\)
which is \(N\mathfrak c-1\) at zero and \(-1\) elsewhere.
This is the smoothing function \(f_{[\mathfrak c]}\) in Section~4.4
of Kings--Sprang. The group \(\Gamma\) is torsion free. Indeed the
norm of \(1-\zeta\) for a nontrivial root of unity is either a unit
or supported at a single rational prime, by the cyclotomic polynomial
at one. The congruences at the two old characteristics exclude
such \(\zeta\). Dirichlet's unit theorem gives \(\Gamma\) rank
\(g-1\) for \(g=[E:\QQ]/2\). In the evaluation use the subgroup
\(\Gamma'=\{\epsilon\in\Gamma:\epsilon\equiv1\bmod\mathfrak c\}\).
It fixes each point of \(D_{\mathfrak c}\) individually.
Its index \(e_J=[\Gamma:\Gamma']\) divides
\(|(\mathcal O_E/\mathfrak c)^\times|=N\mathfrak c-1\), a number
fixed before the auxiliary levels are chosen.

Write the infinity type of \(\alpha\) as \(\beta-\alpha'\) as in
Section~\ref{sec:samecritical731}, and let \(a=|\alpha'|-g\) and
\(b=|\beta|\). Apply the general integral construction of
Kings--Sprang, Theorem~2.20 and Definitions~2.21--2.24, up to
pullback at \(x\) of the \(a\)-th derivative of the class. Up to
this step the construction takes place on an abelian scheme with
the finite étale smoothing kernel, and uses its action by \(\Gamma'\),
its Poincaré bundle and its connection. The CM idempotents and the
moment at \(x\) have not entered. Apply
Lemma~\ref{lem:nonetalemoment733} with \(n=N\mathfrak f\), which
annihilates \(x\). The resulting integral cohomology class belongs to
\[
 H^{g-1}\!\left(\Gamma',\mathcal T(B)\right),\qquad
 \mathcal T(B)={\rm TSym}^a\omega_B\otimes
               {\rm TSym}^b H_B\otimes\bigwedge^g\omega_B .
\]
Here \(\omega_B\) denotes the invariant differentials.
Over the fraction field it is \(n^b\) times the unprojected
characteristic-zero class in Definition~2.24. This identification
uses precisely the moment identity just proved. Its integral
coefficient module is finite free over the local integer ring.
The valuation \(v_\ell(n)=v_\ell(N\mathfrak f_0)\) is fixed as
\(J\) varies. \cite{II-II-kingssprang731}

We next control the projection of this class. Over a number field
containing the embeddings of \(E\), make the CM projection of
Section~2.4 onto the single prescribed component. It projects \(H\)
using its rational CM splitting and selects the prescribed
symmetric tensors and determinant. Trivialize the resulting line
by bases of eigendifferentials and the orientation of the CM type,
as in Proposition~2.27. Before evaluating on a group cycle this
is a linear functional on \(\mathcal T(B)\) over the fraction
field. It is \(\Gamma\)-equivariant to the trivial module because
the selected character of a unit of \(\Gamma\) is trivial. This
also follows by evaluating the infinity type of \(\alpha\) on any
unit congruent to one modulo its conductor.

There is one power \(\ell^k\) making all these linear functionals
integral with the compatible period bases of Section~4.3. To see
this write \(\mathfrak f=\mathfrak f_\ell\mathfrak m\) with
\(\mathfrak f_\ell\) supported above \(\ell\). It is the fixed
\(\ell\)-part of the old conductor. First use the fixed variety
\(\mathcal A(\mathfrak f_\ell)\) and choose its rational bases by
transport from \(\mathcal A\). Over a good integral model multiply
the single functional on its finite free tensor module by a power
\(\ell^k\) so that it is integral. For every new \(B\) the diagram
\[
  \mathcal A(\mathfrak f_\ell)
       \xrightarrow{[\mathfrak b]}
         \mathcal A(\mathfrak f_\ell\mathfrak b^{-1})
       \xleftarrow{[\mathfrak m]}
         \mathcal A(\mathfrak f_\ell\mathfrak b^{-1}\mathfrak m)
          =B
\]
has both degrees prime to \(\ell\). These arrows induce integral
isomorphisms on differentials, on the module \(H\) for the
universal vector extension, and on the displayed tensor modules.
Indeed the inverse is given by an isogeny divided by its
prime-to-\(\ell\) degree on the linear modules. The arrows
intertwine the action of \(\mathcal O_E\), its rational projectors,
and the transported bases. Over \(\mathbb C\) they are induced by
the inclusions of the corresponding period lattices. Composition
with the fixed arrow \([\mathfrak f_\ell]\) gives exactly the
transport through \([\mathfrak f]\) and \([\mathfrak b]\) of
Section~4.3. Hence the same \(\ell^k\) works, including after all
base extensions used for the sections. This proves the denominator
bound for the rational projection even at \(\ell\mid D_E\).

Apply the scaled functional on cohomology and evaluate on the
oriented integral fundamental class of the free abelian group
\(\Gamma'\). This is the top exterior product of an oriented basis
of \(\Gamma'\). Over \(\mathbb C\) the projected evaluation before
multiplication by \(n^b\ell^k\) is the smoothed Eisenstein--Kronecker
value of Corollary~3.28. We give the orbit count for its conversion
to partial values. Write \(\mathfrak a=\mathfrak f\mathfrak b^{-1}\).
The coset \(1+\mathfrak c^{-1}\mathfrak a\) is the disjoint union of
\(1+t+\mathfrak a\) over \(t\in\mathfrak c^{-1}\mathfrak a/
\mathfrak a\). Each coset is \(\Gamma'\)-stable since \(\Gamma'\)
fixes \(x\) and each smoothing torsion point. The smoothed sum is
therefore \(N\mathfrak c\) times the sum for \(1+\mathfrak a\)
minus the sum for \(1+\mathfrak c^{-1}\mathfrak a\), both taken
modulo \(\Gamma'\). The two full cosets in this difference are
also \(\Gamma\)-stable and contain no zero, since
\(\mathfrak b\mathfrak c\) is prime to \(\mathfrak f\).
Multiplication by units acts freely on their nonzero elements.
The summand is invariant under \(\Gamma\) by the infinity type.
Each of the two sums modulo \(\Gamma'\) is thus \(e_J\) times the
corresponding sum modulo \(\Gamma\). The identities hold first in
the region of convergence of the Eisenstein series and then at the
critical point by analytic continuation. Formula (4.1.3) and the
period computation in Section~4.3 give the resulting scalar,
up to the fixed orientation sign, as
\[
 \frac{e_J D}{\alpha(\mathfrak b\mathfrak c)}
   \left(A L_{\mathfrak f}(\alpha,0,b_0)
               - L_{\mathfrak f}(\alpha,0,b_0[\mathfrak c])\right).
\]
Here \(D\ne0\) is the period factor for the original bases and
infinity type, independent of \(\mathfrak f,\mathfrak b\).
Indeed formula (4.1.3) weights the first coset by
\(\alpha(\mathfrak b)\) and the second by
\(\alpha(\mathfrak b\mathfrak c)\). It applies to an ideal character
defined modulo \(\mathfrak f\). Equivalently one may twist it by
a primitive auxiliary character to give exact conductor, then
remove its constant value on each ray class as in
Proposition~\ref{prop:cmraymeasure731}. All comparisons in this
paragraph are over the fraction field or over \(\mathbb C\).
\cite{II-II-kingssprang731}

The multiplier \(\alpha(\mathfrak b\mathfrak c)\) is an algebraic
\(\ell\)-adic unit. Indeed a positive power of the ideal is
principal with generator congruent to one modulo \(\mathfrak f_0\)
and a unit at every place above \(\ell\). The ideal character on
that principal ideal is a product of integral powers of conjugates
of the generator. Choose
\(D'=(N\mathfrak c-1)\ell^{k+b\,v_\ell(N\mathfrak f_0)}D\).
The quotient of this multiplier by \(e_J\ell^k n^b D\) is
\((N\mathfrak c-1)/e_J\) times a rational \(\ell\)-adic unit.
It is integral at \(\ell\). The integral evaluation just proved gives
exactly the assertion about the partial values. Summing the
coefficients in \(\overline{\ZZ}_\ell[H_J]\) and evaluating at a
character gives (*) by the reindexing in
Proposition~\ref{prop:cmraymeasure731}.
\end{proof}

\begin{corollary}[The same critical ratios at all coefficient primes]
\leavevmode\label{cor:allcritical733}
The conclusion on \(R_J(u,z,\vartheta^s)\) in
Theorem~\ref{thm:commoncritical731} holds also when \(\ell\)
divides \(D_E\) or \(N\mathfrak f_0\), provided the old conductor
of each \(u\) contains primes in two distinct rational
characteristics. Retain the other hypotheses on the lattices and
the distinguished old inertia element away from \(\ell\).
In particular for the pairs ramified at both marked levels the
cost in (2) is still at most
\(2v_\ell(Nw_*-1)+\varepsilon\) for any admissible such \(w_*\).
\end{corollary}
\begin{proof}
Choose the smoothing prime by the same finite Frobenius conditions
away from the finitely many rational primes of the old conductors,
and exclude its rational residue characteristic in selecting the
new split pairs. Finite exclusions do not change the Chebotarev
argument or \(d_0\). By Theorem~\ref{thm:badcmmeasure733} we now
have integral polynomials with the same Euler deletion identity.
The new augmentation \(b=v(\Theta_\varnothing)\) is finite and
fixed before \(M,N\). The proof of
Lemma~\ref{lem:criticalwitness731} uses the critical character only
at the new primes and at \(g_*\), all away from \(\ell\).
The specialization in Theorem~\ref{thm:commoncritical731} applies
with this \(b\), and subtracts it on dividing by the original
value. No term involving \(k\), \(n\), the conductor or \(D_E\)
survives this division. At least one of the two old
characteristics differs from \(\ell\), and the asserted inertia
cost follows there when the common marking hypotheses hold.
\end{proof}

This extends the control of the critical-value factors in the
twisted periods. The two outstanding estimates concern the
quantitative size of the local defects as the weights grow and
the full projected theta periods of Section~\ref{sec:missing}.
\clearpage

\subsection{Heights of the critical factors in the chain of lifts}
\label{sec:criticalheights734}

The twists used for integral projections can depend on the coefficient
prime and the embedding. This matters when collecting the critical
values into a determinant. For a family \(y=(y_{\ell,\iota})\) of
nonzero algebraic numbers set
\[
 H_{\rm fin}(y)=\sum_{\ell}\log\ell\,
        \sup_\iota\max\{0,-v_\ell(\iota y_{\ell,\iota})\},
             \qquad v_\ell(\ell)=1 .
\]
The sum is the supremum over finite sets of primes. A bound \(H_{\rm fin}
\le h\) gives the height bound with multiplier \(1\) for a family of
twisted periods. Indeed the multiplier has absolute value \(1\) under
every complex embedding, and the displayed local suprema are valid
exponents for the finite-place inequalities. For algebraic numbers
only the restriction of an embedding to \(\overline{\QQ}\) enters.

\begin{lemma}[The contribution of a tame generator]\leavevmode
\label{lem:tametraceheight734}
Let \(\xi\) have finite order \(D>1\). For every embedding at \(\ell\),
\[
 v_\ell(1-\xi)=
 b_\ell(D):=
 \begin{cases}
   1/\varphi(\ell^a),&D=\ell^a,\quad a\ge1,\\
   0,&D\text{ has a prime divisor different from }\ell .
 \end{cases}
\]
Consequently \(\sum_\ell b_\ell(D)\log\ell\le\log 2\).
\end{lemma}
\begin{proof}
Reduction is injective on roots of unity of order prime to \(\ell\),
by separability and Hensel's lemma. Roots of \(\ell\)-power order
reduce to \(1\). When \(D\) has a nontrivial prime-to-\(\ell\) part,
the reduction of \(\xi\) therefore differs from \(1\).
If \(D=\ell^a\), all the primitive \(D\)-th roots give the same
valuation. In fact \((1-\xi^j)/(1-\xi)=1+\cdots+\xi^{j-1}\)
reduces to \(j\ne0\) for \(1\le j\le D\) prime to \(\ell\).
Multiplication over these primitive powers gives
\(\Phi_{\ell^a}(1)=\ell\), proving the formula. At most one prime
contributes, and \(\log\ell/\varphi(\ell^a)\le\log\ell/(\ell-1)
\le\log2\). The last inequality follows from \(2^{n-1}\ge n\)
for integers \(n\ge2\).
\end{proof}

\begin{theorem}[The height of every ramified critical ratio]\leavevmode
\label{thm:criticalheights734}
At a node of the shorter family of Corollary~\ref{cor:actualmarks714},
fix one marked place \(w_i\) over each of the two rational level
primes \(p_i\), \(i=0,1\). Take the ramified critical pairs: their
tame characters are nontrivial at every marked place, their old
conductors are supported on these places and their conjugates, and
their integers \(z\) satisfy \(1\le z\le t\) and
\(z\le k_\tau\le4b\) at each infinite place, in the notation of
Section~\ref{sec:samecritical731}. For the Hecke character \(u\) of
such a pair write \(D_i(u)>1\) for its tame order at \(w_i\).
Use \(w_0\) when \(\ell\ne p_0\) and \(w_1\) when \(\ell=p_0\).
One may include any of the simultaneous Selmer tests of
Theorem~\ref{thm:twistedselmer730} satisfying its hypotheses for
these chosen inertia generators.

For every \(m\ge1\) and every sequence of positive numbers
\(\epsilon_\ell\) with \(\sum_\ell\epsilon_\ell\log\ell\le\delta\),
there is a family \(\vartheta_{\ell,\iota}\) with auxiliary modulus
prime to \(\ell\) and the prescribed Selmer and new-inertia bounds
of Theorem~\ref{thm:commoncritical731}, with error
\(\epsilon_\ell\) in the Selmer bounds. Simultaneously for all the
pairs and all \(0<|s|\le m\), the inverse of the family
\(R=R_J(u,z,\vartheta_{\ell,\iota}^s)\) satisfies
\[
 H_{\rm fin}(R^{-1})
 \le 2\sum_{\ell\ne p_0} b_\ell(D_0(u))\log\ell
      +2 b_{p_0}(D_1(u))\log p_0+\delta
 \le 4\log 2+\delta .
\]
If both orders have at least two distinct prime divisors, then
\(\max(0,v_\ell(\iota R))\le\epsilon_\ell\) at every place.
\end{theorem}
\begin{proof}
The integers \(z\), the lists of exponents and the tame residue
characters range over finite sets. For fixed such data the quotient
of two pulled-back Hecke characters is trivial at infinity and on all
finite local unit groups, hence factors through the finite ideal class
group of \(E\). Thus the list of pairs needed for all ratios is finite.
At each \((\ell,\iota)\) fix a tame generator at the selected \(w_i\).
It takes every prescribed residue character to a root of exactly its
tame order, since local reciprocity maps tame inertia onto the cyclic
residue unit group in the tame abelian quotient. The place is away
from \(\ell\), so the algebraic weight adds no inertia there. The
order is preserved by \(\iota\). Apply
Theorem~\ref{thm:commoncritical731} with this generator and the
entire list, using Corollary~\ref{cor:allcritical733} at the primes
of bad conductor or discriminant. Its proof with an empty list of
Selmer tests gives the assertion for the critical pairs alone.
The chosen twist works for the whole list at once. Its critical
ratio is positive and algebraic, including at \(z=1\), by
Section~\ref{sec:samecritical731}. With either sign of the twist the
same Euler factors at the auxiliary primes are omitted.

The pointwise estimate \eqref{eq:samebounds731} now gives
\(\max(0,v_\ell(\iota R))\le 2b_\ell(D_i(u))+\epsilon_\ell\).
Take the supremum over the embeddings and sum. The bound on each of
the two old contributions follows from Lemma~\ref{lem:tametraceheight734}.
If both orders have two distinct prime factors then the formula for
\(b_\ell\) gives zero for each old contribution at every prime.
\end{proof}

\begin{corollary}[Arbitrarily small total height for the chain ratios]\leavevmode
\label{cor:quotientchainheight734}
On the shorter fields put \(Q=(q_0d)^{h+6}\), where \(18480\mid q_0\)
and \(2^h=m^{22}\), as in Section~\ref{sec:markedmatrix715}.
Use the prescribed residue character \(\omega_w\) of order \(Q\)
at each marked place. Consider the characters of the successive
lifts in their geometric progression \(\nu_j=\nu_0\xi^j\),
\(1\le j\le m\), where \(\xi\) has tame part \(\omega_w\)
and angular weight of magnitude \(\Delta=2r^8\) at every infinite
place. Write \(t=5r\), with the prescribed \(r\ge2\).
For \(1\le i<j\le m\) form the critical factors for
\(u_{ij}=(\nu_j/\nu_i)^{-1}\). For each \((\ell,\iota)\) put
\[
 P_{\ell,\iota}=
 \prod_{1\le i<j\le m}\ \prod_{z=1}^{t}
       R_J(u_{ij},z,\vartheta_{\ell,\iota}^{-(j-i)}).
\]
Given any \(\delta>0\), the simultaneous twists in
Theorem~\ref{thm:criticalheights734} can be chosen so that
\[
 H_{\rm fin}\!\left(
      \big((N_J^{\,e/2}P_{\ell,\iota})^{-1}\big)_{\ell,\iota}
                 \right)\le\delta
                           \qquad(e\in\ZZ),
\]
where \(N_J\) is the positive norm of the auxiliary modulus for
that twist. The same choice controls the individual critical factors
of the main pairs, whose tame parts are \(\omega_w^j\) for
\(1\le j\le m\). The bound \(\delta\) can be used at every node of
each common interval of Corollary~\ref{cor:fieldrelations712}.
\end{corollary}
\begin{proof}
The local character of order \(Q\) is obtained by taking the
\(Q\)-th power of a character of order \(Q^2\) of the cyclic group
\(\mathbb F_{p_i}^{\times}\), since \(Q^2\mid p_i-1\).
The angular weights of \(u_{ij}\) have magnitude
\((j-i)\Delta\ge t\), so every displayed integer satisfies the
criticality condition. The argument for these factors also follows
by adding their finite list directly in
Theorem~\ref{thm:commoncritical731}, which requires no upper
cutoff on the magnitudes of the angular weights.

For \(1\le a\le m\), the order of \(\omega_w^{\pm a}\) is
\(Q/\gcd(Q,a)\). Both \(2^{h+6}\) and \(3^{h+6}\) divide \(Q\)
and both exceed \(m\). Hence \(2\) and \(3\) both divide this
order. The last assertion of Theorem~\ref{thm:criticalheights734}
applies at every coefficient prime to every factor of the product,
and to every individual main pair. Set \(L=t\binom{m}{2}\) and use
\(\epsilon_\ell=\delta\,2^{-\ell}/((1+L)(1+\log\ell))\).
The positive part of the valuation of a product is at most the
sum of the positive parts. Thus the supremum of that positive part
for \(P_{\ell,\iota}\) is at most \(L\epsilon_\ell\), and its
weighted sum is at most \(\delta\). The norm \(N_J\) is an
\(\ell\)-adic unit by construction. The positive real square root
used in \(N_J^{e/2}\) is algebraic, and its valuation at every
embedding above \(\ell\) is zero by squaring. This proves the claim
also for negative \(e\). The order argument and the choice of the
sum bound apply at every node with the same \(\delta\).
\end{proof}

The estimate concerns the Hecke values in the chain determinant.
The operator clearing the integral projection and the inner products
of projected test forms in Section~\ref{sec:missing} require their
own estimates on these fields.
\clearpage

\subsection{Local lifting with a separate weight interval at each embedding}
\label{sec:labelledlift735}

The automorphic weights need not have the same sign at two
embeddings of one completion. The integral lifting argument can
use a different shift at each embedding of an unramified local
field. This identifies the individual completions contributing to
the trace bound even when some other completion above the same
coefficient prime has a long interval.

\begin{theorem}[Lifting at the labelled widths]\leavevmode
\label{thm:labelledlift735}
Let \(K/\QQ_\ell\) be unramified, \(\ell>2\), and take the lattices
\(T,U_x\) of Lemma~\ref{lem:graphrealization724} over the ring of
integers \(\mathcal O\) of a finite coefficient field.
Suppose that all their rational representations are crystalline.
Fix embeddings of the coefficient field and of \(K\) in
\(\overline{\QQ}_\ell\). At each \(\tau:K\hookrightarrow
\overline{\QQ}_\ell\) let \([A_\tau,B_\tau]\) contain zero and
all labelled filtration jumps of these representations. Suppose
\[
 B_\tau-A_\tau\le\ell-1\quad\hbox{for every }\tau,\qquad
 B_{\tau_0}-A_{\tau_0}\le\ell-2\quad\hbox{for at least one }\tau_0 .
\]
For any uniformizer \(\pi\) of \(\mathcal O\) and any \(n\ge1\),
every cocycle \(z:G_K\to T/\pi^nT\) satisfying the matrix relations
of that lemma lifts, as a cocycle, to \(T\) with crystalline rational
extension. The conclusion also holds for any such cocycle whose
extension, with quotient \(\mathcal O/\pi^n\), is an
\(\mathcal O\)-linear subquotient of torsion reductions of crystalline
lattices in these same labelled intervals.
\end{theorem}
\begin{proof}
We first enlarge the coefficient field to split \(K\) and retain the
ideal generated by the original \(\pi^n\) as precision. Put
\(S=\mathcal O_K\otimes_{\ZZ_\ell}\mathcal O\).
It is a product of copies of \(\mathcal O\) indexed by the
embeddings \(\tau\). The Witt Frobenius permutes the factors
cyclically. Denote the next index by \(\sigma\tau\).
In the convention in which the cyclotomic character has jump
\(-1\), we need a crystalline character of jump \(j_\tau=-A_\tau\)
at each \(\tau\). We construct its lattice \(D\).

Set \(h=\ell-1\). In the \(\tau\)-component of a rank-one
Fontaine--Laffaille module take a basis element \(d_\tau\) of
filtration degree \(j_\tau\). Thus the filtration contains this
whole component through that degree and is zero on it above that
degree. Define the divided Frobenius by
\[
       \phi_i(d_\tau)=\ell^{j_\tau-i}d_{\sigma\tau}
                           \quad(0\le i\le j_\tau).
\]
The relations \(\phi_i=\ell\phi_{i+1}\) on the next filtration
and generation by the images follow directly. Here \(0\le
j_\tau\le B_\tau-A_\tau\le h\), and \(j_{\tau_0}<h\).
The module belongs to the unipotent category at weight \(h\).
To check this, recall that its excluded quotients are the
nonzero modules whose whole filtration is in degree \(h\).
Our component at \(\tau_0\) has zero filtration in that degree.
Strictness forces the \(\tau_0\)-component of any such quotient
to be zero. In the quotient the images of \(\phi_h\) generate,
since each lower divided Frobenius is a multiple of \(\phi_h\).
Moving around the cyclic permutation of components gives zero in
every component. The covariant integral Fontaine--Laffaille
comparison therefore gives a rank-one crystalline lattice \(D\)
with the stated jumps. We use also its exactness, full faithfulness,
closure of the torsion image under subquotients, and the lattice
comparison in the unipotent weight-\(\ell-1\) range.
~\cite[\S2.3]{II-II-II-ghls724}

Tensor the given extension and its presenting lattices with \(D\).
All their rational jumps now belong to \([0,h]\), and at \(\tau_0\)
they are all strictly below \(h\). Each presenting lattice belongs
to the comparison category. Indeed its rational representation
can have no nonzero quotient that is an unramified representation
twisted by the character of constant jump \(h\). Such a quotient
would require jump \(h\) at \(\tau_0\), since the de Rham filtration
on crystalline representations is strict on quotients. This verifies
the extra hypothesis in the endpoint lattice comparison. For the
matrix relations the required presentation is the graph lattice
of Lemma~\ref{lem:graphrealization724} together with \(T\). In that
lemma one may use \(\pi^n\) throughout, by the same elementary
divisor argument. Its rational constituents are copies of the given
representations and the trivial line. For the more general assertion
use the stipulated presentation. Tensoring over \(\mathcal O\)
with the free line preserves subquotients, so the twisted torsion
extension is in the comparison image.

Let \(M\) be the module corresponding to \(T\otimes D\), and
put subscripts \(n\) on reductions at the chosen precision.
The filtered extension of \(D_n\) by \(M_n\) is strict by exactness
in the Fontaine--Laffaille category. Lift \(d_\tau\bmod\pi^n\) to
\(v_\tau\) in filtration \(j_\tau\) in the middle module.
This splits its underlying module componentwise. In each filtration
degree \(i\le j_\tau\), the \(\tau\)-component is the sum of
\(\mathrm{Fil}^i(M_n)_\tau\) and the line generated by \(v_\tau\).
In higher degrees it equals \(\mathrm{Fil}^i(M_n)_\tau\), by
strictness. Thus
\(\phi_{j_\tau}(v_\tau)=v_{\sigma\tau}+a_\tau\) for some
\(a_\tau\in(M_n)_{\sigma\tau}\).
Lift \(a_\tau\) freely to \(\widehat a_\tau\in M_{\sigma\tau}\).
Adjoin an integral generator \(\widehat v_\tau\) in degree
\(j_\tau\) in each component, with
\[
 \phi_i(\widehat v_\tau)=
       \ell^{j_\tau-i}(\widehat v_{\sigma\tau}+\widehat a_\tau)
                                   \quad(0\le i\le j_\tau).
\]
The images generate \(M\) and project onto every quotient line,
so they generate the new module. Its filtration in degree \(h\)
is zero at \(\tau_0\): this holds for \(M\) by its labelled
jumps and for the added line because \(j_{\tau_0}<h\).
The preceding quotient argument shows that the new module is
unipotent. It reduces to the given filtered extension with both
end identifications. Applying the integral comparison, undoing
the twist, and adjusting an integral section by a lift of a
coboundary gives a cocycle lifting \(z\) at every group element.

Finally, for a finite extension \(\mathcal O'/\mathcal O\) of
coefficient rings, \(1\) extends to an \(\mathcal O\)-basis of
\(\mathcal O'\) since its reduction modulo \(\pi\) is nonzero.
Project onto this basis vector. Applying the projection to the
lifted \(T\otimes\mathcal O'\)-valued cocycle gives an integral
lift with values in \(T\). Rationally the operation is restriction
of coefficients, pullback on the quotient along the inclusion
of the original coefficient field, and pushout on the submodule
along the projection. Direct sums and subquotients preserve
crystallinity. This gives the claim over the original field.
\end{proof}

\begin{corollary}[The remaining completions in the trace estimate]\leavevmode
\label{cor:labelledtrace735}
In \eqref{eq:scalarloss732} and \eqref{eq:mainloss732}, form
\([A_\tau,B_\tau]\) separately for each trace family and its target
by taking the minimum and maximum of their jumps and zero. At a
spherical unramified completion \(E_w\) of odd residue characteristic,
omit the term at \(w\) if these labelled widths satisfy
Theorem~\ref{thm:labelledlift735}. If every completion above \(\ell\)
can be omitted, then \(l_e=0\) for the scalar targets in question
and \(\lambda_{T_i\otimes\vartheta^s}/v_\pi(\ell)
\le B_i+\epsilon_\ell\) for the selected main targets. This holds
at every auxiliary level with the prescribed infinity type, for
the same twists as in Theorem~\ref{thm:criticalheights734}.
\end{corollary}
\begin{proof}
The spherical compatibility and the calculation of the jumps in
Corollary~\ref{cor:scalarfl724} apply at \(E_w\) before any width
restriction. Each local matrix relation holds on the restriction
of a global trace cocycle. The twists are trivial on \(G_{E_w}\).
Apply Theorem~\ref{thm:labelledlift735} at this completion. In the
maps used to prove \eqref{eq:scalarloss732} and
\eqref{eq:mainloss732} the image there is zero, since all integral
crystalline evaluation equations hold on the reduction of an
integral crystalline cocycle. The injection into the sum over
the other completions consequently has the same kernel as before.
The numerical Selmer bound on this kernel was imposed at the
same time as the critical-value inequalities. The infinity type
determines the overlying lists of jumps independently of the
auxiliary level, as in Corollary~\ref{cor:scalarfl724}.
\end{proof}

\begin{proposition}[Counting the completions containing a wide label]\leavevmode
\label{prop:orbits735}
Let \(E/\QQ\) be one of the abelian period fields. At an unramified
odd spherical prime \(\ell\) fix an identification of its complex
and \(\ell\)-adic embeddings on the algebraic numbers. Write
\(\mathcal B\) for the set of embeddings of \(E\) where the
labelled width is greater than \(\ell-2\), using the jumps of a
given trace family and its target. The contribution to the two
inequalities in the corollary is supported at completions whose
Frobenius orbits meet \(\mathcal B\). If \(f_\ell\) is the order
of Frobenius in \(\operatorname{Gal}(E/\QQ)\), the sum of the local
degrees of these completions is at most
\(\min([E:\QQ], f_\ell|\mathcal B|)\). An orbit containing a
width at most \(\ell-2\) and having all widths at most \(\ell-1\)
can also be omitted.

In particular, suppose the permitted jumps at a label \(\tau\)
are in
\[
 [-\eta,\eta]\ \cup\
       s_\tau [S_\tau,S_\tau+W],\qquad s_\tau\in\{-1,1\},
 \qquad S_\tau,W,\eta\ge0 .
\]
One may count only the orbits meeting labels for which
\(\max(2\eta,\eta+S_\tau+W)>\ell-2\).
\end{proposition}
\begin{proof}
The decomposition group acts on the embeddings into
\(\overline{\QQ}_\ell\) with one orbit for each completion.
Here the orbits are precisely the cycles of the Frobenius
permutation, each of size \(f_\ell\), since the field is Galois and
unramified. An orbit disjoint from \(\mathcal B\) satisfies the
lifting theorem. Each remaining orbit contains at least one
member of \(\mathcal B\), so the sum of their sizes has the
asserted bound. The endpoint assertion uses the second width
condition in Theorem~\ref{thm:labelledlift735}.
For \(s_\tau=1\) the two displayed intervals and zero lie between
\(-\eta\) and \(\max(\eta,S_\tau+W)\).
For \(s_\tau=-1\) reflect both intervals through zero.
The width estimate follows in both cases without combining
different signs from different embeddings.
\end{proof}

These omissions apply to the trace cocycles actually entering the
determinant. The lengths at the remaining completions, the residual
groups in Section~\ref{sec:missing}, and the full projected inner
products still require quantitative bounds on the common intervals.
\clearpage

\subsection{The sum of the marked Selmer lengths on a common interval}
\label{sec:oddselmer739}

The marked representations in the shorter family have odd rank,
and the cyclotomic string also has odd length. Taking a determinant
in a congruence with a cyclotomic twist exploits both parities. It
restricts the primes contributing to the global term even when the
trace of an individual matrix is zero. The other terms come from
nontrivial roots of unity at the two original levels. We estimate
their sum over coefficient primes before taking any new trace levels.

\begin{proposition}[The determinant restriction on the trace ideal]
\leavevmode\label{prop:oddtraceideal739}
Let \(E=\QQ(\mu_{q_1},\sqrt{-d})\) be as in
Proposition~\ref{prop:oddcyclo737}. Let \(T\) be an integral
polarized representation as in Theorem~\ref{thm:selftwist728},
of odd rank \(h\) with odd string length \(t\). For its trace
ideal \(I_t\), with \(v_\ell(\ell)=1\), one has
\[
           a_\ell(t)\le v_\ell(I_t)\le a_\ell(ht),
\]
where \(a_\ell\) is the explicit cyclotomic valuation of
Proposition~\ref{prop:oddcyclo737}. In particular
\(v_\ell(I_t)=0\) for every \(\ell\nmid q_1d\). The bounds hold
at every coefficient embedding and for finite-order
anticyclotomic twists.
\end{proposition}
\begin{proof}
Let \(J_u\) be the ideal generated by \(1-\kappa(g)^u\),
\(g\in G_E\), in the coefficient ring \(\mathcal O\). Integral
traces give \(I_t\subseteq J_t\). To obtain the other inclusion,
write \(\pi\) for a uniformizer and take any \(j\ge1\) such that
\(I_t\subseteq\pi^j\mathcal O\). At this precision \(T\) and
\(T(t)\) have equal traces on the group. They are isomorphic by
the full-matrix-span argument in Theorem~\ref{thm:selftwist728}.
Taking the determinant of the intertwining identity gives
\[
        \det T(g)\equiv \kappa(g)^{ht}\det T(g)\pmod{\pi^j}
                         \qquad (g\in G_E).
\]
The determinants are units. Thus \(J_{ht}\subseteq\pi^j\mathcal O\)
at every such precision. Since \(\mathcal O\) is a discrete
valuation ring, it follows that \(J_{ht}\subseteq I_t\), also if
\(I_t\) is the unit ideal or the zero ideal. The cyclotomic image
contains an open group of principal units, so \(J_{ht}\ne0\).
The two inclusions and Proposition~\ref{prop:oddcyclo737}, applied
to the odd integers \(t\) and \(ht\), give the bounds. A character
twist multiplies each generator of \(I_t\) by a unit. Its
anticyclotomicity preserves the polarization.
\end{proof}

\begin{lemma}[The primes dividing an inertia order]\leavevmode
\label{lem:orderprimes739}
For an integer \(M\ge1\) and \(y=\max(2,\lceil\log M\rceil)\),
\[
        \sum_{\ell\mid M}\frac{\log\ell}{\ell-1}\le 4\log y+1 .
\]
If \(\rho\ne1\) is any root of unity of order dividing \(M\),
then \(v_\ell(1-\rho)\le\mathbf1_{\ell\mid M}/(\ell-1)\).
\end{lemma}
\begin{proof}
The second assertion follows from Lemma~\ref{lem:tametraceheight734}.
For \(\ell\le y\), at least \(\lfloor y/\ell\rfloor\ge y/(2\ell)\)
factors of \(y!\) are divisible by \(\ell\). Taking valuations and
then logarithms gives
\[
   \sum_{\ell\le y}\frac{\log\ell}{\ell-1}
      \le 2\sum_{\ell\le y}\frac{\log\ell}{\ell}
      \le \frac4y\log(y!)\le 4\log y .
\]
Each \(\ell>y\) has \(\ell-1\ge y\), and the product of the
distinct primes dividing \(M\) divides \(M\). Their remaining
contribution is at most \((\log M)/y\le1\).
\end{proof}

\begin{theorem}[A uniform bound for the marked Selmer groups in odd rank]
\leavevmode\label{thm:oddselmer739}
Use the finitely many distinct node data on one shorter period
interval of Corollary~\ref{cor:fieldrelations712}. At a node \(x\)
write \(E_x=\QQ(\mu_{q_{1,x}},\sqrt{-d_x})\),
\(N_x=q_{1,x}d_x\), and \(p_{x,0},p_{x,1}\) for its levels.
At each \((x,\ell,\iota)\) take a finite list of untwisted marked
main lattices of the following kind. The rational representation
\(V\) is de Rham at all \(w\mid\ell\) and has polarization
\(V^c\simeq V^\vee(1-t)\), with odd rank \(h\le R\) and odd
string length \(t\le U\). As in Corollary~\ref{cor:maincost730},
we use the allowed scaled squarefree
or rank-one spectra excluding \(1\) at each of the common old
marked places away from \(\ell\), and the primitive or rank-one
type at some place above each old level prime for each lattice.
All ramification away from \(\ell\) is at the marked places. Put
\[
\begin{aligned}
    e_*&=\max_x[E_x:\QQ],&
    N_*&=\operatorname{lcm}_{x}N_x,\\
    M_*&=\prod_x\prod_{i=0}^1\prod_{j=1}^R(p_{x,i}^{\,j}-1),&
    y_*&=\max(2,\lceil\log M_*\rceil).
\end{aligned}
\]
For a signed exponent \(0<|s|\le m_x\) let \(b_\ell(T,s)\)
be the normalized length of the divisible Selmer group for
\(T\otimes\vartheta_{x,\ell,\iota}^{\,s}\) with the integral
coefficient-place equations and the ramification allowed in
Theorem~\ref{thm:twistedselmer730}. For any \(\delta>0\), the
twists used in the critical ratios can be chosen with
\[
\begin{split}
 \sum_{\ell}(\log\ell)
   \sup_{x,\iota,T,s} b_\ell(T,s)
 \ \le\ &\log\!\big(N_*\operatorname{lcm}(1,\ldots,RU)\big)\\
        &+(2e_*+2)R\,(4\log y_*+1)+\delta .
\end{split}
\]
The theorem concerns the Selmer groups with the stated integral
equations at every coefficient prime. The same bound applies to
the raw main trace lengths, normalized by \(v_\pi(\ell)\), over
those data for which Corollary~\ref{cor:labelledtrace735} applies
at every coefficient completion, uniformly in the later auxiliary
trace levels and precisions.
\end{theorem}
\begin{proof}
Choose one marked place over each level prime as in
Theorem~\ref{thm:criticalheights734}. For \(g_*\) take a full tame
generator at the first of these lying away from \(\ell\). For each
lattice a possibly different place of primitive or rank-one type
over that same rational prime gives the integral matrix span and a
scalar inertia element with \(y^2-1\) a unit, by
Theorem~\ref{thm:markedmatrix715}. The common element \(g_*\)
has no eigenvalue \(1\) on any lattice by the hypothesis at
every old mark. Its conjugate inverse has the same eigenvalues,
by polarization and the triviality of \(\kappa\) on inertia
away from \(\ell\). Thus the common and individual inertia
conditions in the twisting theorem hold simultaneously. The
unit \(y-1\) makes the normalization of a cocycle at the
individual scalar element unique in its cohomology class,
since changing it by a coboundary changes its value there
by \((y-1)v\).

At a mark \(w\nmid\ell\) above \(p_{x,i}\) each scaled inertia
eigenvalue \(\rho\) has order dividing \(p_{x,i}^{\,j}-1\)
for some \(1\le j\le h\). Indeed an irreducible residue
factor of degree \(j\) has its roots in
\(\mathbb F_{p_{x,i}^{\,j}}^\times\), since the level prime splits
completely in \(E_x\). Multiplication by a residue scalar leaves
the root in that group. Naming the roots in characteristic zero
preserves their orders. The rank-one character also has order
dividing \(p_{x,i}-1\). Therefore for a full tame generator
\(\sigma_w\),
\[
 v_\ell\bigl(\det(1-T(\sigma_w))\bigr)
             \le \frac{R}{\ell-1}\,\mathbf1_{\ell\mid M_*}.
\]
Here we multiplied the factors \(1-\rho\) and applied
Lemma~\ref{lem:orderprimes739}. The same bound holds for the
inverse transpose, because the determinant of \(T(\sigma_w)\)
is a unit. On the finite Tate dual the action of \(\sigma_w\)
is this inverse transpose. Its fixed vectors contain the
invariants of the full decomposition group. The elementary
divisors of the difference from the identity bound their
normalized length at every precision by the displayed
determinant valuation.

Take \(S\) in Theorem~\ref{thm:twistedselmer730} to consist
of the coefficient places and the specified old marked places
with their conjugates. At most \(2e_*\) old places occur.
In the constant \(B_i\) of that theorem, their invariant
groups and the two distinguished determinants together cost
at most \((2e_*+2)R\mathbf1_{\ell\mid M_*}/(\ell-1)\). The
other term of \(B_i\) is \(v_\ell(I_t)\). By
Proposition~\ref{prop:oddtraceideal739} and the explicit
cyclotomic formula it is at most
\[
        v_\ell\!\big(N_*\operatorname{lcm}(1,\ldots,RU)\big).
\]
This follows since \(ht\le RU\), and at any prime dividing
\(N_x\) the conductor part of \(a_\ell(ht)\) is at most
\(v_\ell(N_x)\).

Choose positive \(\epsilon_\ell\) with
\(\sum_\ell\epsilon_\ell\log\ell\le\delta\). At each node and
embedding insert its entire finite list in
Theorem~\ref{thm:criticalheights734}. This supplies the same
twist for these Selmer groups and the critical ratios. It
bounds their normalized lengths by \(B_i+\epsilon_\ell\)
for all the signed powers. Its tame tolerance can also be
chosen as in Corollary~\ref{cor:alltwistcost738}. All the
pointwise majorants use the displayed interval data, so they
still hold after taking the supremum. Factoring the fixed
integer in the cyclotomic bound and applying
Lemma~\ref{lem:orderprimes739} to \(M_*\) proves the inequality.
At the indicated spherical coefficient completions,
Corollary~\ref{cor:labelledtrace735} puts the raw normalized
main cocycles into this same Selmer group at each precision.
This proves the final assertion independently of the later
trace levels.
\end{proof}

The argument bounds the congruence and marked-inertia terms
using only the fields, ranks and two original levels of the
interval. To apply it to all trace cocycles one must also bound
their defect in the local integral equations outside the labelled
lifting range. The complementary-sign strict groups of
Section~\ref{sec:missing} enter a different term of the period
comparison.
\clearpage

\clearpage
\section{References}\label{nav:bibliography-II}
\begingroup\def\section*#1{}\begin{thebibliography}{99}\small
\bibitem{II-II-II-duality13} J. S. Milne, \emph{Arithmetic Duality Theorems}, 2nd ed., 2006, Chapters I and II. \href{https://www.jmilne.org/math/Books/ADTnot.pdf}{Book}.
\bibitem{II-II-II-gp13} C. Greither and C. D. Popescu, \emph{An equivariant main conjecture in Iwasawa theory and applications}, J. Algebraic Geom. \textbf{24} (2015), 629--692. \href{https://arxiv.org/abs/1103.3069}{Authors' paper}.
\bibitem{II-II-II-fw13} B. Ferrero and L. C. Washington, \emph{The Iwasawa invariant \(\mu_p\) vanishes for abelian number fields}, Ann. of Math. \textbf{109} (1979), 377--395. \href{https://doi.org/10.2307/1971116}{Original paper}.
\bibitem{II-II-II-bernoulli13} NIST, \emph{Digital Library of Mathematical Functions}, Section 25.11, formula 25.11.14. \href{https://dlmf.nist.gov/25.11}{Hurwitz zeta function}.
\bibitem{II-II-II-milne} J. S. Milne, \emph{Class Field Theory}, version 4.03, 2020, Chapter V, Theorem 3.27.\enspace\href{https://www.jmilne.org/math/CourseNotes/CFT.pdf}{Course notes}.
\bibitem{II-II-II-weibel15} C. A. Weibel, \emph{The \(K\)-book}, Grad. Stud. Math. 145, AMS, 2013. \href{https://sites.math.rutgers.edu/~weibel/Kbook/Kbook.VI.pdf}{Chapter VI, Section 8}.
\bibitem{II-II-II-kahn15} B. Kahn, \emph{The Quillen--Lichtenbaum conjecture at the prime 2}, Theorem 1. \href{https://webusers.imj-prg.fr/~bruno.kahn/preprints/QL2.pdf}{Author's paper}.
\bibitem{II-II-II-iwasawa15} K. Iwasawa, \emph{On \(\ZZ_\ell\)-extensions of algebraic number fields}, Ann. of Math. \textbf{98} (1973), 246--326.
\bibitem{II-II-II-beh17} E. Bayer-Fluckiger, V. Emery and J. Houriet, \emph{Hermitian lattices and bounds in \(K\)-theory of algebraic integers}, Doc. Math., Extra Vol. Merkurjev (2015), 71--83. \href{https://arxiv.org/abs/1406.5953}{Authors' paper}.
\bibitem{II-II-II-mathew17} A. Mathew, \emph{Torsion exponents in stable homotopy and the Hurewicz homomorphism}, Algebr. Geom. Topol. \textbf{16} (2016), 1025--1041. \href{https://msp.org/agt/2016/16-2/agt-v16-n2-p12-s.pdf}{Original paper}.
\bibitem{II-II-I-jkpsz100} B. W. Jordan, Z. Klagsbrun, B. Poonen, C. Skinner and Y. Zaytman, \emph{Statistics of \(K\)-groups modulo \(p\) for the ring of integers of a varying quadratic number field}, Tunisian J. Math. \textbf{2} (2020). \href{https://msp.org/tunis/2020/2-2/tunis-v2-n2-p02-p.pdf}{Theorems 2.5 and 5.1}.
\bibitem{II-II-V-cft101} J. S. Milne, \emph{Class Field Theory}. \href{https://www.jmilne.org/math/CourseNotes/CFT.pdf}{Local and global Artin reciprocity and the correspondence with idele class quotients}.
\bibitem{II-II-V-stick} L. C. Washington, \emph{Introduction to Cyclotomic Fields}, 2nd ed., Springer, 1997. \href{https://link.springer.com/book/10.1007/978-1-4612-1934-7}{Chapter 6, Stickelberger's theorem}.
\bibitem{II-II-II-gw712} H. Darmon, F. Diamond and R. Taylor, \emph{Fermat's Last Theorem}, Theorems 2.17--2.19, especially the number-field formula in Theorem 2.19. \href{https://www.math.mcgill.ca/darmon/pub/Articles/Expository/05.DDT/paper.pdf}{Authors' text}.
\bibitem{II-II-II-chebotarev713} J. S. Milne, \emph{Algebraic Number Theory}, Theorem 8.31 (Chebotarev density theorem). \href{https://www.jmilne.org/math/CourseNotes/ANT301.pdf}{Author's text}.
\bibitem{II-II-II-bktate717} S. Bloch and K. Kato, \emph{\(L\)-functions and Tamagawa numbers of motives}, The Grothendieck Festschrift, Vol.~I, Progress in Mathematics \textbf{86}, Birkh\"auser, 1990, 333--400, Corollary 3.8.4 and Example 3.9. \href{https://virtualmath1.stanford.edu/~conrad/BSDseminar/refs/BKTamagawa.pdf}{Published chapter}.
\bibitem{II-II-II-fl724} J.-M. Fontaine and G. Laffaille, \emph{Construction de représentations \(p\)-adiques}, Ann. Sci. École Norm. Sup. (4) \textbf{15} (1982), 547--608. \href{https://www.numdam.org/articles/10.24033/asens.1437/}{Published paper}.
\bibitem{II-II-II-ghls724} T. Gee, F. Herzig, T. Liu and D. Savitt, \emph{Potentially crystalline lifts of certain prescribed types}, Section~2.3. \href{https://math.jhu.edu/~savitt/papers/pdfs/FLlifts.pdf}{Authors' paper}.
\bibitem{II-II-II-guerberoff724} L. Guerberoff, \emph{Modularity lifting theorems for Galois representations of unitary type}, Compos. Math. \textbf{147} (2011), 1022--1058. \href{https://doi.org/10.1112/S0010437X10005154}{Published paper}.
\bibitem{II-II-kingssprang731} G. Kings and J. Sprang, \emph{Eisenstein--Kronecker classes, integrality of critical values of Hecke \(L\)-functions and \(p\)-adic interpolation}, Ann. of Math. (2) 202 (2025), 1--109. \href{https://annals.math.princeton.edu/2025/202-1/p01}{Published paper}. \href{https://arxiv.org/pdf/1912.03657\#page=52}{Partial ray-class integrality, Corollary 4.13 of the authors' September 2024 text}.
\bibitem{II-II-liulimit732} T. Liu, \emph{Torsion \(p\)-adic Galois representations and a conjecture of Fontaine}, Ann. Sci. Éc. Norm. Supér. (4) 40 (2007), 633--674. \href{https://www.numdam.org/item/ASENS_2007_4_40_4_633_0/}{Article}.
\bibitem{II-II-gaoliulimit732} H. Gao and T. Liu, Appendix B to H. Gao, \emph{Breuil--Kisin modules and integral \(p\)-adic Hodge theory}, J. Eur. Math. Soc. 25 (2023), 3979--4032. \href{https://doi.org/10.4171/JEMS/1278}{Theorem B.0.1 and its corrected proof}.
\end{thebibliography}\endgroup
\end{document}
